
\Data\ID{1003020401}
\Updated{2023-06-20 09:06:32}
\Level{A}
\SubArea{Conic Sections}
\LANGen{
\Question{
The equation of an ellipse is given by \( 5x^2+9y^2-30x-18y+9=0 \). The length of its semi-major axis is:}
{\( 3 \)}{\( 5 \)}{\( 9 \)}{\( \sqrt5 \)}{\( 14 \)}{}}
\LANGpl{
\Question{
Dana jest elipsa o równaniu \( 5x^2+9y^2-30x-18y+9=0 \). Długość półosi wielkiej to:}
{\( 3 \)}{\( 5 \)}{\( 9 \)}{\( \sqrt5 \)}{\( 14 \)}{}}
\LANGcs{
\Question{
Elipsa je dána rovnicí \( 5x^2+9y^2-30x-18y+9=0 \). Délka hlavní poloosy této elipsy je:}
{\( 3 \)}{\( 5 \)}{\( 9 \)}{\( \sqrt5 \)}{\( 14 \)}{}}
\LANGsk{
\Question{
Elipsa je daná rovnicou \( 5x^2+9y^2-30x-18y+9=0 \). Dĺžka hlavnej polosi tejto elipsy je:}
{\( 3 \)}{\( 5 \)}{\( 9 \)}{\( \sqrt5 \)}{\( 14 \)}{}}
\LANGes{
\Question{
Dada la ecuación de la elipse \( 5x^2+9y^2-30x-18y+9=0 \), ¿Cuánto mide el semieje mayor?}
{\( 3 \)}{\( 5 \)}{\( 9 \)}{\( \sqrt5 \)}{\( 14 \)}{}}
\End

\Data\ID{1003020402}
\Updated{2022-08-25 10:08:06}
\Level{A}
\SubArea{Conic Sections}
\LANGen{
\Question{
The equation of an ellipse is given by \( x^2+4y^2-6x+32y+48=0 \). The length of its semi-minor axis is:}
{\( \frac52 \)}{\( 5 \)}{\( 25 \)}{\( \frac{25}4 \)}{\( \frac{15}2 \)}{}}
\LANGpl{
\Question{
Dana jest elipsa o równaniu  \( x^2+4y^2-6x+32y+48=0 \). Długość półosi małej to:}
{\( \frac52 \)}{\( 5 \)}{\( 25 \)}{\( \frac{25}4 \)}{\( \frac{15}2 \)}{}}
\LANGcs{
\Question{
Elipsa je dána rovnicí \( x^2+4y^2-6x+32y+48=0 \). Délka vedlejší poloosy této elipsy je:}
{\( \frac52 \)}{\( 5 \)}{\( 25 \)}{\( \frac{25}4 \)}{\( \frac{15}2 \)}{}}
\LANGsk{
\Question{
Elipsa je daná rovnicou \( x^2+4y^2-6x+32y+48=0 \). Dĺžka vedľajšej polosi tejto elipsy je:}
{\( \frac52 \)}{\( 5 \)}{\( 25 \)}{\( \frac{25}4 \)}{\( \frac{15}2 \)}{}}
\LANGes{
\Question{
Dada la ecuación de la elipse \( x^2+4y^2-6x+32y+48=0 \), ¿Cuánto mide el semieje menor?}
{\( \frac52 \)}{\( 5 \)}{\( 25 \)}{\( \frac{25}4 \)}{\( \frac{15}2 \)}{}}
\End

\Data\ID{1003020403}
\Updated{2022-08-25 10:08:06}
\Level{A}
\SubArea{Conic Sections}
\LANGen{
\Question{
An ellipse has the  centre at \( S=[-1;3] \), the semi major axis with the length of \( 3 \) and the semi minor axis with the length of \( 2 \). The semi major axis is parallel to \( y \) axis. The equation of the ellipse is:}
{\( \frac{(x+1)^2}4+ \frac{(y-3)^2}9 = 1 \)}{\( \frac{(x+1)^2}9 + \frac{(y-3)^2}4 = 1 \)}{\( \frac{(x-1)^2}4+ \frac{(y+3)^2}9 = 1 \)}{\( \frac{(x+1)^2}4-\frac{(y-3)^2}9 = 1 \)}{\( \frac{(x+1)^2}4 + \frac{(y-3)^2}9 = -1 \)}{}}
\LANGpl{
\Question{
Elipsa ma środek w punkcie \( S=[-1;3] \), długość półosi wielkiej to  \( 3 \), długość półosi małej jest równa \( 2 \). Półoś wielka jest równoległa do osi \( y \). Wskaż równanie elipsy.}
{\( \frac{(x+1)^2}4+ \frac{(y-3)^2}9 = 1 \)}{\( \frac{(x+1)^2}9 + \frac{(y-3)^2}4 = 1 \)}{\( \frac{(x-1)^2}4+ \frac{(y+3)^2}9 = 1 \)}{\( \frac{(x+1)^2}4-\frac{(y-3)^2}9 = 1 \)}{\( \frac{(x+1)^2}4 + \frac{(y-3)^2}9 = -1 \)}{}}
\LANGcs{
\Question{
Elipsa má střed \( S=[-1;3] \), hlavní poloosu délky  \( 3 \) a vedlejší poloosu délky \( 2 \). Hlavní poloosa je přitom rovnoběžná s osou \( y \). Rovnice této elipsy je:}
{\( \frac{(x+1)^2}4+ \frac{(y-3)^2}9 = 1 \)}{\( \frac{(x+1)^2}9 + \frac{(y-3)^2}4 = 1 \)}{\( \frac{(x-1)^2}4+ \frac{(y+3)^2}9 = 1 \)}{\( \frac{(x+1)^2}4-\frac{(y-3)^2}9 = 1 \)}{\( \frac{(x+1)^2}4 + \frac{(y-3)^2}9 = -1 \)}{}}
\LANGsk{
\Question{
Elipsa má stred \( S=[-1;3] \), hlavnú poloos dĺžky  \( 3 \) a vedľajšiu poloos dĺžky \( 2 \). Hlavná poloos je pritom rovnobežná s osou \( y \). Rovnica tejto elipsy je:}
{\( \frac{(x+1)^2}4+ \frac{(y-3)^2}9 = 1 \)}{\( \frac{(x+1)^2}9 + \frac{(y-3)^2}4 = 1 \)}{\( \frac{(x-1)^2}4+ \frac{(y+3)^2}9 = 1 \)}{\( \frac{(x+1)^2}4-\frac{(y-3)^2}9 = 1 \)}{\( \frac{(x+1)^2}4 + \frac{(y-3)^2}9 = -1 \)}{}}
\LANGes{
\Question{
Una elipse tiene el centro en \( S=(-1;3) \), el semieje mayor mide \( 3 \) y el semieje menor mide \( 2 \). El semieje mayor es paralelo al eje \( y \). La ecuación de la elipse es:}
{\( \frac{(x+1)^2}4+ \frac{(y-3)^2}9 = 1 \)}{\( \frac{(x+1)^2}9 + \frac{(y-3)^2}4 = 1 \)}{\( \frac{(x-1)^2}4+ \frac{(y+3)^2}9 = 1 \)}{\( \frac{(x+1)^2}4-\frac{(y-3)^2}9 = 1 \)}{\( \frac{(x+1)^2}4 + \frac{(y-3)^2}9 = -1 \)}{}}
\End

\Data\ID{1003024002}
\Updated{2022-08-25 10:08:36}
\Level{A}
\SubArea{Conic Sections}
\LANGen{
\Question{
The equation of an ellipse is given by \( \frac{x^2}{48} +\frac{y^2}{36} = 1\). The point \( [-2;5] \)  represents:}
{a point inside the ellipse}{vertex of the ellipse}{co-vertex of the ellipse}{a point outside the ellipse}{the centre of the ellipse}{}}
\LANGpl{
\Question{
Elipsa jest opisana równaniem \( \frac{x^2}{48} +\frac{y^2}{36} = 1\). Punkt \( [-2;5] \)  to:}
{punkt leżący wewnątrz elipsy}{wierzchołek elipsy}{punkt przecięcia elipsy z osią małą}{punkt leżący na zewnątrz elipsy}{środek elipsy}{}}
\LANGcs{
\Question{
Elipsa je dána rovnicí \( \frac{x^2}{48} +\frac{y^2}{36} = 1\). Bod \( [-2;5] \)  představuje:}
{bod ležící uvnitř elipsy}{hlavní vrchol elipsy}{vedlejší vrchol elipsy}{bod ležící vně elipsy}{střed elipsy}{}}
\LANGsk{
\Question{
Elipsa je daná rovnicou \( \frac{x^2}{48} +\frac{y^2}{36} = 1\). Bod \( [-2;5] \)  predstavuje:}
{bod ležiaci vnútri elipsy}{hlavný vrchol elipsy}{vedľajší vrchol elipsy}{bod ležiaci von z elipsy}{stred elipsy}{}}
\LANGes{
\Question{
Dada la ecuación de una elipse \( \frac{x^2}{48} +\frac{y^2}{36} = 1\). El punto \( (-2;5) \)  representa:}
{un punto que se encuentra en el interior de la elipse}{vértice de la elipse}{co-vértice de la elipse}{un punto que se encuentra fuera de la elipse}{el centro de la elipse}{}}
\End

\Data\ID{1103040101}
\Updated{2021-09-23 05:09:59}
\Level{A}
\SubArea{Conic Sections}
\LANGen{
\Question{
The equation of an ellipse is given by \( \frac{(x-2)^2}4+y^2=1 \).
Which of the pictures below shows the given ellipse in the  rectangular coordinate system?}
{}{}{}{}{}{}}
\LANGpl{
\Question{
Elipsa jest opisana równaniem \( \frac{(x-2)^2}4+y^2=1 \).
Które z poniższych rysunków pokazuje daną  elipsę w kartezjańskim układzie współrzędnych?}
{}{}{}{}{}{}}
\LANGcs{
\Question{
Elipsa je dána rovnicí \( \frac{(x-2)^2}4+y^2=1 \). 
Obrázek této elipsy v kartézské soustavě souřadnic je:}
{}{}{}{}{}{}}
\LANGsk{
\Question{
Elipsa je daná rovnicou \( \frac{(x-2)^2}4+y^2=1 \). 
Obrázok tejto elipsy v kartézskej sústave súradníc je:}
{}{}{}{}{}{}}
\LANGes{
\Question{
Dada la ecuación de una elipse \( \frac{(x-2)^2}4+y^2=1 \). ¿Cuál de las imágenes de abajo representa la elipse en el sistema de coordenadas cartesianas?}
{}{}{}{}{}{}}

\IMGanswer{1103040101a_0.pdf}
\IMGanswer{1103040101b_0.pdf}
\IMGanswer{1103040101c.pdf}
\IMGanswer{1103040101d.pdf}\End

\Data\ID{1103040102}
\Updated{2021-09-25 08:09:31}
\Level{A}
\SubArea{Conic Sections}
\IMG{1103040102.pdf}
\LANGen{
\Question{
The diagram of an ellipse in the rectangular coordinate system is shown in the picture. Find the standard form equation of this ellipse:}
{\( \frac{(x-3)^2}4+\frac{(y-3)^2}9=1 \)}{\( \frac{(x-3)^2}9+\frac{(y-3)^2}4=1 \)}{\( \frac{(x+3)^2}4+\frac{(y+3)^2}9=1 \)}{\( \frac{(x+3)^2}9+\frac{(y+3)^2}4=1 \)}{}{}}
\LANGpl{
\Question{
Rysunek przedstawia elipsę na kartezjańskim układzie współrzędnych. Wskaż postać kanoniczną równania danej elipsy.}
{\( \frac{(x-3)^2}4+\frac{(y-3)^2}9=1 \)}{\( \frac{(x-3)^2}9+\frac{(y-3)^2}4=1 \)}{\( \frac{(x+3)^2}4+\frac{(y+3)^2}9=1 \)}{\( \frac{(x+3)^2}9+\frac{(y+3)^2}4=1 \)}{}{}}
\LANGcs{
\Question{
Elipsa je dána obrázkem v kartézské soustavě souřadnic. Středový tvar rovnice této elipsy je:}
{\( \frac{(x-3)^2}4+\frac{(y-3)^2}9=1 \)}{\( \frac{(x-3)^2}9+\frac{(y-3)^2}4=1 \)}{\( \frac{(x+3)^2}4+\frac{(y+3)^2}9=1 \)}{\( \frac{(x+3)^2}9+\frac{(y+3)^2}4=1 \)}{}{}}
\LANGsk{
\Question{
Elipsa je daná obrázkom v kartézskej sústave súradníc. Stredový tvar rovnice tejto elipsy je:}
{\( \frac{(x-3)^2}4+\frac{(y-3)^2}9=1 \)}{\( \frac{(x-3)^2}9+\frac{(y-3)^2}4=1 \)}{\( \frac{(x+3)^2}4+\frac{(y+3)^2}9=1 \)}{\( \frac{(x+3)^2}9+\frac{(y+3)^2}4=1 \)}{}{}}
\LANGes{
\Question{
En la imagen se representa una elipse en el sistema de coordenadas cartesianas. encuentra su ecuación estándar .}
{\( \frac{(x-3)^2}4+\frac{(y-3)^2}9=1 \)}{\( \frac{(x-3)^2}9+\frac{(y-3)^2}4=1 \)}{\( \frac{(x+3)^2}4+\frac{(y+3)^2}9=1 \)}{\( \frac{(x+3)^2}9+\frac{(y+3)^2}4=1 \)}{}{}}
\End

\Data\ID{1103040103}
\Updated{2021-09-25 08:09:34}
\Level{A}
\SubArea{Conic Sections}
\IMG{1103040103.pdf}
\LANGen{
\Question{
The diagram of an ellipse in the rectangular coordinate system is shown in the picture.What is the semi-major axis of this ellipse?}
{The segment \( SN \)}{The segment \( SK \)}{The line \( EF \)}{The line \( KL \)}{}{}}
\LANGpl{
\Question{
Rysunek przedstawia elipsę na kartezjańskim układzie współrzędnych. Wskaż półoś wielką danej elipsy.}
{Odcinek \( SN \)}{Odcinek \( SK \)}{Prosta \( EF \)}{Prosta \( KL \)}{}{}}
\LANGcs{
\Question{
Elipsa je dána obrázkem v kartézské soustavě souřadnic. Hlavní poloosa této elipsy je:}
{úsečka \( SN \)}{úsečka \( SK \)}{přímka \( EF \)}{přímka \( KL \)}{}{}}
\LANGsk{
\Question{
Elipsa je daná obrázkom v kartézskej sústave súradníc. Hlavná poloos tejto elipsy je:}
{úsečka \( SN \)}{úsečka \( SK \)}{priamka \( EF \)}{priamka \( KL \)}{}{}}
\LANGes{
\Question{
En la imagen está representada una elipse en un sistema de coordenadas cartesianas. ¿Cuál es el semieje mayor de esta elipse?}
{El segmento \( SN \)}{El segmento\( SK \)}{La recta \( EF \)}{La recta \( KL \)}{}{}}
\End

\Data\ID{1103040106}
\Updated{2021-09-25 08:09:08}
\Level{A}
\SubArea{Conic Sections}
\IMG{1103040106.pdf}
\LANGen{
\Question{
A circle in the rectangular coordinate system is shown in the picture. What is the standard form equation of this circle?}
{\( x^2+(y-3)^2=4 \)}{\( x^2+(y+3)^2=4 \)}{\( x^2+(y-3)^2=2 \)}{\( x^2+(y+3)^2=2 \)}{}{}}
\LANGpl{
\Question{
Rysunek przedstawia okrąg  na kartezjańskim układzie współrzędnych. Wskaż postać kanoniczną równania danego  okręgu.}
{\( x^2+(y-3)^2=4 \)}{\( x^2+(y+3)^2=4 \)}{\( x^2+(y-3)^2=2 \)}{\( x^2+(y+3)^2=2 \)}{}{}}
\LANGcs{
\Question{
Kružnice je dána obrázkem v kartézské soustavě souřadnic. Středový tvar rovnice této kružnice je:}
{\( x^2+(y-3)^2=4 \)}{\( x^2+(y+3)^2=4 \)}{\( x^2+(y-3)^2=2 \)}{\( x^2+(y+3)^2=2 \)}{}{}}
\LANGsk{
\Question{
Kružnica je daná obrázkom v kartézskej sústave súradníc. Stredový tvar rovnice tejto kružnice je:}
{\( x^2+(y-3)^2=4 \)}{\( x^2+(y+3)^2=4 \)}{\( x^2+(y-3)^2=2 \)}{\( x^2+(y+3)^2=2 \)}{}{}}
\LANGes{
\Question{
En la imagen se representa una circunferencia en el sistema de coordenadas cartesianas. Identifica su ecuación estándar.}
{\( x^2+(y-3)^2=4 \)}{\( x^2+(y+3)^2=4 \)}{\( x^2+(y-3)^2=2 \)}{\( x^2+(y+3)^2=2 \)}{}{}}
\End

\Data\ID{1103040107}
\Updated{2021-09-25 08:09:25}
\Level{A}
\SubArea{Conic Sections}
\IMG{1103040107_0.pdf}
\LANGen{
\Question{
The picture below shows an ellipse drawn in the rectangular coordinate system. What is the semi-minor axis of this ellipse?}
{The segment \( SM \)}{The segment \( SL \)}{The line \( MN \)}{The line \( EF \)}{}{}}
\LANGpl{
\Question{
Rysunek przedstawia elipsę na kartezjańskim układzie współrzędnych. Wskaż półoś małą danej elipsy.}
{Odcinek \( SM \)}{Odcinek \( SL \)}{Prosta \( MN \)}{Prosta \( EF \)}{}{}}
\LANGcs{
\Question{
Elipsa je dána obrázkem v kartézské soustavě souřadnic. Vedlejší poloosa této elipsy je:}
{úsečka \( SM \)}{úsečka \( SL \)}{přímka \( MN \)}{přímka \( EF \)}{}{}}
\LANGsk{
\Question{
Elipsa je daná obrázkom v kartézskej sústave súradníc. Vedľajšia poloos tejto elipsy je:}
{úsečka \( SM \)}{úsečka \( SL \)}{priamka \( MN \)}{priamka \( EF \)}{}{}}
\LANGes{
\Question{
En la imagen se representa una elipse en el sistema de coordenadas cartesianas. Identifica el semieje menor de esta elipse.}
{El segmento \( SM \)}{El segmento \( SL \)}{La recta \( MN \)}{La recta \( EF \)}{}{}}
\End

\Data\ID{2010005908}
\Updated{2022-05-16 11:05:35}
\Level{A}
\SubArea{Conic Sections}
\LANGen{
\Question{
Which of the given points is one of the co-vertices on the minor axis of the ellipse  \(25x^{2} + 9y^{2} + 100x - 54y - 44 = 0\)?}
{\([-5;3]\)}{\([-2;0]\)}{\([-2;-2]\)}{\([3;1]\)}{}{}}
\LANGpl{
\Question{
Który z podanych punktów jest jednym z wierzchołków równoległych na małej osi elipsy \(25x^{2} + 9y^{2} + 100x - 54y - 44 = 0\)?}
{\([-5;3]\)}{\([-2;0]\)}{\([-2;-2]\)}{\([3;1]\)}{}{}}
\LANGcs{
\Question{
Který z bodů je vedlejším vrcholem elipsy  \(25x^{2} + 9y^{2} + 100x - 54y - 44 = 0\)?}
{\([-5;3]\)}{\([-2;0]\)}{\([-2;-2]\)}{\([3;1]\)}{}{}}
\LANGsk{
\Question{
Elipsa je daná rovnicou  \(25x^{2} + 9y^{2} + 100x - 54y - 44 = 0\). Jeden z jej vedľajších vrcholov má súradnice:}
{\([-5;3]\)}{\([-2;0]\)}{\([-2;-2]\)}{\([3;1]\)}{}{}}
\LANGes{
\Question{
Encuentra las coordenadas de uno de los vértices en el eje menor de la elipse \(25x^{2} + 9y^{2} + 100x - 54y - 44 = 0\).}
{\([-5;3]\)}{\([-2;0]\)}{\([-2;-2]\)}{\([3;1]\)}{}{}}
\End

\Data\ID{2010006007}
\Updated{2022-05-16 11:05:35}
\Level{A}
\SubArea{Conic Sections}
\LANGen{
\Question{
Consider the circle \(x^{2} + y^{2} +4x+ 6y + 10 = 0\).
Find the radius of the circle.}
{\(\sqrt{3}\)}{\(2\)}{\(3\)}{\(9\)}{}{}}
\LANGpl{
\Question{
Rozważmy okrąg \(x^{2} + y^{2} +4x+ 6y + 10 = 0\).
Znajdź promień okręgu.}
{\(\sqrt{3}\)}{\(2\)}{\(3\)}{\(9\)}{}{}}
\LANGcs{
\Question{
Je dána kružnice \(x^{2} + y^{2} +4x+ 6y + 10 = 0\).
Určete její poloměr.}
{\(\sqrt{3}\)}{\(2\)}{\(3\)}{\(9\)}{}{}}
\LANGsk{
\Question{
Je daná kružnica \(x^{2} + y^{2} +4x+ 6y + 10 = 0\).
Polomer tejto kružnice je rovný:}
{\(\sqrt{3}\)}{\(2\)}{\(3\)}{\(9\)}{}{}}
\LANGes{
\Question{
Consideremos la circunferencia \(x^{2} + y^{2} +4x+ 6y + 10 = 0\).
Encuentra su radio.}
{\(\sqrt{3}\)}{\(2\)}{\(3\)}{\(9\)}{}{}}
\End

\Data\ID{2010006301}
\Updated{2022-05-16 11:05:39}
\Level{A}
\SubArea{Conic Sections}
\LANGen{
\Question{
The equation of an ellipse is given by \( 9x^2+5y^2+18x-30y+9=0 \). The length of its semi-minor axis is:}
{\( \sqrt5\)}{\(5\)}{\(9\)}{\(3\)}{\(2\)}{}}
\LANGpl{
\Question{
Równanie elipsy jest podane przez \( 9x^2+5y^2+18x-30y+9=0 \). Długość jej półosi małej to:}
{\( \sqrt5\)}{\(5\)}{\(9\)}{\(3\)}{\(2\)}{}}
\LANGcs{
\Question{
Určete délku vedlejší poloosy elipsy dané rovnicí \( 9x^2+5y^2+18x-30y+9=0 \).}
{\( \sqrt5\)}{\(5\)}{\(9\)}{\(3\)}{\(2\)}{}}
\LANGsk{
\Question{
Elipsa je daná rovnicou \( 9x^2+5y^2+18x-30y+9=0 \). Dĺžka vedľajšej polosi tejto elipsy je:}
{\( \sqrt5\)}{\(5\)}{\(9\)}{\(3\)}{\(2\)}{}}
\LANGes{
\Question{
Dada la ecuación de la elipse \( 9x^2+5y^2+18x-30y+9=0 \). ¿Cuánto mide el semieje menor?}
{\( \sqrt5\)}{\(5\)}{\(9\)}{\(3\)}{\(2\)}{}}
\End

\Data\ID{9000091201}
\Updated{2021-09-25 08:09:39}
\Level{A}
\SubArea{Conic Sections}
\LANGen{
\Question{
Consider the circle \(k\colon x^{2} - 4x + y^{2} + 6y + 12 = 0\).
Find the radius of the circle.}
{\(1\)}{\(2\)}{\(3\)}{\(4\)}{}{}}
\LANGpl{
\Question{
Dany jest okrąg  \(k\colon x^{2} - 4x + y^{2} + 6y + 12 = 0\).
Wskaż promień okręgu.}
{\(1\)}{\(2\)}{\(3\)}{\(4\)}{}{}}
\LANGcs{
\Question{
Je dána kružnice \(k\colon x^{2} - 4x + y^{2} + 6y + 12 = 0\).
Poloměr této kružnice je roven:}
{\(1\)}{\(2\)}{\(3\)}{\(4\)}{}{}}
\LANGsk{
\Question{
Je daná kružnica \(k\colon x^{2} - 4x + y^{2} + 6y + 12 = 0\).
Polomer tejto kružnice je rovný:}
{\(1\)}{\(2\)}{\(3\)}{\(4\)}{}{}}
\LANGes{
\Question{
Consideremos la circunferencia \(k\colon x^{2} - 4x + y^{2} + 6y + 12 = 0\).
Halla su radio.}
{\(1\)}{\(2\)}{\(3\)}{\(4\)}{}{}}
\End

\Data\ID{9000091202}
\Updated{2021-09-25 08:09:08}
\Level{A}
\SubArea{Conic Sections}
\LANGen{
\Question{
Consider the circle \(k\colon x^{2} + 6x + y^{2} + 2y + 6 = 0\).
Find the radius of the circle.}
{\(2\)}{\(\sqrt{2}\)}{\(3\)}{\(4\)}{}{}}
\LANGpl{
\Question{
Dany jest okrąg  \(k\colon x^{2} + 6x + y^{2} + 2y + 6 = 0\).
 Wskaż promień okręgu.}
{\(2\)}{\(\sqrt{2}\)}{\(3\)}{\(4\)}{}{}}
\LANGcs{
\Question{
Je dána kružnice \(k\colon x^{2} + 6x + y^{2} + 2y + 6 = 0\).
Poloměr této kružnice je roven:}
{\(2\)}{\(\sqrt{2}\)}{\(3\)}{\(4\)}{}{}}
\LANGsk{
\Question{
Je daná kružnica \(k\colon x^{2} + 6x + y^{2} + 2y + 6 = 0\).
Polomer tejto kružnice je rovný:}
{\(2\)}{\(\sqrt{2}\)}{\(3\)}{\(4\)}{}{}}
\LANGes{
\Question{
Consideremos la circunferencia \(k\colon x^{2} + 6x + y^{2} + 2y + 6 = 0\).
Encuentra su radio.}
{\(2\)}{\(\sqrt{2}\)}{\(3\)}{\(4\)}{}{}}
\End

\Data\ID{9000091203}
\Updated{2022-04-23 05:04:26}
\Level{A}
\SubArea{Conic Sections}
\LANGen{
\Question{
Consider the circle \(x^{2} - 2x + y^{2} - 6y + 8 = 0\).
Find the radius of the circle.}
{\(\sqrt{2}\)}{\(2\)}{\(3\)}{\(4\)}{}{}}
\LANGpl{
\Question{
Dany jest okrąg \(x^{2} - 2x + y^{2} - 6y + 8 = 0\).
Wskaż promień okręgu.}
{\(\sqrt{2}\)}{\(2\)}{\(3\)}{\(4\)}{}{}}
\LANGcs{
\Question{
Je dána kružnice \(x^{2} - 2x + y^{2} - 6y + 8 = 0\).
Poloměr této kružnice je roven:}
{\(\sqrt{2}\)}{\(2\)}{\(3\)}{\(4\)}{}{}}
\LANGsk{
\Question{
Je daná kružnica \(x^{2} - 2x + y^{2} - 6y + 8 = 0\).
Polomer tejto kružnice je rovný:}
{\(\sqrt{2}\)}{\(2\)}{\(3\)}{\(4\)}{}{}}
\LANGes{
\Question{
Consideremos la circunferencia \(x^{2} - 2x + y^{2} - 6y + 8 = 0\).
Encuentra su radio.}
{\(\sqrt{2}\)}{\(2\)}{\(3\)}{\(4\)}{}{}}
\End

\Data\ID{9000091204}
\Updated{2021-10-11 11:10:48}
\Level{A}
\SubArea{Conic Sections}
\LANGen{
\Question{
Consider the circle \(k\colon x^{2} + 4x + y^{2} - 8y + 11 = 0\).
Find the radius of the circle.}
{\(3\)}{\(1\)}{\(2\)}{\(4\)}{}{}}
\LANGpl{
\Question{
Dany jest okrąg \(k\colon x^{2} + 4x + y^{2} - 8y + 11 = 0\).
Wskaż promień okręgu.}
{\(3\)}{\(1\)}{\(2\)}{\(4\)}{}{}}
\LANGcs{
\Question{
Je dána kružnice \(k\colon x^{2} + 4x + y^{2} - 8y + 11 = 0\).
Poloměr této kružnice je roven:}
{\(3\)}{\(1\)}{\(2\)}{\(4\)}{}{}}
\LANGsk{
\Question{
Je daná kružnica \(k\colon x^{2} + 4x + y^{2} - 8y + 11 = 0\).
Polomer tejto kružnice je rovný:}
{\(3\)}{\(1\)}{\(2\)}{\(4\)}{}{}}
\LANGes{
\Question{
Consideremos la circunferencia \(k\colon x^{2} + 4x + y^{2} - 8y + 11 = 0\).
Encuentra su radio.}
{\(3\)}{\(1\)}{\(2\)}{\(4\)}{}{}}
\End

\Data\ID{9000091205}
\Updated{2021-10-11 11:10:17}
\Level{A}
\SubArea{Conic Sections}
\LANGen{
\Question{
Consider the circle \(k\colon x^{2} - 10x + y^{2} + 10y + 34 = 0\).
Find the radius of the circle.}
{\(4\)}{\(2\)}{\(3\)}{\(1\)}{}{}}
\LANGpl{
\Question{
Dany jest okrąg \(k\colon x^{2} - 10x + y^{2} + 10y + 34 = 0\).
Wskaż promień okręgu.}
{\(4\)}{\(2\)}{\(3\)}{\(1\)}{}{}}
\LANGcs{
\Question{
Je dána kružnice \(k\colon x^{2} - 10x + y^{2} + 10y + 34 = 0\).
Poloměr této kružnice je roven:}
{\(4\)}{\(2\)}{\(3\)}{\(1\)}{}{}}
\LANGsk{
\Question{
Je daná kružnica \(k\colon x^{2} - 10x + y^{2} + 10y + 34 = 0\).
Polomer tejto kružnice je rovný:}
{\(4\)}{\(2\)}{\(3\)}{\(1\)}{}{}}
\LANGes{
\Question{
Consideremos la circunferencia \(k\colon x^{2} - 10x + y^{2} + 10y + 34 = 0\).
Encuentra su radio.}
{\(4\)}{\(2\)}{\(3\)}{\(1\)}{}{}}
\End

\Data\ID{9000091206}
\Updated{2021-10-11 11:10:40}
\Level{A}
\SubArea{Conic Sections}
\LANGen{
\Question{
Consider the circle \(k\colon x^{2} - 4x + y^{2} + 6y + 11 = 0\).
Find the center of the circle.}
{\(S = [2;-3]\)}{\(S = [2;3]\)}{\(S = [-2;3]\)}{\(S = [-2;-3]\)}{}{}}
\LANGpl{
\Question{
Dany jest okrąg \(k\colon x^{2} - 4x + y^{2} + 6y + 11 = 0\).
Wyznacz środek okręgu.}
{\(S = [2;-3]\)}{\(S = [2;3]\)}{\(S = [-2;3]\)}{\(S = [-2;-3]\)}{}{}}
\LANGcs{
\Question{
Je dána kružnice \(k\colon x^{2} - 4x + y^{2} + 6y + 11 = 0\).
Střed této kružnice je:}
{\(S = [2;-3]\)}{\(S = [2;3]\)}{\(S = [-2;3]\)}{\(S = [-2;-3]\)}{}{}}
\LANGsk{
\Question{
Je daná kružnica \(k\colon x^{2} - 4x + y^{2} + 6y + 11 = 0\).
Stred tejto kružnice je:}
{\(S = [2;-3]\)}{\(S = [2;3]\)}{\(S = [-2;3]\)}{\(S = [-2;-3]\)}{}{}}
\LANGes{
\Question{
Consideremos la circunferencia \(k\colon x^{2} - 4x + y^{2} + 6y + 11 = 0\).
Encuentra su radio.}
{\(S = [2;-3]\)}{\(S = [2;3]\)}{\(S = [-2;3]\)}{\(S = [-2;-3]\)}{}{}}
\End

\Data\ID{9000091207}
\Updated{2021-10-11 11:10:05}
\Level{A}
\SubArea{Conic Sections}
\LANGen{
\Question{
Consider the circle \(k\colon x^{2} - 6x + y^{2} + 2y + 6 = 0\).
Find the center of the circle.}
{\(S = [3;-1]\)}{\(S = [-3;-1]\)}{\(S = [3;1]\)}{\(S = [-3;1]\)}{}{}}
\LANGpl{
\Question{
Dany jest okrąg \(k\colon x^{2} - 6x + y^{2} + 2y + 6 = 0\).
Wyznacz środek okręgu.}
{\(S = [3;-1]\)}{\(S = [-3;-1]\)}{\(S = [3;1]\)}{\(S = [-3;1]\)}{}{}}
\LANGcs{
\Question{
Je dána kružnice \(k\colon x^{2} - 6x + y^{2} + 2y + 6 = 0\).
Střed této kružnice je:}
{\(S = [3;-1]\)}{\(S = [-3;-1]\)}{\(S = [3;1]\)}{\(S = [-3;1]\)}{}{}}
\LANGsk{
\Question{
Je daná kružnica \(k\colon x^{2} - 6x + y^{2} + 2y + 6 = 0\).
Stred tejto kružnice je:}
{\(S = [3;-1]\)}{\(S = [-3;-1]\)}{\(S = [3;1]\)}{\(S = [-3;1]\)}{}{}}
\LANGes{
\Question{
Consideremos la circunferencia \(k\colon x^{2} - 6x + y^{2} + 2y + 6 = 0\).
Encuentra su radio.}
{\(S = [3;-1]\)}{\(S = [-3;-1]\)}{\(S = [3;1]\)}{\(S = [-3;1]\)}{}{}}
\End

\Data\ID{9000091208}
\Updated{2021-10-11 11:10:34}
\Level{A}
\SubArea{Conic Sections}
\LANGen{
\Question{
Consider the circle \(k\colon x^{2} + 2x + y^{2} - 4y + 2 = 0\).
Find the center of the circle.}
{\(S = [-1;2]\)}{\(S = [-1;-2]\)}{\(S = [1;-2]\)}{\(S = [1;2]\)}{}{}}
\LANGpl{
\Question{
Dany jest okrąg \(k\colon x^{2} + 2x + y^{2} - 4y + 2 = 0\).
Wyznacz środek okręgu.}
{\(S = [-1;2]\)}{\(S = [-1;-2]\)}{\(S = [1;-2]\)}{\(S = [1;2]\)}{}{}}
\LANGcs{
\Question{
Je dána kružnice \(k\colon x^{2} + 2x + y^{2} - 4y + 2 = 0\).
Střed této kružnice je:}
{\(S = [-1;2]\)}{\(S = [-1;-2]\)}{\(S = [1;-2]\)}{\(S = [1;2]\)}{}{}}
\LANGsk{
\Question{
Je daná kružnica \(k\colon x^{2} + 2x + y^{2} - 4y + 2 = 0\).
Stred tejto kružnice je:}
{\(S = [-1;2]\)}{\(S = [-1;-2]\)}{\(S = [1;-2]\)}{\(S = [1;2]\)}{}{}}
\LANGes{
\Question{
Consideremos la circunferencia \(k\colon x^{2} + 2x + y^{2} - 4y + 2 = 0\).
Encuentra su centro.}
{\(S = [-1;2]\)}{\(S = [-1;-2]\)}{\(S = [1;-2]\)}{\(S = [1;2]\)}{}{}}
\End

\Data\ID{9000091209}
\Updated{2021-10-11 11:10:13}
\Level{A}
\SubArea{Conic Sections}
\LANGen{
\Question{
Consider the circle \(k\colon x^{2} - 6x + y^{2} + 4y + 9 = 0\).
Find the center of the circle.}
{\(S = [3;-2]\)}{\(S = [-3;-2]\)}{\(S = [3;2]\)}{\(S = [-3;2]\)}{}{}}
\LANGpl{
\Question{
Dany jest okrąg \(k\colon x^{2} - 6x + y^{2} + 4y + 9 = 0\).
Wyznacz środek okręgu.}
{\(S = [3;-2]\)}{\(S = [-3;-2]\)}{\(S = [3;2]\)}{\(S = [-3;2]\)}{}{}}
\LANGcs{
\Question{
Je dána kružnice \(k\colon x^{2} - 6x + y^{2} + 4y + 9 = 0\).
Střed této kružnice je:}
{\(S = [3;-2]\)}{\(S = [-3;-2]\)}{\(S = [3;2]\)}{\(S = [-3;2]\)}{}{}}
\LANGsk{
\Question{
Je daná kružnica \(k\colon x^{2} - 6x + y^{2} + 4y + 9 = 0\).
Stred tejto kružnice je:}
{\(S = [3;-2]\)}{\(S = [-3;-2]\)}{\(S = [3;2]\)}{\(S = [-3;2]\)}{}{}}
\LANGes{
\Question{
Consideremos la circunferencia \(k\colon x^{2} - 6x + y^{2} + 4y + 9 = 0\).
Encuentra su centro .}
{\(S = [3;-2]\)}{\(S = [-3;-2]\)}{\(S = [3;2]\)}{\(S = [-3;2]\)}{}{}}
\End

\Data\ID{9000091210}
\Updated{2021-10-11 11:10:46}
\Level{A}
\SubArea{Conic Sections}
\LANGen{
\Question{
Consider the circle \(k\colon x^{2} + 8x + y^{2} - 2y + 12 = 0\).
Find the center of the circle.}
{\(S = [-4;1]\)}{\(S = [-4;-1]\)}{\(S = [4;1]\)}{\(S = [4;-1]\)}{}{}}
\LANGpl{
\Question{
Dany jest okrąg  \(k\colon x^{2} + 8x + y^{2} - 2y + 12 = 0\). Wyznacz środek okręgu.}
{\(S = [-4;1]\)}{\(S = [-4;-1]\)}{\(S = [4;1]\)}{\(S = [4;-1]\)}{}{}}
\LANGcs{
\Question{
Je dána kružnice \(k\colon x^{2} + 8x + y^{2} - 2y + 12 = 0\).
Střed této kružnice je:}
{\(S = [-4;1]\)}{\(S = [-4;-1]\)}{\(S = [4;1]\)}{\(S = [4;-1]\)}{}{}}
\LANGsk{
\Question{
Je daná kružnica \(k\colon x^{2} + 8x + y^{2} - 2y + 12 = 0\).
Stred tejto kružnice je:}
{\(S = [-4;1]\)}{\(S = [-4;-1]\)}{\(S = [4;1]\)}{\(S = [4;-1]\)}{}{}}
\LANGes{
\Question{
Consideremos la circunferencia \(k\colon x^{2} + 8x + y^{2} - 2y + 12 = 0\).
Encuentra su centro.}
{\(S = [-4;1]\)}{\(S = [-4;-1]\)}{\(S = [4;1]\)}{\(S = [4;-1]\)}{}{}}
\End

\Data\ID{9000149701}
\Updated{2020-08-30 07:08:30}
\Level{A}
\SubArea{Conic Sections}
\LANGen{
\Question{
Find the center of the following circle. 
\[ 
 x^{2} + y^{2} - 4x + 6y - 12 = 0
\]}
{\([2;-3]\)}{\([-2;3]\)}{\([2;3]\)}{\([-2;-3]\)}{}{}}
\LANGpl{
\Question{
Wyznacz środek okręgu
\[ 
 x^{2} + y^{2} - 4x + 6y - 12 = 0
\]}
{\([2;-3]\)}{\([-2;3]\)}{\([2;3]\)}{\([-2;-3]\)}{}{}}
\LANGcs{
\Question{
Kružnice je dána rovnicí \(x^{2} + y^{2} - 4x + 6y - 12 = 0\).
Její střed má souřadnice:}
{\([2;-3]\)}{\([-2;3]\)}{\([2;3]\)}{\([-2;-3]\)}{}{}}
\LANGsk{
\Question{
Kružnice je daná rovnicou  \(x^{2} + y^{2} - 4x + 6y - 12 = 0\).
Jej stred má súradnice:}
{\([2;-3]\)}{\([-2;3]\)}{\([2;3]\)}{\([-2;-3]\)}{}{}}
\LANGes{
\Question{
Halla el centro de la siguiente circunferencia.  
\[ 
 x^{2} + y^{2} - 4x + 6y - 12 = 0
\]}
{\([2;-3]\)}{\([-2;3]\)}{\([2;3]\)}{\([-2;-3]\)}{}{}}
\End

\Data\ID{9000149702}
\Updated{2020-08-30 07:08:59}
\Level{A}
\SubArea{Conic Sections}
\LANGen{
\Question{
Find the center of the following circle. 
\[ 
 x^{2} + y^{2} + 2x - 8y + 13 = 0
\]}
{\([-1;4]\)}{\([-1;-4]\)}{\([1;4]\)}{\([1;-4]\)}{}{}}
\LANGpl{
\Question{
Wyznacz środek okręgu. 
\[ 
 x^{2} + y^{2} + 2x - 8y + 13 = 0
\]}
{\([-1;4]\)}{\([-1;-4]\)}{\([1;4]\)}{\([1;-4]\)}{}{}}
\LANGcs{
\Question{
Kružnice je dána rovnicí \(x^{2} + y^{2} + 2x - 8y + 13 = 0\).
Její střed má souřadnice:}
{\([-1;4]\)}{\([-1;-4]\)}{\([1;4]\)}{\([1;-4]\)}{}{}}
\LANGsk{
\Question{
Kružnice je daná rovnicou \(x^{2} + y^{2} + 2x - 8y + 13 = 0\).
Jej stred má súradnice:}
{\([-1;4]\)}{\([-1;-4]\)}{\([1;4]\)}{\([1;-4]\)}{}{}}
\LANGes{
\Question{
Halla el centro de la siguiente circunferencia.  
\[ 
 x^{2} + y^{2} + 2x - 8y + 13 = 0
\]}
{\([-1;4]\)}{\([-1;-4]\)}{\([1;4]\)}{\([1;-4]\)}{}{}}
\End

\Data\ID{9000149703}
\Updated{2020-08-30 07:08:28}
\Level{A}
\SubArea{Conic Sections}
\LANGen{
\Question{
Find the center of the following circle. 
\[ 
 x^{2} + y^{2} - 10x - 2y + 10 = 0
\]}
{\([5;1]\)}{\([5;-1]\)}{\([-5;1]\)}{\([-5;-1]\)}{}{}}
\LANGpl{
\Question{
Wyznacz środek okręgu. 
\[ 
 x^{2} + y^{2} - 10x - 2y + 10 = 0
\]}
{\([5;1]\)}{\([5;-1]\)}{\([-5;1]\)}{\([-5;-1]\)}{}{}}
\LANGcs{
\Question{
Kružnice je dána rovnicí \(x^{2} + y^{2} - 10x - 2y + 10 = 0\).
Její střed má souřadnice:}
{\([5;1]\)}{\([5;-1]\)}{\([-5;1]\)}{\([-5;-1]\)}{}{}}
\LANGsk{
\Question{
Kružnica je daná rovnicou \(x^{2} + y^{2} - 10x - 2y + 10 = 0\).
Jej stred má súradnice:}
{\([5;1]\)}{\([5;-1]\)}{\([-5;1]\)}{\([-5;-1]\)}{}{}}
\LANGes{
\Question{
Halla el centro de la siguiente circunferencia.  
\[ 
 x^{2} + y^{2} - 10x - 2y + 10 = 0
\]}
{\([5;1]\)}{\([5;-1]\)}{\([-5;1]\)}{\([-5;-1]\)}{}{}}
\End

\Data\ID{9000149704}
\Updated{2020-08-30 08:08:06}
\Level{A}
\SubArea{Conic Sections}
\LANGen{
\Question{
Find the center of the following ellipse. 
\[ 
 9x^{2} + 4y^{2} + 54x - 32y + 109 = 0
\]}
{\([-3;4]\)}{\([-3;-4]\)}{\([3;4]\)}{\([3;-4]\)}{}{}}
\LANGpl{
\Question{
Wyznacz środek elipsy.
\[ 
 9x^{2} + 4y^{2} + 54x - 32y + 109 = 0
\]}
{\([-3;4]\)}{\([-3;-4]\)}{\([3;4]\)}{\([3;-4]\)}{}{}}
\LANGcs{
\Question{
Elipsa je dána rovnicí \(9x^{2} + 4y^{2} + 54x - 32y + 109 = 0\).
Její střed má souřadnice:}
{\([-3;4]\)}{\([-3;-4]\)}{\([3;4]\)}{\([3;-4]\)}{}{}}
\LANGsk{
\Question{
Elipsa je daná rovnicou \(9x^{2} + 4y^{2} + 54x - 32y + 109 = 0\).
Jej stred má súradnice:}
{\([-3;4]\)}{\([-3;-4]\)}{\([3;4]\)}{\([3;-4]\)}{}{}}
\LANGes{
\Question{
Halla el centro de la siguiente elipse.
\[ 
 9x^{2} + 4y^{2} + 54x - 32y + 109 = 0
\]}
{\([-3;4]\)}{\([-3;-4]\)}{\([3;4]\)}{\([3;-4]\)}{}{}}
\End

\Data\ID{9000149705}
\Updated{2020-08-30 08:08:40}
\Level{A}
\SubArea{Conic Sections}
\LANGen{
\Question{
Find the center of the following ellipse. 
\[ 
 16x^{2} + 9y^{2} - 32x - 54y - 47 = 0
\]}
{\([1;3]\)}{\([1;-3]\)}{\([-1;3]\)}{\([-1;-3]\)}{}{}}
\LANGpl{
\Question{
Wyznacz środek elipsy.
\[ 
 16x^{2} + 9y^{2} - 32x - 54y - 47 = 0
\]}
{\([1;3]\)}{\([1;-3]\)}{\([-1;3]\)}{\([-1;-3]\)}{}{}}
\LANGcs{
\Question{
Elipsa je dána rovnicí \(16x^{2} + 9y^{2} - 32x - 54y - 47 = 0\).
Její střed má souřadnice:}
{\([1;3]\)}{\([1;-3]\)}{\([-1;3]\)}{\([-1;-3]\)}{}{}}
\LANGsk{
\Question{
Elipsa je daná rovnicou \(16x^{2} + 9y^{2} - 32x - 54y - 47 = 0\).
Jej stred má súradnice:}
{\([1;3]\)}{\([1;-3]\)}{\([-1;3]\)}{\([-1;-3]\)}{}{}}
\LANGes{
\Question{
Halla el centro de la siguiente elipse.  
\[ 
 16x^{2} + 9y^{2} - 32x - 54y - 47 = 0
\]}
{\([1;3]\)}{\([1;-3]\)}{\([-1;3]\)}{\([-1;-3]\)}{}{}}
\End

\Data\ID{9000168701}
\Updated{2022-04-23 05:04:20}
\Level{A}
\SubArea{Conic Sections}
\LANGen{
\Question{
Which of the given points is one of the vertices on the major axis of the ellipse  \(9x^{2} + 16y^{2} - 18x + 96y + 9 = 0\)?}
{\([5;-3]\)}{\([4;-3]\)}{\([1;1]\)}{\([1;0]\)}{}{}}
\LANGpl{
\Question{
Dana jest elipsa  \(9x^{2} + 16y^{2} - 18x + 96y + 9 = 0\),
wskaż współrzędne jednego z ich wierzchołków na głównej osi.}
{\([5;-3]\)}{\([4;-3]\)}{\([1;1]\)}{\([1;0]\)}{}{}}
\LANGcs{
\Question{
Elipsa je dána rovnicí \(9x^{2} + 16y^{2} - 18x + 96y + 9 = 0\).
Jeden z jejich hlavních vrcholů má souřadnice:}
{\([5;-3]\)}{\([4;-3]\)}{\([1;1]\)}{\([1;0]\)}{}{}}
\LANGsk{
\Question{
Elipsa je daná rovnicou \(9x^{2} + 16y^{2} - 18x + 96y + 9 = 0\).
Jeden z ich hlavných vrcholov má súradnice:}
{\([5;-3]\)}{\([4;-3]\)}{\([1;1]\)}{\([1;0]\)}{}{}}
\LANGes{
\Question{
Dada la ecuación de elipse \(9x^{2} + 16y^{2} - 18x + 96y + 9 = 0\),
encuentra las coordenadas de uno de los vértices en el eje mayor.}
{\([5;-3]\)}{\([4;-3]\)}{\([1;1]\)}{\([1;0]\)}{}{}}
\End

\Data\ID{9000168702}
\Updated{2020-11-07 03:11:25}
\Level{A}
\SubArea{Conic Sections}
\LANGen{
\Question{
Given the ellipse \(4x^{2} + 9y^{2} + 16x - 18y - 11 = 0\),
find the coordinates of the vertex on the minor axis.}
{\([-2;3]\)}{\([-2;4]\)}{\([0;1]\)}{\([1;1]\)}{}{}}
\LANGpl{
\Question{
Dana jest elipsa  \(4x^{2} + 9y^{2} + 16x - 18y - 11 = 0\),
wskaż współrzędne wierzchołka na osi małej.}
{\([-2;3]\)}{\([-2;4]\)}{\([0;1]\)}{\([1;1]\)}{}{}}
\LANGcs{
\Question{
Elipsa je dána rovnicí \(4x^{2} + 9y^{2} + 16x - 18y - 11 = 0\).
Její vedlejší vrchol má souřadnice:}
{\([-2;3]\)}{\([-2;4]\)}{\([0;1]\)}{\([1;1]\)}{}{}}
\LANGsk{
\Question{
Elipsa je daná rovnicou \(4x^{2} + 9y^{2} + 16x - 18y - 11 = 0\).
Jej vedľajší vrchol má súradnice:}
{\([-2;3]\)}{\([-2;4]\)}{\([0;1]\)}{\([1;1]\)}{}{}}
\LANGes{
\Question{
Dada la ecuación de elipse \(4x^{2} + 9y^{2} + 16x - 18y - 11 = 0\),
encuentra las coordenadas del vértice en el eje menor.}
{\([-2;3]\)}{\([-2;4]\)}{\([0;1]\)}{\([1;1]\)}{}{}}
\End

\Data\ID{9000168703}
\Updated{2020-11-07 04:11:03}
\Level{A}
\SubArea{Conic Sections}
\LANGen{
\Question{
Given the ellipse \(25x^{2} + 9y^{2} - 150x + 18y + 9 = 0\),
find the vertex on the major axis.}
{\([3;4]\)}{\([3;2]\)}{\([8;-1]\)}{\([6;-1]\)}{}{}}
\LANGpl{
\Question{
Dana jest elipsa \(25x^{2} + 9y^{2} - 150x + 18y + 9 = 0\),
wskaż współrzędne wierzchołka na osi głównej.}
{\([3;4]\)}{\([3;2]\)}{\([8;-1]\)}{\([6;-1]\)}{}{}}
\LANGcs{
\Question{
Elipsa je dána rovnicí \(25x^{2} + 9y^{2} - 150x + 18y + 9 = 0\).
Její hlavní vrchol má souřadnice:}
{\([3;4]\)}{\([3;2]\)}{\([8;-1]\)}{\([6;-1]\)}{}{}}
\LANGsk{
\Question{
Elipsa je daná rovnicou \(25x^{2} + 9y^{2} - 150x + 18y + 9 = 0\).
Jej hlavný vrchol má súradnice:}
{\([3;4]\)}{\([3;2]\)}{\([8;-1]\)}{\([6;-1]\)}{}{}}
\LANGes{
\Question{
Dada la elipse \(25x^{2} + 9y^{2} - 150x + 18y + 9 = 0\),
encuentra el vértice del eje mayor.}
{\([3;4]\)}{\([3;2]\)}{\([8;-1]\)}{\([6;-1]\)}{}{}}
\End

\Data\ID{9000168704}
\Updated{2020-11-07 04:11:21}
\Level{A}
\SubArea{Conic Sections}
\LANGen{
\Question{
Given the ellipse \(9x^{2} + 4y^{2} + 36x - 24y + 36 = 0\),
find the vertex on the minor axis.}
{\([-4;3]\)}{\([-5;3]\)}{\([-2;0]\)}{\([-2;1]\)}{}{}}
\LANGpl{
\Question{
Dana jest elipsa \(9x^{2} + 4y^{2} + 36x - 24y + 36 = 0\),
wskaż współrzędne wierzchołka na osi małej.}
{\([-4;3]\)}{\([-5;3]\)}{\([-2;0]\)}{\([-2;1]\)}{}{}}
\LANGcs{
\Question{
Elipsa je dána rovnicí \(9x^{2} + 4y^{2} + 36x - 24y + 36 = 0\).
Její vedlejší vrchol má souřadnice:}
{\([-4;3]\)}{\([-5;3]\)}{\([-2;0]\)}{\([-2;1]\)}{}{}}
\LANGsk{
\Question{
Elipsa je daná rovnicou \(9x^{2} + 4y^{2} + 36x - 24y + 36 = 0\).
Jej vedľajší vrchol má súradnice:}
{\([-4;3]\)}{\([-5;3]\)}{\([-2;0]\)}{\([-2;1]\)}{}{}}
\LANGes{
\Question{
Dada la elipse \(9x^{2} + 4y^{2} + 36x - 24y + 36 = 0\),
encuentra el vértice del eje menor.}
{\([-4;3]\)}{\([-5;3]\)}{\([-2;0]\)}{\([-2;1]\)}{}{}}
\End

\Data\ID{9000168705}
\Updated{2020-11-07 04:11:36}
\Level{A}
\SubArea{Conic Sections}
\LANGen{
\Question{
Given the ellipse \(16x^{2} + 9y^{2} + 32x - 36y - 92 = 0\),
find the vertex on the major axis.}
{\([-1;-2]\)}{\([-1;-1]\)}{\([3;2]\)}{\([2;2]\)}{}{}}
\LANGpl{
\Question{
Dana jest elipsa \(16x^{2} + 9y^{2} + 32x - 36y - 92 = 0\),
wskaż współrzędne wierzchołka na osi głównej.}
{\([-1;-2]\)}{\([-1;-1]\)}{\([3;2]\)}{\([2;2]\)}{}{}}
\LANGcs{
\Question{
Elipsa je dána rovnicí \(16x^{2} + 9y^{2} + 32x - 36y - 92 = 0\).
Její hlavní vrchol má souřadnice:}
{\([-1;-2]\)}{\([-1;-1]\)}{\([3;2]\)}{\([2;2]\)}{}{}}
\LANGsk{
\Question{
Elipsa je daná rovnicou \(16x^{2} + 9y^{2} + 32x - 36y - 92 = 0\).
Jej hlavný vrchol má súradnice:}
{\([-1;-2]\)}{\([-1;-1]\)}{\([3;2]\)}{\([2;2]\)}{}{}}
\LANGes{
\Question{
Dada la elipse \(16x^{2} + 9y^{2} + 32x - 36y - 92 = 0\),
encuentra el vértice del eje mayor.}
{\([-1;-2]\)}{\([-1;-1]\)}{\([3;2]\)}{\([2;2]\)}{}{}}
\End

\Data\ID{1003020404}
\Updated{2022-08-25 10:08:09}
\Level{B}
\SubArea{Conic Sections}
\LANGen{
\Question{
The equation of a hyperbola is given by \( 9x^2-16y^2-108x+96y+36=0\). The length of its semi-major axis is:}
{\( 4 \)}{\( 16 \)}{\( 3 \)}{\( 9 \)}{\( 25 \)}{}}
\LANGpl{
\Question{
Dane jest równanie hiperboli \( 9x^2-16y^2-108x+96y+36=0\). Długość półosi wielkiej to:}
{\( 4 \)}{\( 16 \)}{\( 3 \)}{\( 9 \)}{\( 25 \)}{}}
\LANGcs{
\Question{
Hyperbola je dána rovnicí \( 9x^2-16y^2-108x+96y+36=0\). Délka hlavní poloosy této hyperboly je:}
{\( 4 \)}{\( 16 \)}{\( 3 \)}{\( 9 \)}{\( 25 \)}{}}
\LANGsk{
\Question{
Hyperbola je daná rovnicou \( 9x^2-16y^2-108x+96y+36=0\). Dĺžka hlavnej polosi  tejto hyperboly je:}
{\( 4 \)}{\( 16 \)}{\( 3 \)}{\( 9 \)}{\( 25 \)}{}}
\LANGes{
\Question{
Dada la ecuación de la hipérbola \( 9x^2-16y^2-108x+96y+36=0\), ¿Cuánto mide su semieje mayor?}
{\( 4 \)}{\( 16 \)}{\( 3 \)}{\( 9 \)}{\( 25 \)}{}}
\End

\Data\ID{1003020405}
\Updated{2022-08-25 10:08:09}
\Level{B}
\SubArea{Conic Sections}
\LANGen{
\Question{
A hyperbola has two asymptotes with the equations \( y=\frac13(x+2) \) and \(y=-\frac13(x+2)\). What are  the coordinates of its centre?}
{\( [-2;0] \)}{\( [2;0] \)}{\( \left[-\frac13;0\right] \)}{\( [0;2] \)}{\( \left[\frac13;2\right] \)}{}}
\LANGpl{
\Question{
Hiperbola ma dwie asymptoty o równaniach  \( y=\frac13(x+2) \) i \(y=-\frac13(x+2)\). Wyznacz współrzędne jej środka.}
{\( [-2;0] \)}{\( [2;0] \)}{\( \left[-\frac13;0\right] \)}{\( [0;2] \)}{\( \left[\frac13;2\right] \)}{}}
\LANGcs{
\Question{
Hyperbola má rovnice asymptot \( y=\frac13(x+2) \) a \(y=-\frac13(x+2)\). Střed této hyperboly má souřadnice:}
{\( [-2;0] \)}{\( [2;0] \)}{\( \left[-\frac13;0\right] \)}{\( [0;2] \)}{\( \left[\frac13;2\right] \)}{}}
\LANGsk{
\Question{
Hyperbola má rovnice asymptot \( y=\frac13(x+2) \) a \(y=-\frac13(x+2)\). Stred tejto hyperboly má súradnice:}
{\( [-2;0] \)}{\( [2;0] \)}{\( \left[-\frac13;0\right] \)}{\( [0;2] \)}{\( \left[\frac13;2\right] \)}{}}
\LANGes{
\Question{
Una hipérbola tiene dos asíntotas cuyas ecuaciones son \( y=\frac13(x+2) \) y \(y=-\frac13(x+2)\). Define las coordenadas del centro de esta hipérbola.}
{\( [-2;0] \)}{\( [2;0] \)}{\( \left[-\frac13;0\right] \)}{\( [0;2] \)}{\( \left[\frac13;2\right] \)}{}}
\End

\Data\ID{1003020406}
\Updated{2022-08-25 10:08:09}
\Level{B}
\SubArea{Conic Sections}
\LANGen{
\Question{
The equation of a hyperbola is given by \( \frac{x^2}4-\frac{y^2}3=1 \). Which of the following are the asymptote equations of the hyperbola?}
{\( y=\pm\frac{\sqrt3}2x \)}{\( x=\pm\frac{\sqrt3}2y \)}{\( y=\pm\frac34 x \)}{\( y=\pm\frac43 x \)}{\( y=\pm\frac{2\sqrt3}3 x \)}{}}
\LANGpl{
\Question{
Dane jest równanie hiperboli \( \frac{x^2}4-\frac{y^2}3=1 \). Wskaż równanie asymptoty hiperboli.}
{\( y=\pm\frac{\sqrt3}2x \)}{\( x=\pm\frac{\sqrt3}2y \)}{\( y=\pm\frac34 x \)}{\( y=\pm\frac43 x \)}{\( y=\pm\frac{2\sqrt3}3 x \)}{}}
\LANGcs{
\Question{
Hyperbola je dána rovnicí \( \frac{x^2}4-\frac{y^2}3=1 \). Rovnice asymptot této hyperboly jsou:}
{\( y=\pm\frac{\sqrt3}2x \)}{\( x=\pm\frac{\sqrt3}2y \)}{\( y=\pm\frac34 x \)}{\( y=\pm\frac43 x \)}{\( y=\pm\frac{2\sqrt3}3 x \)}{}}
\LANGsk{
\Question{
Hyperbola je daná rovnicou \( \frac{x^2}4-\frac{y^2}3=1 \). Rovnice asymptot tejto hyperboly sú:}
{\( y=\pm\frac{\sqrt3}2x \)}{\( x=\pm\frac{\sqrt3}2y \)}{\( y=\pm\frac34 x \)}{\( y=\pm\frac43 x \)}{\( y=\pm\frac{2\sqrt3}3 x \)}{}}
\LANGes{
\Question{
Dada la ecuación de la hipérbola \( \frac{x^2}4-\frac{y^2}3=1 \). Identifica las ecuaciones de sus asíntotas.}
{\( y=\pm\frac{\sqrt3}2x \)}{\( x=\pm\frac{\sqrt3}2y \)}{\( y=\pm\frac34 x \)}{\( y=\pm\frac43 x \)}{\( y=\pm\frac{2\sqrt3}3 x \)}{}}
\End

\Data\ID{1003020407}
\Updated{2022-08-25 10:08:09}
\Level{B}
\SubArea{Conic Sections}
\LANGen{
\Question{
The equation of a hyperbola is given by \( \frac{(x-5)^2}{12}-\frac{(y+3)^2}8=1 \).  The centre of the hyperbola has the coordinates:}
{\( [5;-3] \)}{\( [-5;3] \)}{\( [-5;-3] \)}{\( [12;8] \)}{\( [8;12] \)}{}}
\LANGpl{
\Question{
Dane jest równanie hiperboli \( \frac{(x-5)^2}{12}-\frac{(y+3)^2}8=1 \).  Wskaż współrzędne środka hiperboli.}
{\( [5;-3] \)}{\( [-5;3] \)}{\( [-5;-3] \)}{\( [12;8] \)}{\( [8;12] \)}{}}
\LANGcs{
\Question{
Hyperbola je dána rovnicí \( \frac{(x-5)^2}{12}-\frac{(y+3)^2}8=1 \).  Střed této hyperboly má souřadnice:}
{\( [5;-3] \)}{\( [-5;3] \)}{\( [-5;-3] \)}{\( [12;8] \)}{\( [8;12] \)}{}}
\LANGsk{
\Question{
Hyperbola je daná rovnicou \( \frac{(x-5)^2}{12}-\frac{(y+3)^2}8=1 \).  Stred tejto hyperboly má súradnice:}
{\( [5;-3] \)}{\( [-5;3] \)}{\( [-5;-3] \)}{\( [12;8] \)}{\( [8;12] \)}{}}
\LANGes{
\Question{
Dada la ecuación de la hipérbola \( \frac{(x-5)^2}{12}-\frac{(y+3)^2}8=1 \).  Identifica las coordenadas del centro de la hipérbola:}
{\( [5;-3] \)}{\( [-5;3] \)}{\( [-5;-3] \)}{\( [12;8] \)}{\( [8;12] \)}{}}
\End

\Data\ID{1003020408}
\Updated{2022-08-25 10:08:09}
\Level{B}
\SubArea{Conic Sections}
\LANGen{
\Question{
The equation of a hyperbola is given by \( \frac{x^2}9-\frac{y^2}{16}=1 \). What is the distance between its foci?}
{\( 10 \)}{\( 5 \)}{\( 16 \)}{\( 9 \)}{\( 4 \)}{}}
\LANGpl{
\Question{
Dane jest równanie hiperboli \( \frac{x^2}9-\frac{y^2}{16}=1 \). Wskaż odległość pomiędzy ogniskami danej hiperboli.}
{\( 10 \)}{\( 5 \)}{\( 16 \)}{\( 9 \)}{\( 4 \)}{}}
\LANGcs{
\Question{
Hyperbola je dána rovnicí \( \frac{x^2}9-\frac{y^2}{16}=1 \). Vzdálenost ohnisek této hyperboly je:}
{\( 10 \)}{\( 5 \)}{\( 16 \)}{\( 9 \)}{\( 4 \)}{}}
\LANGsk{
\Question{
Hyperbola je daná rovnicou \( \frac{x^2}9-\frac{y^2}{16}=1 \). Vzdialenosť ohnísk tejto hyperboly je:}
{\( 10 \)}{\( 5 \)}{\( 16 \)}{\( 9 \)}{\( 4 \)}{}}
\LANGes{
\Question{
Dada la ecuación de la hipérbola \( \frac{x^2}9-\frac{y^2}{16}=1 \). Halla la distancia entre sus focos.}
{\( 10 \)}{\( 5 \)}{\( 16 \)}{\( 9 \)}{\( 4 \)}{}}
\End

\Data\ID{1003020409}
\Updated{2022-08-25 10:08:09}
\Level{B}
\SubArea{Conic Sections}
\LANGen{
\Question{
The equation of a parabola is given by \( y^2 = 4x \). What is the parameter value of this parabola?}
{\( 2 \)}{\( 4 \)}{\( \frac12 \)}{\( \frac14 \)}{\( 1 \)}{}}
\LANGpl{
\Question{
Parabola jest opisana równaniem \( y^2 = 4x \). Wyznacz wartość parametru paraboli.}
{\( 2 \)}{\( 4 \)}{\( \frac12 \)}{\( \frac14 \)}{\( 1 \)}{}}
\LANGcs{
\Question{
Parabola je dána rovnicí \( y^2 = 4x \). Parametr této paraboly je:}
{\( 2 \)}{\( 4 \)}{\( \frac12 \)}{\( \frac14 \)}{\( 1 \)}{}}
\LANGsk{
\Question{
Parabola je daná rovnicou \( y^2 = 4x \). Parameter tejto paraboly je:}
{\( 2 \)}{\( 4 \)}{\( \frac12 \)}{\( \frac14 \)}{\( 1 \)}{}}
\LANGes{
\Question{
Dada la ecuación de la parábola \( y^2 = 4x \). Identifica el parámetro de la parábola.}
{\( 2 \)}{\( 4 \)}{\( \frac12 \)}{\( \frac14 \)}{\( 1 \)}{}}
\End

\Data\ID{1003020410}
\Updated{2022-08-25 10:08:09}
\Level{B}
\SubArea{Conic Sections}
\LANGen{
\Question{
The equation of a parabola is given by \( x^2+4y-6x+3=0 \). What is the parameter value of this parabola?}
{\( 2 \)}{\( 4 \)}{\( -2 \)}{\( -4 \)}{\( \frac32 \)}{}}
\LANGpl{
\Question{
Parabola jest opisana równaniem \( x^2+4y-6x+3=0 \). Wyznacz wartość parametru danej paraboli.}
{\( 2 \)}{\( 4 \)}{\( -2 \)}{\( -4 \)}{\( \frac32 \)}{}}
\LANGcs{
\Question{
Parabola je dána rovnicí \( x^2+4y-6x+3=0 \). Parametr této paraboly je:}
{\( 2 \)}{\( 4 \)}{\( -2 \)}{\( -4 \)}{\( \frac32 \)}{}}
\LANGsk{
\Question{
Parabola je daná rovnicou \( x^2+4y-6x+3=0 \). Parameter tejto paraboly je:}
{\( 2 \)}{\( 4 \)}{\( -2 \)}{\( -4 \)}{\( \frac32 \)}{}}
\LANGes{
\Question{
Dada la ecuación de la parábola \( x^2+4y-6x+3=0 \). Identifica el parámetro de la parábola.}
{\( 2 \)}{\( 4 \)}{\( -2 \)}{\( -4 \)}{\( \frac32 \)}{}}
\End

\Data\ID{1003020411}
\Updated{2023-06-05 09:06:34}
\Level{B}
\SubArea{Conic Sections}
\LANGen{
\Question{
A parabola has its focus \( F=[5;0] \) and the  directrix with the equation \( x=0 \). What are the coordinates of its vertex?}
{\( \left[\frac52;0\right] \)}{\( [0;0] \)}{\( \left[0;\frac52\right] \)}{\( [-5;0] \)}{\( \left[-\frac52;0\right] \)}{}}
\LANGpl{
\Question{
Parabola ma ognisko \( F=[5;0] \) oraz kierownicę o równaniu \( x=0 \). Wyznacz współrzędne wierzchołka paraboli.}
{\( \left[\frac52;0\right] \)}{\( [0;0] \)}{\( \left[0;\frac52\right] \)}{\( [-5;0] \)}{\( \left[-\frac52;0\right] \)}{}}
\LANGcs{
\Question{
Parabola má ohnisko \( F=[5;0] \) a řídící přímku o rovnici \( x=0 \). Vrchol této paraboly má souřadnice:}
{\( \left[\frac52;0\right] \)}{\( [0;0] \)}{\( \left[0;\frac52\right] \)}{\( [-5;0] \)}{\( \left[-\frac52;0\right] \)}{}}
\LANGsk{
\Question{
Parabola má ohnisko \( F=[5;0] \) a riadiacu priamku s rovnicou \( x=0 \). Vrchol tejto paraboly má súradnice:}
{\( \left[\frac52;0\right] \)}{\( [0;0] \)}{\( \left[0;\frac52\right] \)}{\( [-5;0] \)}{\( \left[-\frac52;0\right] \)}{}}
\LANGes{
\Question{
Una parábola tiene su foco en \( F=(5;0) \) y su directriz es \( x=0 \). Halla las coordenadas de su vértice.}
{\( \left[\frac52;0\right] \)}{\( [0;0] \)}{\( \left[0;\frac52\right] \)}{\( [-5;0] \)}{\( \left[-\frac52;0\right] \)}{}}
\End

\Data\ID{1003020412}
\Updated{2022-08-25 10:08:09}
\Level{B}
\SubArea{Conic Sections}
\LANGen{
\Question{
A parabola passes through the points \( K=[0;0] \) and \( L=[6;-2] \) and it is symmetric with respect to \( y \)-axis. What are the coordinates of its vertex?}
{\( [0;0] \)}{\( [6;-2] \)}{\( [3;0] \)}{\( [-6;-2] \)}{\( [0;-1] \)}{}}
\LANGpl{
\Question{
Parabola przechodzi przez punkty \( K=[0;0] \) i \( L=[6;-2] \) i jest symetryczna względem osi  \( y \). Wyznacz współrzędne jej wierzchołka.}
{\( [0;0] \)}{\( [6;-2] \)}{\( [3;0] \)}{\( [-6;-2] \)}{\( [0;-1] \)}{}}
\LANGcs{
\Question{
Parabola, která prochází body \( K=[0;0] \) a \( L=[6;-2] \) a je souměrná podle osy \( y \), má vrchol o souřadnicích:}
{\( [0;0] \)}{\( [6;-2] \)}{\( [3;0] \)}{\( [-6;-2] \)}{\( [0;-1] \)}{}}
\LANGsk{
\Question{
Parabola, ktorá prechádza bodmi \( K=[0;0] \) a \( L=[6;-2] \) a je súmerná podľa osi \( y \), má vrchol so súradnicami:}
{\( [0;0] \)}{\( [6;-2] \)}{\( [3;0] \)}{\( [-6;-2] \)}{\( [0;-1] \)}{}}
\LANGes{
\Question{
Una parábola pasa por los puntos \( K=[0;0] \) y \( L=[6;-2] \) y es simétrica respecto al eje \( y \). Halla las coordenadas de su vértice.}
{\( [0;0] \)}{\( [6;-2] \)}{\( [3;0] \)}{\( [-6;-2] \)}{\( [0;-1] \)}{}}
\End

\Data\ID{1003020413}
\Updated{2022-08-25 10:08:09}
\Level{B}
\SubArea{Conic Sections}
\LANGen{
\Question{
The equation of a parabola is given by \( x^2+6x+y+10=0 \). What are  the coordinates of its vertex?}
{\( [-3;-1] \)}{\( [3;1] \)}{\( [3;-1] \)}{\( [-1;-3] \)}{\( [-1;3] \)}{}}
\LANGpl{
\Question{
Parabola jest opisana równaniem \( x^2+6x+y+10=0 \). Wyznacz współrzędne jej wierzchołka.}
{\( [-3;-1] \)}{\( [3;1] \)}{\( [3;-1] \)}{\( [-1;-3] \)}{\( [-1;3] \)}{}}
\LANGcs{
\Question{
Parabola je dána rovnicí \( x^2+6x+y+10=0 \). Vrchol této paraboly má souřadnice:}
{\( [-3;-1] \)}{\( [3;1] \)}{\( [3;-1] \)}{\( [-1;-3] \)}{\( [-1;3] \)}{}}
\LANGsk{
\Question{
Parabola je daná rovnicou \( x^2+6x+y+10=0 \). Vrchol tejto paraboly má súradnice:}
{\( [-3;-1] \)}{\( [3;1] \)}{\( [3;-1] \)}{\( [-1;-3] \)}{\( [-1;3] \)}{}}
\LANGes{
\Question{
Dada la ecuación de la parábola \( x^2+6x+y+10=0 \). Halla las coordenadas de su vértice.}
{\( [-3;-1] \)}{\( [3;1] \)}{\( [3;-1] \)}{\( [-1;-3] \)}{\( [-1;3] \)}{}}
\End

\Data\ID{1003024001}
\Updated{2022-08-25 10:08:36}
\Level{B}
\SubArea{Conic Sections}
\LANGen{
\Question{
The equation of a hyperbola is given by \( 4x^2-y^2-12=0 \). Which of the following points lies on such hyperbola?}
{\( [-2;2] \)}{\( [2;1] \)}{\( [0;12] \)}{\( [3;0] \)}{\( [\sqrt3;12] \)}{}}
\LANGpl{
\Question{
Hiperbola jest opisana równaniem \( 4x^2-y^2-12=0 \). Wskaż punkt, który leży na hiperboli.}
{\( [-2;2] \)}{\( [2;1] \)}{\( [0;12] \)}{\( [3;0] \)}{\( [\sqrt3;12] \)}{}}
\LANGcs{
\Question{
Hyperbola je dána rovnicí \( 4x^2-y^2-12=0 \). Určete, který z následujících bodů leží na této hyperbole:}
{\( [-2;2] \)}{\( [2;1] \)}{\( [0;12] \)}{\( [3;0] \)}{\( [\sqrt3;12] \)}{}}
\LANGsk{
\Question{
Hyperbola je daná rovnicou \( 4x^2-y^2-12=0 \). Určte, ktorý z následujúcich bodov leží na tejto hyperbole:}
{\( [-2;2] \)}{\( [2;1] \)}{\( [0;12] \)}{\( [3;0] \)}{\( [\sqrt3;12] \)}{}}
\LANGes{
\Question{
Dada la ecuación de la hipérbola \( 4x^2-y^2-12=0 \). ¿Cuáles de los siguientes puntos se encuentran en la hipérbola?}
{\( [-2;2] \)}{\( [2;1] \)}{\( [0;12] \)}{\( [3;0] \)}{\( [\sqrt3;12] \)}{}}
\End

\Data\ID{1003024003}
\Updated{2023-06-05 08:06:51}
\Level{B}
\SubArea{Conic Sections}
\LANGen{
\Question{
Let \( A=[0;0] \), \( B=[1;7] \), \( C=[-1;-5] \) be points that  all lie on a parabola. Find the equation of the parabola.}
{\( y=x^2+6x \)}{\( y=x^2 \)}{\( y=7x^2 \)}{\( y=-x^2+4x \)}{\( x=y^2 \)}{}}
\LANGpl{
\Question{
Punkty \( A=[0;0] \), \( B=[1;7] \), \( C=[-1;-5] \) leżą na  paraboli. Wskaż równanie paraboli.}
{\( y=x^2+6x \)}{\( y=x^2 \)}{\( y=7x^2 \)}{\( y=-x^2+4x \)}{\( x=y^2 \)}{}}
\LANGcs{
\Question{
Jsou dány body \( A=[0;0] \), \( B=[1;7] \), \( C=[-1;-5] \). Všechny tyto tři body leží na parabole o rovnici:}
{\( y=x^2+6x \)}{\( y=x^2 \)}{\( y=7x^2 \)}{\( y=-x^2+4x \)}{\( x=y^2 \)}{}}
\LANGsk{
\Question{
Sú dané body \( A=[0;0] \), \( B=[1;7] \), \( C=[-1;-5] \). Všetky tieto tri body ležia na parabole s rovnicou:}
{\( y=x^2+6x \)}{\( y=x^2 \)}{\( y=7x^2 \)}{\( y=-x^2+4x \)}{\( x=y^2 \)}{}}
\LANGes{
\Question{
Dados los puntos \( A=(0;0) \), \( B=(1;7) \), \( C=(-1;-5) \) que son puntos de la parábola. Halla la ecuación de dicha parábola.}
{\( y=x^2+6x \)}{\( y=x^2 \)}{\( y=7x^2 \)}{\( y=-x^2+4x \)}{\( x=y^2 \)}{}}
\End

\Data\ID{1003024004}
\Updated{2022-08-25 10:08:40}
\Level{B}
\SubArea{Conic Sections}
\LANGen{
\Question{
Decide whether the equation \( 4x^2+9y^2-8x+18y+13=0 \) (in \( x \)-\( y \) plane) describes:}
{a point}{an ellipse}{a parabola}{a hyperbola}{an empty set}{}}
\LANGpl{
\Question{
Równanie  \( 4x^2+9y^2-8x+18y+13=0 \) (na płaszczyźnie \( x \)-\( y \) ) opisuje:}
{punkt}{elipsę}{parabolę}{hiperbolę}{zbiór pusty}{}}
\LANGcs{
\Question{
Rovnice \( 4x^2+9y^2-8x+18y+13=0 \) (v rovině \( x \)\( y \)) popisuje:}
{bod}{elipsu}{parabolu}{hyperbolu}{prázdnou množinu}{}}
\LANGsk{
\Question{
Rovnica \( 4x^2+9y^2-8x+18y+13=0 \) (v rovine \( x \)\( y \)) popisuje:}
{bod}{elipsu}{parabolu}{hyperbolu}{prázdnu množinu}{}}
\LANGes{
\Question{
Decide qué lugar geométrico describe la ecuación \( 4x^2+9y^2-8x+18y+13=0 \) (en el plano \( x \)-\( y \)):}
{el punto}{una elipse}{una parábola}{una hipérbola}{el conjunto vacío}{}}
\End

\Data\ID{1003024005}
\Updated{2022-08-25 10:08:40}
\Level{B}
\SubArea{Conic Sections}
\LANGen{
\Question{
Decide whether the equation \( 4x^2+4y-8x=1\) (in \( x \)-\( y \) plane) describes:}
{a parabola}{an ellipse}{a circle}{a hyperbola}{point \( \left[0;\frac14\right] \)}{}}
\LANGpl{
\Question{
Równanie \( 4x^2+4y-8x=1\) (na płaszczyźnie \( x \)-\( y \)) opisuje:}
{parabolę}{elipsę}{okrąg}{hiperbolę}{punkt \( \left[0;\frac14\right] \)}{}}
\LANGcs{
\Question{
Rovnice \( 4x^2+4y-8x=1\) (v rovině \( x \)\( y \)) popisuje:}
{parabolu}{elipsu}{kružnici}{hyperbolu}{bod \( \left[0;\frac14\right] \)}{}}
\LANGsk{
\Question{
Rovnica \( 4x^2+4y-8x=1\) (v rovine \( x \)\( y \)) popisuje:}
{parabolu}{elipsu}{kružnicu}{hyperbolu}{bod \( \left[0;\frac14\right] \)}{}}
\LANGes{
\Question{
Decide qué lugar geométrico describe la ecuación \( 4x^2+4y-8x=1\) (en el plano \( x \)-\( y \)):}
{una parábola}{una elipse}{una circunferencia}{una hipérbola}{el punto \( \left[0;\frac14\right] \)}{}}
\End

\Data\ID{1003024006}
\Updated{2022-08-25 10:08:40}
\Level{B}
\SubArea{Conic Sections}
\LANGen{
\Question{
Decide whether the equation \( (x+3)^2+(y+4)^2=5 \) (in \( x \)-\( y \) plane) describes:}
{a circle with center \( [-3;-4] \)}{a circle with center \( [3;4] \)}{point \( [-3;-4] \)}{a hyperbola}{an ellipse with center  \( [3;4] \)}{}}
\LANGpl{
\Question{
Równanie \( (x+3)^2+(y+4)^2=5 \) (na płaszczyźnie \( x \)-\( y \)) opisuje:}
{okrąg o środku \( [-3;-4] \)}{okrąg o środku \( [3;4] \)}{punkt \( [-3;-4] \)}{hiperbolę}{elipsę o środku  \( [3;4] \)}{}}
\LANGcs{
\Question{
Rovnice \( (x+3)^2+(y+4)^2=5 \) (v rovině \( x \)\( y \)) popisuje:}
{kružnici se středem \( [-3;-4] \)}{kružnici se středem \( [3;4] \)}{bod \( [-3;-4] \)}{hyperbolu}{elipsu se středem  \( [3;4] \)}{}}
\LANGsk{
\Question{
Rovnica \( (x+3)^2+(y+4)^2=5 \) (v rovine \( x \)\( y \)) popisuje:}
{kružnicu so stredom \( [-3;-4] \)}{kružnicu so stredom \( [3;4] \)}{bod \( [-3;-4] \)}{hyperbolu}{elipsu so stredom  \( [3;4] \)}{}}
\LANGes{
\Question{
Decide qué lugar geométrico describe la ecuación \( (x+3)^2+(y+4)^2=5 \) (en el plano \( x \)-\( y \)):}
{una circunferencia con el centro \( [-3;-4] \)}{una circunferencia con el centro \( [3;4] \)}{el punto \( [-3;-4] \)}{una hipérbola}{una elipse con el centro \( [3;4] \)}{}}
\End

\Data\ID{1003024007}
\Updated{2022-08-25 10:08:40}
\Level{B}
\SubArea{Conic Sections}
\LANGen{
\Question{
Decide whether the equation \( 2x^2-4y^2-6x-12y=0 \) (in \( x \)-\( y \) plane) describes:}
{a hyperbola}{a circle}{a parabola}{an ellipse}{point \( \left[\frac32;\frac{-3}2\right] \)}{}}
\LANGpl{
\Question{
Równanie \( 2x^2-4y^2-6x-12y=0 \) (na płaszczyźnie \( x \)-\( y \)) opisuje:}
{hiperbolę}{okrąg}{parabolę}{elipsę}{punkt \( \left[\frac32;\frac{-3}2\right] \)}{}}
\LANGcs{
\Question{
Rovnice \( 2x^2-4y^2-6x-12y=0 \) (v rovině \( x \)\( y \)) popisuje:}
{hyperbolu}{kružnici}{parabolu}{elipsu}{bod \( \left[\frac32;\frac{-3}2\right] \)}{}}
\LANGsk{
\Question{
Rovnica \( 2x^2-4y^2-6x-12y=0 \) (v rovine \( x \)\( y \)) popisuje:}
{hyperbolu}{kružnicu}{parabolu}{elipsu}{bod \( \left[\frac32;\frac{-3}2\right] \)}{}}
\LANGes{
\Question{
Decide qué lugar geométrico describe la ecuación \( 2x^2-4y^2-6x-12y=0 \) (en el plano \( x \)-\( y \)):}
{una hipérbola}{una circunferencia}{una parábola}{una elipse}{el punto \( \left[\frac32;\frac{-3}2\right] \)}{}}
\End

\Data\ID{1003024008}
\Updated{2022-08-25 10:08:40}
\Level{B}
\SubArea{Conic Sections}
\LANGen{
\Question{
Decide which of the given equations describes a hyperbola.}
{\( 4x^2-9y^2+18y-45=0 \)}{\( 4x^2+9y^2-8x-36y+4=0 \)}{\( x^2+y^2-2x+4y-4=0 \)}{\( x^2+y^2-12x+40=0 \)}{\( x^2-2x-4y+1=0 \)}{}}
\LANGpl{
\Question{
Wskaż, które z poniższych równań opisuje hiperbolę.}
{\( 4x^2-9y^2+18y-45=0 \)}{\( 4x^2+9y^2-8x-36y+4=0 \)}{\( x^2+y^2-2x+4y-4=0 \)}{\( x^2+y^2-12x+40=0 \)}{\( x^2-2x-4y+1=0 \)}{}}
\LANGcs{
\Question{
Která z daných rovnic představuje rovnici hyperboly?}
{\( 4x^2-9y^2+18y-45=0 \)}{\( 4x^2+9y^2-8x-36y+4=0 \)}{\( x^2+y^2-2x+4y-4=0 \)}{\( x^2+y^2-12x+40=0 \)}{\( x^2-2x-4y+1=0 \)}{}}
\LANGsk{
\Question{
Ktorá z daných rovníc predstavuje rovnicu hyperboly?}
{\( 4x^2-9y^2+18y-45=0 \)}{\( 4x^2+9y^2-8x-36y+4=0 \)}{\( x^2+y^2-2x+4y-4=0 \)}{\( x^2+y^2-12x+40=0 \)}{\( x^2-2x-4y+1=0 \)}{}}
\LANGes{
\Question{
Decide cuál de las ecuaciones describe una hipérbola.}
{\( 4x^2-9y^2+18y-45=0 \)}{\( 4x^2+9y^2-8x-36y+4=0 \)}{\( x^2+y^2-2x+4y-4=0 \)}{\( x^2+y^2-12x+40=0 \)}{\( x^2-2x-4y+1=0 \)}{}}
\End

\Data\ID{1003024101}
\Updated{2022-08-25 10:08:40}
\Level{B}
\SubArea{Conic Sections}
\LANGen{
\Question{
The equation of a hyperbola with the center \( S=[-1;3] \),  the focus \( F=[4;3] \), and the vertex \( A=[2;3] \) is given by:}
{\( \frac{(x+1)^2}{9}-\frac{(y-3)^2}{16} =1 \)}{\( \frac{(x-1)^2}{9}-\frac{(y+3)^2}{16} =1 \)}{\( \frac{(x+1)^2}{9}+\frac{(y-3)^2}{16} =1 \)}{\( \frac{(x-1)^2}{16}-\frac{(y+3)^2}{9} =1 \)}{\( \frac{(y-3)^2}{16}-\frac{(x+1)^2}{9} =1 \)}{}}
\LANGpl{
\Question{
Wskaż równanie opisujące hiperbolę o środku \( S=[-1;3] \),  ognisku \( F=[4;3] \) oraz wierzchołku \( A=[2;3] \).}
{\( \frac{(x+1)^2}{9}-\frac{(y-3)^2}{16} =1 \)}{\( \frac{(x-1)^2}{9}-\frac{(y+3)^2}{16} =1 \)}{\( \frac{(x+1)^2}{9}+\frac{(y-3)^2}{16} =1 \)}{\( \frac{(x-1)^2}{16}-\frac{(y+3)^2}{9} =1 \)}{\( \frac{(y-3)^2}{16}-\frac{(x+1)^2}{9} =1 \)}{}}
\LANGcs{
\Question{
Rovnice hyperboly, která má střed \( S=[-1;3] \),  ohnisko \( F=[4;3] \) a vrchol \( A=[2;3] \), je:}
{\( \frac{(x+1)^2}{9}-\frac{(y-3)^2}{16} =1 \)}{\( \frac{(x-1)^2}{9}-\frac{(y+3)^2}{16} =1 \)}{\( \frac{(x+1)^2}{9}+\frac{(y-3)^2}{16} =1 \)}{\( \frac{(x-1)^2}{16}-\frac{(y+3)^2}{9} =1 \)}{\( \frac{(y-3)^2}{16}-\frac{(x+1)^2}{9} =1 \)}{}}
\LANGsk{
\Question{
Rovnica hyperboly, ktorá má stred \( S=[-1;3] \),  ohnisko \( F=[4;3] \) a vrchol \( A=[2;3] \), je:}
{\( \frac{(x+1)^2}{9}-\frac{(y-3)^2}{16} =1 \)}{\( \frac{(x-1)^2}{9}-\frac{(y+3)^2}{16} =1 \)}{\( \frac{(x+1)^2}{9}+\frac{(y-3)^2}{16} =1 \)}{\( \frac{(x-1)^2}{16}-\frac{(y+3)^2}{9} =1 \)}{\( \frac{(y-3)^2}{16}-\frac{(x+1)^2}{9} =1 \)}{}}
\LANGes{
\Question{
Elige la ecuación de la hipérbola con centro \( S=(-1;3) \), foco \( F=(4;3) \) y  vértice \( A=(2;3) \):}
{\( \frac{(x+1)^2}{9}-\frac{(y-3)^2}{16} =1 \)}{\( \frac{(x-1)^2}{9}-\frac{(y+3)^2}{16} =1 \)}{\( \frac{(x+1)^2}{9}+\frac{(y-3)^2}{16} =1 \)}{\( \frac{(x-1)^2}{16}-\frac{(y+3)^2}{9} =1 \)}{\( \frac{(y-3)^2}{16}-\frac{(x+1)^2}{9} =1 \)}{}}
\End

\Data\ID{1003024102}
\Updated{2022-08-25 10:08:40}
\Level{B}
\SubArea{Conic Sections}
\LANGen{
\Question{
The equation of a parabola is given by \( 3y^2+x-12y+14=0 \). Find the equation of its directrix.}
{\( x=-\frac{23}{12} \)}{\( x=\frac{23}{12} \)}{\( y=-\frac{23}{12} \)}{\( y=\frac{23}{12} \)}{\( x=-\frac{11}{6} \)}{}}
\LANGpl{
\Question{
Parabola jest opisana równaniem \( 3y^2+x-12y+14=0 \). Wskaż równanie kierownicy danej paraboli.}
{\( x=-\frac{23}{12} \)}{\( x=\frac{23}{12} \)}{\( y=-\frac{23}{12} \)}{\( y=\frac{23}{12} \)}{\( x=-\frac{11}{6} \)}{}}
\LANGcs{
\Question{
Parabola je dána rovnicí \( 3y^2+x-12y+14=0 \). Rovnice řídící přímky této paraboly je:}
{\( x=-\frac{23}{12} \)}{\( x=\frac{23}{12} \)}{\( y=-\frac{23}{12} \)}{\( y=\frac{23}{12} \)}{\( x=-\frac{11}{6} \)}{}}
\LANGsk{
\Question{
Parabola je daná rovnicou \( 3y^2+x-12y+14=0 \). Rovnica riadiacej priamky tejto paraboly je:}
{\( x=-\frac{23}{12} \)}{\( x=\frac{23}{12} \)}{\( y=-\frac{23}{12} \)}{\( y=\frac{23}{12} \)}{\( x=-\frac{11}{6} \)}{}}
\LANGes{
\Question{
Dada la ecuación de la parábola \( 3y^2+x-12y+14=0 \). Halla la ecuación de su directriz.}
{\( x=-\frac{23}{12} \)}{\( x=\frac{23}{12} \)}{\( y=-\frac{23}{12} \)}{\( y=\frac{23}{12} \)}{\( x=-\frac{11}{6} \)}{}}
\End

\Data\ID{1103040104}
\Updated{2020-11-11 09:11:19}
\Level{B}
\SubArea{Conic Sections}
\IMG{1103040104.pdf}
\LANGen{
\Question{
The picture shows a hyperbola drawn in the rectangular coordinate system. What is the eccentricity of this hyperbola?}
{The distance of points \( S \) and \( F \)}{The distance of points \( S \) and \( A \)}{The distance of points \( A \) and \( B \)}{The distance of points \( E \) and \( F \)}{}{}}
\LANGpl{
\Question{
Rysunek przedstawia hiperbolę  na kartezjańskim układzie współrzędnych. Wskaż mimośród hiperboli.}
{Odległość między punktami \( S \) i \( F \)}{Odległość między punktami \( S \) i \( A \)}{Odległość między punktami  \( A \) i \( B \)}{Odległość między punktami \( E \) i \( F \)}{}{}}
\LANGcs{
\Question{
Hyperbola je dána obrázkem v kartézské soustavě souřadnic. Excentricita této hyperboly je:}
{vzdálenost bodů \( S \) a \( F \)}{vzdálenost bodů \( S \) a \( A \)}{vzdálenost bodů \( A \) a \( B \)}{vzdálenost bodů \( E \) a \( F \)}{}{}}
\LANGsk{
\Question{
Hyperbola je daná obrázkom v kartézskej sústave súradníc. Excentricita tejto hyperboly je:}
{vzdialenosť bodov \( S \) a \( F \)}{vzdialenosť bodov \( S \) a \( A \)}{vzdialenosť bodov \( A \) a \( B \)}{vzdialenosť bodov \( E \) a \( F \)}{}{}}
\LANGes{
\Question{
La imagen representa una hipérbola en el sistema de coordenadas cartesianas. Define la excentricidad de esta hipérbola.}
{La distancia entre los puntos \( S \) y \( F \)}{La distancia entre los puntos \( S \) y \( A \)}{La distancia entre los puntos \( A \) y \( B \)}{La distancia entre los puntos \( E \) y \( F \)}{}{}}
\End

\Data\ID{1103040105}
\Updated{2020-11-07 07:11:00}
\Level{B}
\SubArea{Conic Sections}
\IMG{1103040105.pdf}
\LANGen{
\Question{
A parabola is given by the picture. What is the parameter of this parabola?}
{The distance of the point \( F \) from the line \( d \)}{The distance of points \( V \) and \( F \)}{A half of the segment \( DV \)}{Two times the distance of the point \( F \) from the line \( d \)}{}{}}
\LANGpl{
\Question{
Rysunek przedstawia parabolę. Wskaż parametr tej paraboli.}
{Odległość  punktu \( F \) od prostej \( d \)}{Odległość punktów \( V \) i \( F \)}{Połowa odcinka  \( DV \)}{Dwukrotna odległość punktu \( F \) od prostej \( d \)}{}{}}
\LANGcs{
\Question{
Parabola je dána obrázkem. Parametr této paraboly je:}
{vzdálenost bodu \( F \) od přímky \( d \)}{vzdálenost bodů \( V \) a \( F \)}{polovina délky úsečky  \( DV \)}{dvojnásobek vzdálenosti bodu  \( F \) od přímky \( d \)}{}{}}
\LANGsk{
\Question{
Parabola je daná obrázkom. Parameter tejto paraboly je:}
{vzdialenosť bodu \( F \) od priamky \( d \)}{vzdialenosť bodov \( V \) a \( F \)}{polovica dľžky úsečky  \( DV \)}{dvojnásobok vzdialenosti bodu  \( F \) od priamky \( d \)}{}{}}
\LANGes{
\Question{
En la imagen se representa una parábola. ¿Cuál es el parámetro de esta parábola?}
{La diferencia del punto \( F \) de la recta \( d \)}{La distancia entre los puntos \( V \) y \( F \)}{La mitad del segmento \( DV \)}{La distancia doble entre el punto \( F \) de la recta \( d \)}{}{}}
\End

\Data\ID{1103040108}
\Updated{2021-09-25 08:09:19}
\Level{B}
\SubArea{Conic Sections}
\IMG{1103040108.pdf}
\LANGen{
\Question{
A parabola in the rectangular coordinate system is shown in the picture. What is the vertex form equation of this parabola?}
{\( x^2 = 4(y-1) \)}{\( x^2 = 4(y+1) \)}{\( y^2 = 4(x-1) \)}{\( y^2 = 4(x+1) \)}{}{}}
\LANGpl{
\Question{
Rysunek przedstawia parabolę na kartezjańskim układzie współrzędnych. Wskaż równanie wierzchołka danej paraboli.}
{\( x^2 = 4(y-1) \)}{\( x^2 = 4(y+1) \)}{\( y^2 = 4(x-1) \)}{\( y^2 = 4(x+1) \)}{}{}}
\LANGcs{
\Question{
Parabola je dána obrázkem v kartézské soustavě souřadnic. Vrcholový tvar rovnice této paraboly je:}
{\( x^2 = 4(y-1) \)}{\( x^2 = 4(y+1) \)}{\( y^2 = 4(x-1) \)}{\( y^2 = 4(x+1) \)}{}{}}
\LANGsk{
\Question{
Parabola je daná obrázkom v kartézskej sústave súradníc. Vrcholový tvar rovnice tejto paraboly je:}
{\( x^2 = 4(y-1) \)}{\( x^2 = 4(y+1) \)}{\( y^2 = 4(x-1) \)}{\( y^2 = 4(x+1) \)}{}{}}
\LANGes{
\Question{
En la imagen se representa una parábola en el sistema de coordenadas cartesianas. Identifica su ecuación.}
{\( x^2 = 4(y-1) \)}{\( x^2 = 4(y+1) \)}{\( y^2 = 4(x-1) \)}{\( y^2 = 4(x+1) \)}{}{}}
\End

\Data\ID{1103040109}
\Updated{2020-11-08 06:11:25}
\Level{B}
\SubArea{Conic Sections}
\IMG{1103040109.pdf}
\LANGen{
\Question{
The picture shows a hyperbola drawn in the rectangular coordinate system. What is the minor axis of this hyperbola?}
{The \( y \)-axis}{The segment \( EF \)}{The \( x \)-axis}{The segment \( AB \)}{}{}}
\LANGpl{
\Question{
Rysunek przedstawia hiperbolę na kartezjańskim układzie współrzędnych. Wskaż oś małą hiperboli.}
{Oś \( y \)}{Odcinek  \( EF \)}{Oś \( x \)}{Odcinek \( AB \)}{}{}}
\LANGcs{
\Question{
Hyperbola je dána obrázkem v kartézské soustavě souřadnic. Vedlejší osa této hyperboly je:}
{osa \( y \)}{úsečka \( EF \)}{osa \( x \)}{úsečka \( AB \)}{}{}}
\LANGsk{
\Question{
Hyperbola je daná obrázkom v kartézskej sústave súradníc. Vedľajšia os tejto hyperboly je:}
{os \( y \)}{úsečka \( EF \)}{os \( x \)}{úsečka \( AB \)}{}{}}
\LANGes{
\Question{
En la imagen se representa una hipérbola en el sistema de coordenadas cartesianas. Identifica el eje menor de esta hipérbola.}
{El eje \( y \)}{El segmento \( EF \)}{El eje \( x \)}{El segmento \( AB \)}{}{}}
\End

\Data\ID{1103040110}
\Updated{2020-11-08 07:11:55}
\Level{B}
\SubArea{Conic Sections}
\IMG{1103040110.pdf}
\LANGen{
\Question{
The picture shows a hyperbola drawn in the rectangular coordinate system. What is the major axis of this hyperbola?}
{The \( y \)-axis}{The segment \( AB \)}{The segment \( EF \)}{The \( x \)-axis}{}{}}
\LANGpl{
\Question{
Rysunek przedstawia hiperbolę na kartezjańskim układzie współrzędnych. Wskaż oś wielką danej hiperboli.}
{Oś \( y \)}{Odcinek \( AB \)}{Odcinek \( EF \)}{Oś \( x \)}{}{}}
\LANGcs{
\Question{
Hyperbola je dána obrázkem v kartézské soustavě souřadnic. Hlavní osa této hyperboly je:}
{osa \( y \)}{úsečka \( AB \)}{úsečka \( EF \)}{osa \( x \)}{}{}}
\LANGsk{
\Question{
Hyperbola je daná obrázkom v kartézskej sústave súradníc. Hlavná os tejto hyperboly je:}
{os \( y \)}{úsečka \( AB \)}{úsečka \( EF \)}{os \( x \)}{}{}}
\LANGes{
\Question{
En la imagen se representa una hipérbola en el sistema de coordenadas cartesianas. Identifica el eje mayor de esta hipérbola.}
{El eje \( y \)}{El segmento \( AB \)}{El segmento \( EF \)}{El eje \( x \)}{}{}}
\End

\Data\ID{2010005901}
\Updated{2022-05-16 11:05:31}
\Level{B}
\SubArea{Conic Sections}
\LANGen{
\Question{
Find the distance between the points where the
\(y\)-axis
intersects the following hyperbola. 
\[ 
 H\colon \frac{\left (y+3\right )^{2}}
 {36} -\frac{\left (x+4\right )^{2}} 
 {9} = 1
\]}
{\(20\)}{\(16\)}{\(10\)}{\(8\)}{}{}}
\LANGpl{
\Question{
Wyznacz odległość między punktami, w których
oś \(y\) przecina następującą hiperbolę.
\[ 
 H\colon \frac{\left (y+3\right )^{2}}
 {36} -\frac{\left (x+4\right )^{2}} 
 {9} = 1
\]}
{\(20\)}{\(16\)}{\(10\)}{\(8\)}{}{}}
\LANGcs{
\Question{
Určete vzdálenost mezi body, ve kterých osa 
\(y\) protíná hyperbolu:
\[ 
 H\colon \frac{\left (y+3\right )^{2}}
 {36} -\frac{\left (x+4\right )^{2}} 
 {9} = 1
\]}
{\(20\)}{\(16\)}{\(10\)}{\(8\)}{}{}}
\LANGsk{
\Question{
Daná je hyperbola \( H\colon \frac{\left (y+3\right )^{2}} {36} -\frac{\left (x+4\right )^{2}}  {9} = 1. \) Vzdialenosť priesečníkov tejto hyperboly s osou \(y\) je rovná:}
{\(20\)}{\(16\)}{\(10\)}{\(8\)}{}{}}
\LANGes{
\Question{
Encuentra la distancia entre los puntos en los que el eje \(y\)
corta la siguiente hipérbola.  
\[ 
 H\colon \frac{\left (y+3\right )^{2}}
 {36} -\frac{\left (x+4\right )^{2}} 
 {9} = 1
\]}
{\(20\)}{\(16\)}{\(10\)}{\(8\)}{}{}}
\End

\Data\ID{2010005902}
\Updated{2022-05-16 11:05:31}
\Level{B}
\SubArea{Conic Sections}
\LANGen{
\Question{
Find the distance between the points of intersection of the given hyperbola with the given straight line $q$.
\[ 
 H\colon \frac{\left (y+6\right )^{2}}
 {10} -\frac{\left (x-5\right )^{2}} 
 {6} = 1;\quad q\colon y+1 = 0
\]}
{\(6\)}{\(8\)}{\(10\)}{\(12\)}{}{}}
\LANGpl{
\Question{
Wyznacz odległość między punktami przecięcia danej hiperboli z daną linią prostą $q$.
\[ 
 H\colon \frac{\left (y+6\right )^{2}}
 {10} -\frac{\left (x-5\right )^{2}} 
 {6} = 1;\quad q\colon y+1 = 0
\]}
{\(6\)}{\(8\)}{\(10\)}{\(12\)}{}{}}
\LANGcs{
\Question{
Určete vzdálenost průsečíků dané hyperboly s přímkou $q$.
\[ 
 H\colon \frac{\left (y+6\right )^{2}}
 {10} -\frac{\left (x-5\right )^{2}} 
 {6} = 1;\quad q\colon y+1 = 0
\]}
{\(6\)}{\(8\)}{\(10\)}{\(12\)}{}{}}
\LANGsk{
\Question{
Daná je hyperbola \(H\colon \frac{\left (y+6\right )^{2}} {10} -\frac{\left (x-5\right )^{2}}  {6} = 1 \) a priamka \(q\colon y+1 = 0.\) Vzdialenosť priesečníkov tejto hyperboly s priamkou \(q\) sa rovná:}
{\(6\)}{\(8\)}{\(10\)}{\(12\)}{}{}}
\LANGes{
\Question{
Encuentra la distancia entre las intresecciones de la siguiente hipérbola y la recta $q$. 
\[ 
 H\colon \frac{\left (y+6\right )^{2}}
 {10} -\frac{\left (x-5\right )^{2}} 
 {6} = 1;\quad q\colon y+1 = 0
\]}
{\(6\)}{\(8\)}{\(10\)}{\(12\)}{}{}}
\End

\Data\ID{2010005907}
\Updated{2022-05-16 11:05:31}
\Level{B}
\SubArea{Conic Sections}
\LANGen{
\Question{
Find the vertex of the following parabola. 
\[ 
 y^{2} + 12x - 6y - 15 = 0
\]}
{\([2;3]\)}{\([-2;3]\)}{\([2;-3]\)}{\([-2;-3]\)}{}{}}
\LANGpl{
\Question{
Znajdź wierzchołek następującej paraboli.
\[ 
 y^{2} + 12x - 6y - 15 = 0
\]}
{\([2;3]\)}{\([-2;3]\)}{\([2;-3]\)}{\([-2;-3]\)}{}{}}
\LANGcs{
\Question{
Určete vrchol paraboly:
\[ 
 y^{2} + 12x - 6y - 15 = 0
\]}
{\([2;3]\)}{\([-2;3]\)}{\([2;-3]\)}{\([-2;-3]\)}{}{}}
\LANGsk{
\Question{
Vrchol paraboly, danej rovnicou \(  y^{2} + 12x - 6y - 15 = 0,\) má súradnice:}
{\([2;3]\)}{\([-2;3]\)}{\([2;-3]\)}{\([-2;-3]\)}{}{}}
\LANGes{
\Question{
Halla el vértice de la siguiente parábola. 
\[ 
 y^{2} + 12x - 6y - 15 = 0
\]}
{\([2;3]\)}{\([-2;3]\)}{\([2;-3]\)}{\([-2;-3]\)}{}{}}
\End

\Data\ID{2010005909}
\Updated{2022-05-16 11:05:35}
\Level{B}
\SubArea{Conic Sections}
\LANGen{
\Question{
Which of the given points is one of the vertices of the hyperbola  \(25y^{2} - 4x^{2} - 24x + 50y - 111 = 0\)?}
{\([-3;1]\)}{\([3;1]\)}{\([-3;4]\)}{\([3;4]\)}{}{}}
\LANGpl{
\Question{
Który z podanych punktów jest jednym z wierzchołków hiperboli \(25y^{2} - 4x^{2} - 24x + 50y - 111 = 0\)?}
{\([-3;1]\)}{\([3;1]\)}{\([-3;4]\)}{\([3;4]\)}{}{}}
\LANGcs{
\Question{
Který z bodů je vrcholem hyperboly  \(25y^{2} - 4x^{2} - 24x + 50y - 111 = 0\)?}
{\([-3;1]\)}{\([3;1]\)}{\([-3;4]\)}{\([3;4]\)}{}{}}
\LANGsk{
\Question{
Hyperbola je daná rovnicou \(25y^{2} - 4x^{2} - 24x + 50y - 111 = 0\). Jeden z ich vrcholov má súradnice:}
{\([-3;1]\)}{\([3;1]\)}{\([-3;4]\)}{\([3;4]\)}{}{}}
\LANGes{
\Question{
Halla el vértice de la hipérbola \(25y^{2} - 4x^{2} - 24x + 50y - 111 = 0\).}
{\([-3;1]\)}{\([3;1]\)}{\([-3;4]\)}{\([3;4]\)}{}{}}
\End

\Data\ID{2010006008}
\Updated{2022-05-16 11:05:35}
\Level{B}
\SubArea{Conic Sections}
\LANGen{
\Question{
Parabola is a set of points that are equidistant from the point (focus)
and the line (directrix). Find the equation of the directrix of the parabola
\(x^{2} + 4x +8y-20= 0\).}
{\(y-5 = 0\)}{\(y-1 = 0\)}{\(x = 0\)}{\(x+4 = 0\)}{}{}}
\LANGpl{
\Question{
Parabola to zbiór punktów, które są równoodległe od punktu (ognisko)
i prostej (kierownica). Znajdź równanie kierownicy paraboli
\(x^{2} + 4x +8y-20= 0\).}
{\(y-5 = 0\)}{\(y-1 = 0\)}{\(x = 0\)}{\(x+4 = 0\)}{}{}}
\LANGcs{
\Question{
Parabola je množina bodů, které mají stejnou vzdálenost od bodu (ohnisko) a přímky (řidící přímka). Určete rovnici řídící přímky paraboly
\(x^{2} + 4x +8y-20= 0\).}
{\(y-5 = 0\)}{\(y-1 = 0\)}{\(x = 0\)}{\(x+4 = 0\)}{}{}}
\LANGsk{
\Question{
Je daná parabola \(x^{2} + 4x +8y-20= 0.\) Rovnica riadiacej priamky tejto paraboly je:}
{\(y-5 = 0\)}{\(y-1 = 0\)}{\(x = 0\)}{\(x+4 = 0\)}{}{}}
\LANGes{
\Question{
Una parábola es el conjunto de puntos del plano equidistantes a un punto fijo (llamado foco) y a una recta fija (llamada directriz). Encuentra la ecuación de la directriz de la parábola 
\(x^{2} + 4x +8y-20= 0\).}
{\(y-5 = 0\)}{\(y-1 = 0\)}{\(x = 0\)}{\(x+4 = 0\)}{}{}}
\End

\Data\ID{2010006009}
\Updated{2022-05-16 11:05:39}
\Level{B}
\SubArea{Conic Sections}
\LANGen{
\Question{
Find the distance between the point \([5;5]\) 
and the focus of the parabola \(y^{2} +12x + 6y -15 = 0\).}
{\(10\)}{\(8\)}{\(3\sqrt{10}\)}{\(2\sqrt{41}\)}{}{}}
\LANGpl{
\Question{
Znajdź odległość między punktem \([5;5]\)
i ogniskiem paraboli \(y^{2} +12x + 6y -15 = 0\).}
{\(10\)}{\(8\)}{\(3\sqrt{10}\)}{\(2\sqrt{41}\)}{}{}}
\LANGcs{
\Question{
Určete vzdálenost bodu \([5;5]\) od ohniska paraboly \(y^{2} +12x + 6y -15 = 0\).}
{\(10\)}{\(8\)}{\(3\sqrt{10}\)}{\(2\sqrt{41}\)}{}{}}
\LANGsk{
\Question{
Je daná parabola \(y^{2} +12x + 6y -15 = 0\). Vzdialenosť ohniska tejto paraboly od bodu \([5;5]\) je rovná:}
{\(10\)}{\(8\)}{\(3\sqrt{10}\)}{\(2\sqrt{41}\)}{}{}}
\LANGes{
\Question{
Encuentra la distancia entre el punto \([5;5]\) 
y el foco de la parábola \(y^{2} +12x + 6y -15 = 0\).}
{\(10\)}{\(8\)}{\(3\sqrt{10}\)}{\(2\sqrt{41}\)}{}{}}
\End

\Data\ID{2010006302}
\Updated{2022-05-16 11:05:39}
\Level{B}
\SubArea{Conic Sections}
\LANGen{
\Question{
The equation of a hyperbola is given by \( -16x^2+9y^2-96x+108y+36=0\). The length of its semi-minor axis is:}
{\( 3 \)}{\( 9\)}{\( 4 \)}{\( 16 \)}{\( 5 \)}{}}
\LANGpl{
\Question{
Równanie hiperboli jest podane przez \( -16x^2+9y^2-96x+108y+36=0\). Długość jego półosi małej to:}
{\( 3 \)}{\( 9\)}{\( 4 \)}{\( 16 \)}{\( 5 \)}{}}
\LANGcs{
\Question{
Určete délku vedlejší poloosy hyperboly dané rovnicí \( -16x^2+9y^2-96x+108y+36=0\).}
{\( 3 \)}{\( 9\)}{\( 4 \)}{\( 16 \)}{\( 5 \)}{}}
\LANGsk{
\Question{
Hyperbola je daná rovnicou \( -16x^2+9y^2-96x+108y+36=0\). Dĺžka vedľajšej polosi tejto hyperboly je:}
{\( 3 \)}{\( 9\)}{\( 4 \)}{\( 16 \)}{\( 5 \)}{}}
\LANGes{
\Question{
Dada la ecuación de la hipérbola \( -16x^2+9y^2-96x+108y+36=0\). ¿Cuánto mide su semieje menor?}
{\( 3 \)}{\( 9\)}{\( 4 \)}{\( 16 \)}{\( 5 \)}{}}
\End

\Data\ID{2010006303}
\Updated{2022-05-16 11:05:39}
\Level{B}
\SubArea{Conic Sections}
\LANGen{
\Question{
A parabola passes through the points \( A=[5;0] \) and \( B=[-6;2] \) and it is symmetric with respect to \( x\)-axis. What are the coordinates of its vertex?}
{\( [5;0] \)}{\( [6;-2] \)}{\( [0;5] \)}{\( [-6;2] \)}{\( [-1;1] \)}{}}
\LANGpl{
\Question{
Parabola przechodzi przez punkty \( A=[5;0] \) i \( B=[-6;2] \) i jest symetryczna względem osi \( x\). Jakie są współrzędne jej wierzchołka?}
{\( [5;0] \)}{\( [6;-2] \)}{\( [0;5] \)}{\( [-6;2] \)}{\( [-1;1] \)}{}}
\LANGcs{
\Question{
Parabola prochází body  \( A=[5;0] \) a \( B=[-6;2] \) a je souměrná podle osy \( x\). Jaké jsou souřadnice jejího vrcholu?}
{\( [5;0] \)}{\( [6;-2] \)}{\( [0;5] \)}{\( [-6;2] \)}{\( [-1;1] \)}{}}
\LANGsk{
\Question{
Parabola prechádza bodmi \( A=[5;0] \), \( B=[-6;2] \) a je súmerná podľa osi \( x\). Jej vrchol má súradnice:}
{\( [5;0] \)}{\( [6;-2] \)}{\( [0;5] \)}{\( [-6;2] \)}{\( [-1;1] \)}{}}
\LANGes{
\Question{
Una parábola pasa por los puntos \( A=[5;0] \) y \( B=[-6;2] \) y es simétrica respecto al eje \( x\). Halla las coordenadas de su vértice.}
{\( [5;0] \)}{\( [6;-2] \)}{\( [0;5] \)}{\( [-6;2] \)}{\( [-1;1] \)}{}}
\End

\Data\ID{2010006304}
\Updated{2022-05-16 11:05:39}
\Level{B}
\SubArea{Conic Sections}
\LANGen{
\Question{
The equation of a parabola is given by \( y^2-x-6y+10=0 \). What are  the coordinates of its vertex?}
{\( [1;3] \)}{\( [3;1] \)}{\( [3;-1] \)}{\( [-1;-3] \)}{\( [-3;-1] \)}{}}
\LANGpl{
\Question{
Równanie paraboli ma postać \( y^2-x-6y+10=0 \). Jakie są współrzędne jego wierzchołka?}
{\( [1;3] \)}{\( [3;1] \)}{\( [3;-1] \)}{\( [-1;-3] \)}{\( [-3;-1] \)}{}}
\LANGcs{
\Question{
Parabola je dána rovnicí \( y^2-x-6y+10=0 \). Jaké jsou souřadnice jejího vrcholu?}
{\( [1;3] \)}{\( [3;1] \)}{\( [3;-1] \)}{\( [-1;-3] \)}{\( [-3;-1] \)}{}}
\LANGsk{
\Question{
Parabola je daná rovnicou \( y^2-x-6y+10=0 \). Vrchol tejto paraboly má súradnice:}
{\( [1;3] \)}{\( [3;1] \)}{\( [3;-1] \)}{\( [-1;-3] \)}{\( [-3;-1] \)}{}}
\LANGes{
\Question{
Dada la ecuación de la parábola \( y^2-x-6y+10=0 \). Halla las coordenadas de su vértice.}
{\( [1;3] \)}{\( [3;1] \)}{\( [3;-1] \)}{\( [-1;-3] \)}{\( [-3;-1] \)}{}}
\End

\Data\ID{2010006305}
\Updated{2022-05-16 11:05:29}
\Level{B}
\SubArea{Conic Sections}
\LANGen{
\Question{
Decide whether the equation \( 9x^2+4y^2-18x+8y+14=0 \) (in \( x \)-\( y \) plane) describes:}
{an empty set}{an ellipse}{a parabola}{a hyperbola}{a point}{}}
\LANGpl{
\Question{
Zdecyduj, czy równanie \( 9x^2+4y^2-18x+8y+14=0 \) (w płaszczyźnie \( x \)-\( y \)) opisuje:}
{pusty zbiór}{elipsę}{parabolę}{hiperbolę}{punkt}{}}
\LANGcs{
\Question{
Rovnice \( 9x^2+4y^2-18x+8y+14=0 \) (v rovině \( x \)\( y \)) popisuje:}
{prázdnou množinu}{elipsu}{parabolu}{hyperbolu}{bod}{}}
\LANGsk{
\Question{
Rovnica \( 9x^2+4y^2-18x+8y+14=0 \) (v  rovine \( xy \)) popisuje:}
{prázdnu množinu}{elipsu}{parabolu}{hyperbolu}{bod}{}}
\LANGes{
\Question{
Decide qué lugar geométrico describe la ecuación \( 9x^2+4y^2-18x+8y+14=0 \) (en el plano \( x \)-\( y \)):}
{un conjunto vacío}{una elipse}{una parábola}{una hipérbola}{un punto}{}}
\End

\Data\ID{2010006306}
\Updated{2022-05-16 11:05:29}
\Level{B}
\SubArea{Conic Sections}
\LANGen{
\Question{
The equation of a hyperbola with the center \( S=[1;-3] \),  the focus \( F=[1;2] \), and the vertex \( A=[1;0] \) is given by:}
{\( \frac{(y+3)^2}{9}-\frac{(x-1)^2}{16} =1 \)}{\( \frac{(x-1)^2}{16}-\frac{(y+3)^2}{9} =1 \)}{\( \frac{(x-1)^2}{9}-\frac{(y+3)^2}{16} =1 \)}{\( \frac{(y+3)^2}{16}-\frac{(x-1)^2}{9} =1 \)}{\( \frac{(x+1)^2}{16}-\frac{(y-3)^2}{9} =1 \)}{}}
\LANGpl{
\Question{
Równanie hiperboli ze środkiem \( S=[1;-3] \), ogniskiem \( F=[1;2] \) i wierzchołkiem \( A=[1;0] \) jest wyrażone przez:}
{\( \frac{(y+3)^2}{9}-\frac{(x-1)^2}{16} =1 \)}{\( \frac{(x-1)^2}{16}-\frac{(y+3)^2}{9} =1 \)}{\( \frac{(x-1)^2}{9}-\frac{(y+3)^2}{16} =1 \)}{\( \frac{(y+3)^2}{16}-\frac{(x-1)^2}{9} =1 \)}{\( \frac{(x+1)^2}{16}-\frac{(y-3)^2}{9} =1 \)}{}}
\LANGcs{
\Question{
Hyperbola se středem \( S=[1;-3] \), ohniskem \( F=[1;2] \) a vrcholem \( A=[1;0] \) má rovnici:}
{\( \frac{(y+3)^2}{9}-\frac{(x-1)^2}{16} =1 \)}{\( \frac{(x-1)^2}{16}-\frac{(y+3)^2}{9} =1 \)}{\( \frac{(x-1)^2}{9}-\frac{(y+3)^2}{16} =1 \)}{\( \frac{(y+3)^2}{16}-\frac{(x-1)^2}{9} =1 \)}{\( \frac{(x+1)^2}{16}-\frac{(y-3)^2}{9} =1 \)}{}}
\LANGsk{
\Question{
Rovnica hyperboly, ktorá má stred \( S=[1;-3] \),  ohnisko \( F=[1;2] \) a vrchol \( A=[1;0] \) je:}
{\( \frac{(y+3)^2}{9}-\frac{(x-1)^2}{16} =1 \)}{\( \frac{(x-1)^2}{16}-\frac{(y+3)^2}{9} =1 \)}{\( \frac{(x-1)^2}{9}-\frac{(y+3)^2}{16} =1 \)}{\( \frac{(y+3)^2}{16}-\frac{(x-1)^2}{9} =1 \)}{\( \frac{(x+1)^2}{16}-\frac{(y-3)^2}{9} =1 \)}{}}
\LANGes{
\Question{
Elige la ecuación de la hipérbola con centro \( S=[1;-3] \), foco \( F=[1;2] \), y vértice \( A=[1;0] \):}
{\( \frac{(y+3)^2}{9}-\frac{(x-1)^2}{16} =1 \)}{\( \frac{(x-1)^2}{16}-\frac{(y+3)^2}{9} =1 \)}{\( \frac{(x-1)^2}{9}-\frac{(y+3)^2}{16} =1 \)}{\( \frac{(y+3)^2}{16}-\frac{(x-1)^2}{9} =1 \)}{\( \frac{(x+1)^2}{16}-\frac{(y-3)^2}{9} =1 \)}{}}
\End

\Data\ID{2010006307}
\Updated{2022-05-16 11:05:29}
\Level{B}
\SubArea{Conic Sections}
\LANGen{
\Question{
The equation of a parabola is given by \( 4x^2+16x-y+18=0 \). Find the equation of its directrix.}
{\( x=\frac{31}{16} \)}{\( x=-\frac{33}{16} \)}{\( x=\frac{33}{16} \)}{\( x=-\frac{31}{16} \)}{\( x=\frac{15}{8} \)}{}}
\LANGpl{
\Question{
Równanie paraboli ma postać \( 4x^2+16x-y+18=0 \). Znajdź równanie jego kierownicy.}
{\( x=\frac{31}{16} \)}{\( x=-\frac{33}{16} \)}{\( x=\frac{33}{16} \)}{\( x=-\frac{31}{16} \)}{\( x=\frac{15}{8} \)}{}}
\LANGcs{
\Question{
Parabola je dána rovnicí \( 4x^2+16x-y+18=0 \). Určete rovnici řídící přímky.}
{\( x=\frac{31}{16} \)}{\( x=-\frac{33}{16} \)}{\( x=\frac{33}{16} \)}{\( x=-\frac{31}{16} \)}{\( x=\frac{15}{8} \)}{}}
\LANGsk{
\Question{
Parabola je daná rovnicou \( 4x^2+16x-y+18=0 \). Rovnica riadiacej priamky tejto paraboly je:}
{\( x=\frac{31}{16} \)}{\( x=-\frac{33}{16} \)}{\( x=\frac{33}{16} \)}{\( x=-\frac{31}{16} \)}{\( x=\frac{15}{8} \)}{}}
\LANGes{
\Question{
Dada la ecuación de la parábola \( 4x^2+16x-y+18=0 \). Halla la ecuación de su directriz.}
{\( x=\frac{31}{16} \)}{\( x=-\frac{33}{16} \)}{\( x=\frac{33}{16} \)}{\( x=-\frac{31}{16} \)}{\( x=\frac{15}{8} \)}{}}
\End

\Data\ID{9000097001}
\Updated{2020-11-08 09:11:42}
\Level{B}
\SubArea{Conic Sections}
\LANGen{
\Question{
Parabola is a set of the points that are equidistant from the point
(focus) and the line (directrix). Find the directrix of the parabola
\((x - 3)^{2} = 8y\).}
{\(y = -2\)}{\(x = 3\)}{\(x = 0\)}{\(y = 0\)}{}{}}
\LANGpl{
\Question{
Parabola jest zbiorem punktów jednakowo odległych od punktu zwanego ogniskiem
 i prostej zwanej kierownicą.  Wyznacz kierownicę paraboli
\((x - 3)^{2} = 8y\).}
{\(y = -2\)}{\(x = 3\)}{\(x = 0\)}{\(y = 0\)}{}{}}
\LANGcs{
\Question{
Je dána parabola \((x - 3)^{2} = 8y\).
Řídící přímka této paraboly je dána předpisem:}
{\(y = -2\)}{\(x = 3\)}{\(x = 0\)}{\(y = 0\)}{}{}}
\LANGsk{
\Question{
Je daná parabola \((x - 3)^{2} = 8y\).
Riadiaca priamka tejto paraboly je daná predpisom:}
{\(y = -2\)}{\(x = 3\)}{\(x = 0\)}{\(y = 0\)}{}{}}
\LANGes{
\Question{
Una parábola es el conjunto de puntos en el plano equidistante de un punto fijo (llamado foco) y una línea fija (llamada directriz). Encuentra la directriz de la parábola \((x - 3)^{2} = 8y\).}
{\(y = -2\)}{\(x = 3\)}{\(x = 0\)}{\(y = 0\)}{}{}}
\End

\Data\ID{9000097002}
\Updated{2020-11-08 08:11:19}
\Level{B}
\SubArea{Conic Sections}
\LANGen{
\Question{
Consider the parabola \((x - 4)^{2} = 8(y + 1)\).
Find the vertex of this parabola.}
{\([4;-1]\)}{\([4;1]\)}{\([1;4]\)}{\([4;8]\)}{}{}}
\LANGpl{
\Question{
Dana jest  parabola \((x - 4)^{2} = 8(y + 1)\).
Wskaż współrzędne wierzchołka paraboli.}
{\([4;-1]\)}{\([4;1]\)}{\([1;4]\)}{\([4;8]\)}{}{}}
\LANGcs{
\Question{
Je dána parabola \((x - 4)^{2} = 8(y + 1)\).
Vrchol této paraboly má souřadnice:}
{\([4;-1]\)}{\([4;1]\)}{\([1;4]\)}{\([4;8]\)}{}{}}
\LANGsk{
\Question{
Je daná parabola \((x - 4)^{2} = 8(y + 1)\).
Vrchol tejto paraboly má súradnice:}
{\([4;-1]\)}{\([4;1]\)}{\([1;4]\)}{\([4;8]\)}{}{}}
\LANGes{
\Question{
Consideremos la parábola \((x - 4)^{2} = 8(y + 1)\).
Encuentra el vértice de esta parábola.}
{\([4;-1]\)}{\([4;1]\)}{\([1;4]\)}{\([4;8]\)}{}{}}
\End

\Data\ID{9000097003}
\Updated{2020-11-08 08:11:23}
\Level{B}
\SubArea{Conic Sections}
\LANGen{
\Question{
Consider the parabola \((x + 1)^{2} = 4(y - 3)\).
Find the focus of this parabola.}
{\([-1;4]\)}{\([-1;3]\)}{\([3;-1]\)}{\([0;3]\)}{}{}}
\LANGpl{
\Question{
Dana jest  parabola \((x + 1)^{2} = 4(y - 3)\).
Wyznacz współrzędne ogniska paraboli.}
{\([-1;4]\)}{\([-1;3]\)}{\([3;-1]\)}{\([0;3]\)}{}{}}
\LANGcs{
\Question{
Je dána parabola \((x + 1)^{2} = 4(y - 3)\).
Ohnisko této paraboly má souřadnice:}
{\([-1;4]\)}{\([-1;3]\)}{\([3;-1]\)}{\([0;3]\)}{}{}}
\LANGsk{
\Question{
Je daná parabola \((x + 1)^{2} = 4(y - 3)\).
Ohnisko tejto paraboly má súradnice:}
{\([-1;4]\)}{\([-1;3]\)}{\([3;-1]\)}{\([0;3]\)}{}{}}
\LANGes{
\Question{
Consideremos la parábola \((x + 1)^{2} = 4(y - 3)\).
Encuentra el foco de la parábola.}
{\([-1;4]\)}{\([-1;3]\)}{\([3;-1]\)}{\([0;3]\)}{}{}}
\End

\Data\ID{9000097004}
\Updated{2020-11-08 09:11:32}
\Level{B}
\SubArea{Conic Sections}
\LANGen{
\Question{
Parabola is a set of the points that are equidistant from the point
(focus) and the line (directrix). Find the directrix of the parabola
\((x + 2)^{2} = -8(y - 1)\).}
{\(y = 3\)}{\(x = -2\)}{\(x = 1\)}{\(y = -4\)}{}{}}
\LANGpl{
\Question{
Parabola jest zbiorem punktów jednakowo odległych od punktu zwanego ogniskiem
i prostej zwanej kierownicą.  Wyznacz kierownicę paraboli.
\((x + 2)^{2} = -8(y - 1)\).}
{\(y = 3\)}{\(x = -2\)}{\(x = 1\)}{\(y = -4\)}{}{}}
\LANGcs{
\Question{
Je dána parabola \((x + 2)^{2} = -8(y - 1)\).
Řídící přímka této paraboly je dána předpisem:}
{\(y = 3\)}{\(x = -2\)}{\(x = 1\)}{\(y = -4\)}{}{}}
\LANGsk{
\Question{
Je daná parabola \((x + 2)^{2} = -8(y - 1)\).
Riadiaca priamka tejto paraboly je daná predpisom:}
{\(y = 3\)}{\(x = -2\)}{\(x = 1\)}{\(y = -4\)}{}{}}
\LANGes{
\Question{
Una parábola es el conjunto de puntos en el plano equidistante de un punto fijo (llamado foco) y una línea fija (llamada directriz). Encuentra la directriz de la parábola
\((x + 2)^{2} = -8(y - 1)\).}
{\(y = 3\)}{\(x = -2\)}{\(x = 1\)}{\(y = -4\)}{}{}}
\End

\Data\ID{9000097005}
\Updated{2020-11-07 08:11:09}
\Level{B}
\SubArea{Conic Sections}
\LANGen{
\Question{
Consider the parabola \((x - 5)^{2} = -8(y - 3)\).
Find the vertex of this parabola.}
{\([5;3]\)}{\([3;5]\)}{\([5;-4]\)}{\([-4;3]\)}{}{}}
\LANGpl{
\Question{
Dana jest parabola \((x - 5)^{2} = -8(y - 3)\).
Wskaż współrzędne wierzchołka paraboli.}
{\([5;3]\)}{\([3;5]\)}{\([5;-4]\)}{\([-4;3]\)}{}{}}
\LANGcs{
\Question{
Je dána parabola \((x - 5)^{2} = -8(y - 3)\).
Vrchol této paraboly má souřadnice:}
{\([5;3]\)}{\([3;5]\)}{\([5;-4]\)}{\([-4;3]\)}{}{}}
\LANGsk{
\Question{
Je daná parabola \((x - 5)^{2} = -8(y - 3)\).
Vrchol tejto paraboly má súradnice:}
{\([5;3]\)}{\([3;5]\)}{\([5;-4]\)}{\([-4;3]\)}{}{}}
\LANGes{
\Question{
Consideremos la parábola \((x - 5)^{2} = -8(y - 3)\).
Halla el vértice de esta parábola.}
{\([5;3]\)}{\([3;5]\)}{\([5;-4]\)}{\([-4;3]\)}{}{}}
\End

\Data\ID{9000097006}
\Updated{2020-11-08 08:11:14}
\Level{B}
\SubArea{Conic Sections}
\LANGen{
\Question{
Consider the parabola \((y - 1)^{2} = 12(x - 1)\).
Find the focus of this parabola.}
{\([4;1]\)}{\([1;6]\)}{\([3;1]\)}{\([1;1]\)}{}{}}
\LANGpl{
\Question{
Dana jest  \((y - 1)^{2} = 12(x - 1)\).
Wyznacz współrzędne ogniska paraboli.}
{\([4;1]\)}{\([1;6]\)}{\([3;1]\)}{\([1;1]\)}{}{}}
\LANGcs{
\Question{
Je dána parabola \((y - 1)^{2} = 12(x - 1)\).
Ohnisko této paraboly má souřadnice:}
{\([4;1]\)}{\([1;6]\)}{\([3;1]\)}{\([1;1]\)}{}{}}
\LANGsk{
\Question{
Je daná parabola \((y - 1)^{2} = 12(x - 1)\).
Ohnisko tejto paraboly má súradnice:}
{\([4;1]\)}{\([1;6]\)}{\([3;1]\)}{\([1;1]\)}{}{}}
\LANGes{
\Question{
Consideremos la parábola \((y - 1)^{2} = 12(x - 1)\).
Encuentra el foco de la parábola.}
{\([4;1]\)}{\([1;6]\)}{\([3;1]\)}{\([1;1]\)}{}{}}
\End

\Data\ID{9000097007}
\Updated{2020-11-08 09:11:06}
\Level{B}
\SubArea{Conic Sections}
\LANGen{
\Question{
Parabola is a set of the points that are equidistant from the point
(focus) and the line (directrix). Find the directrix of the parabola
\((y - 4)^{2} = 8(x - 1)\).}
{\(x = -1\)}{\(x = 4\)}{\(y = 4\)}{\(y = 2\)}{}{}}
\LANGpl{
\Question{
Parabola jest zbiorem punktów jednakowo odległych od punktu zwanego ogniskiem
i prostej zwanej kierownicą.  Wyznacz kierownicę paraboli.
\((y - 4)^{2} = 8(x - 1)\).}
{\(x = -1\)}{\(x = 4\)}{\(y = 4\)}{\(y = 2\)}{}{}}
\LANGcs{
\Question{
Je dána parabola \((y - 4)^{2} = 8(x - 1)\).
Řídící přímka této paraboly je dána předpisem:}
{\(x = -1\)}{\(x = 4\)}{\(y = 4\)}{\(y = 2\)}{}{}}
\LANGsk{
\Question{
Je daná parabola \((y - 4)^{2} = 8(x - 1)\).
Riadiaca priamka tejto paraboly je daná predpisom:}
{\(x = -1\)}{\(x = 4\)}{\(y = 4\)}{\(y = 2\)}{}{}}
\LANGes{
\Question{
Una parábola es el conjunto de puntos en el plano equidistante de un punto fijo (llamado foco) y una línea fija (llamada directriz). Encuentra la directriz de la parábola
\((y - 4)^{2} = 8(x - 1)\).}
{\(x = -1\)}{\(x = 4\)}{\(y = 4\)}{\(y = 2\)}{}{}}
\End

\Data\ID{9000097008}
\Updated{2020-11-08 08:11:46}
\Level{B}
\SubArea{Conic Sections}
\LANGen{
\Question{
Consider the parabola \((y - 2)^{2} = -12(x + 1)\).
Find the vertex of this parabola.}
{\([-1;2]\)}{\([2;-1]\)}{\([2;-6]\)}{\([-6;-1]\)}{}{}}
\LANGpl{
\Question{
Dana jest parabola \((y - 2)^{2} = -12(x + 1)\).
Wskaż współrzędne wierzchołka paraboli.}
{\([-1;2]\)}{\([2;-1]\)}{\([2;-6]\)}{\([-6;-1]\)}{}{}}
\LANGcs{
\Question{
Je dána parabola \((y - 2)^{2} = -12(x + 1)\).
Vrchol této paraboly má souřadnice:}
{\([-1;2]\)}{\([2;-1]\)}{\([2;-6]\)}{\([-6;-1]\)}{}{}}
\LANGsk{
\Question{
Je daná parabola \((y - 2)^{2} = -12(x + 1)\).
Vrchol tejto paraboly má súradnice:}
{\([-1;2]\)}{\([2;-1]\)}{\([2;-6]\)}{\([-6;-1]\)}{}{}}
\LANGes{
\Question{
Consideremos la parábola \((y - 2)^{2} = -12(x + 1)\).
Encuentra el vértice de esta parábola.}
{\([-1;2]\)}{\([2;-1]\)}{\([2;-6]\)}{\([-6;-1]\)}{}{}}
\End

\Data\ID{9000097009}
\Updated{2020-11-08 09:11:17}
\Level{B}
\SubArea{Conic Sections}
\LANGen{
\Question{
Consider the parabola \((y - 3)^{2} = -4(x + 2)\).
Find the focus of this parabola}
{\([-3;3]\)}{\([-2;3]\)}{\([3;-2]\)}{\([3;-4]\)}{}{}}
\LANGpl{
\Question{
Dana jest parabola \((y - 3)^{2} = -4(x + 2)\).
Wyznacz współrzędne ogniska paraboli.}
{\([-3;3]\)}{\([-2;3]\)}{\([3;-2]\)}{\([3;-4]\)}{}{}}
\LANGcs{
\Question{
Je dána parabola \((y - 3)^{2} = -4(x + 2)\).
Ohnisko této paraboly má souřadnice:}
{\([-3;3]\)}{\([-2;3]\)}{\([3;-2]\)}{\([3;-4]\)}{}{}}
\LANGsk{
\Question{
Je daná parabola \((y - 3)^{2} = -4(x + 2)\).
Ohnisko tejto paraboly má súradnice:}
{\([-3;3]\)}{\([-2;3]\)}{\([3;-2]\)}{\([3;-4]\)}{}{}}
\LANGes{
\Question{
Consideremos la parábola \((y - 3)^{2} = -4(x + 2)\).
Encuentra el foco de la parábola.}
{\([-3;3]\)}{\([-2;3]\)}{\([3;-2]\)}{\([3;-4]\)}{}{}}
\End

\Data\ID{9000097010}
\Updated{2020-11-08 09:11:55}
\Level{B}
\SubArea{Conic Sections}
\LANGen{
\Question{
Parabola is a set of the points that are equidistant from the point
(focus) and the line (directrix). Find the directrix of the parabola
\((y + 3)^{2} = -8(x + 4)\).}
{\(x = -2\)}{\(x = -4\)}{\(y = -3\)}{\(y = -2\)}{}{}}
\LANGpl{
\Question{
Parabola jest zbiorem punktów jednakowo odległych od punktu zwanego ogniskiem
i prostej zwanej kierownicą.  Wyznacz kierownicę paraboli.
\((y + 3)^{2} = -8(x + 4)\).}
{\(x = -2\)}{\(x = -4\)}{\(y = -3\)}{\(y = -2\)}{}{}}
\LANGcs{
\Question{
Je dána parabola \((y + 3)^{2} = -8(x + 4)\).
Řídící přímka této paraboly je dána předpisem:}
{\(x = -2\)}{\(x = -4\)}{\(y = -3\)}{\(y = -2\)}{}{}}
\LANGsk{
\Question{
Je daná parabola \((y + 3)^{2} = -8(x + 4)\).
Riadiaca priamka tejto paraboly je daná predpisom:}
{\(x = -2\)}{\(x = -4\)}{\(y = -3\)}{\(y = -2\)}{}{}}
\LANGes{
\Question{
Una parábola es el conjunto de puntos en el plano equidistante de un punto fijo (llamado foco) y una línea fija (llamada directriz). Encuentra la directriz de la parábola
\((y + 3)^{2} = -8(x + 4)\).}
{\(x = -2\)}{\(x = -4\)}{\(y = -3\)}{\(y = -2\)}{}{}}
\End

\Data\ID{9000105401}
\Updated{2021-10-14 06:10:01}
\Level{B}
\SubArea{Conic Sections}
\LANGen{
\Question{
The parabola \(P\colon x^{2} - 6x - 4y + 5 = 0\) has
two \(x\)-intercepts.
Find their distance.}
{\(4\)}{\(6\)}{\(8\)}{\(10\)}{}{}}
\LANGpl{
\Question{
Parabola \(P\colon x^{2} - 6x - 4y + 5 = 0\) przecina oś  \(x\) w dwóch punktach.
Wskaż odległość pomiędzy nimi.}
{\(4\)}{\(6\)}{\(8\)}{\(10\)}{}{}}
\LANGcs{
\Question{
Parabola \(P\colon x^{2} - 6x - 4y + 5 = 0\) 
protíná osu \(x\) 
ve dvou bodech. Jejich vzdálenost je:}
{\(4\)}{\(6\)}{\(8\)}{\(10\)}{}{}}
\LANGsk{
\Question{
Parabola \(P\colon x^{2} - 6x - 4y + 5 = 0\) 
pretína os \(x\) 
v dvoch bodoch. Ich vzdialenosť je:}
{\(4\)}{\(6\)}{\(8\)}{\(10\)}{}{}}
\LANGes{
\Question{
La parábola \(P\colon x^{2} - 6x - 4y + 5 = 0\) corta el eje \(x\) en dos puntos. Calcula la distancia entre estos dos puntos.}
{\(4\)}{\(6\)}{\(8\)}{\(10\)}{}{}}
\End

\Data\ID{9000105402}
\Updated{2021-10-14 06:10:17}
\Level{B}
\SubArea{Conic Sections}
\LANGen{
\Question{
The parabola \(P\colon x^{2} - 4x - 10y - 21 = 0\) has
two \(x\)-intercepts.
Find their distance.}
{\(10\)}{\(12\)}{\(8\)}{\(6\)}{}{}}
\LANGpl{
\Question{
Parabola \(P\colon x^{2} - 4x - 10y - 21 = 0\) przecina oś \(x\) w dwóch punktach.
Wskaż odległość pomiędzy nimi.}
{\(10\)}{\(12\)}{\(8\)}{\(6\)}{}{}}
\LANGcs{
\Question{
Parabola \(P\colon x^{2} - 4x - 10y - 21 = 0\) 
protíná osu \(x\) 
ve dvou bodech. Jejich vzdálenost je:}
{\(10\)}{\(12\)}{\(8\)}{\(6\)}{}{}}
\LANGsk{
\Question{
Parabola \(P\colon x^{2} - 4x - 10y - 21 = 0\) 
pretína os \(x\) 
v dvoch bodoch. Ich vzdialenosť je:}
{\(10\)}{\(12\)}{\(8\)}{\(6\)}{}{}}
\LANGes{
\Question{
La parábola \(P\colon x^{2} - 4x - 10y - 21 = 0\) corta el eje \(x\) en dos puntos. Calcula la distancia entre estos dos puntos.}
{\(10\)}{\(12\)}{\(8\)}{\(6\)}{}{}}
\End

\Data\ID{9000105403}
\Updated{2021-10-14 06:10:41}
\Level{B}
\SubArea{Conic Sections}
\LANGen{
\Question{
The parabola \(P\colon y^{2} + 2y + 10x - 24 = 0\) has
two \(y\)-intercepts.
Find their distance.}
{\(10\)}{\(8\)}{\(6\)}{\(4\)}{}{}}
\LANGpl{
\Question{
Parabola \(P\colon y^{2} + 2y + 10x - 24 = 0\) przecina
oś  \(y\) w dwóch punktach.
Wskaż odległość pomiędzy nimi.}
{\(10\)}{\(8\)}{\(6\)}{\(4\)}{}{}}
\LANGcs{
\Question{
Parabola \(P\colon y^{2} + 2y + 10x - 24 = 0\) 
protíná osu \(y\) 
ve dvou bodech. Jejich vzdálenost je:}
{\(10\)}{\(8\)}{\(6\)}{\(4\)}{}{}}
\LANGsk{
\Question{
Parabola \(P\colon y^{2} + 2y + 10x - 24 = 0\) 
pretína os \(y\) 
v dvoch bodoch. Ich vzdialenosť je:}
{\(10\)}{\(8\)}{\(6\)}{\(4\)}{}{}}
\LANGes{
\Question{
La parábola \(P\colon y^{2} + 2y + 10x - 24 = 0\) corta el eje \(y\) en dos puntos. Calcula la distancia entre estos dos puntos.}
{\(10\)}{\(8\)}{\(6\)}{\(4\)}{}{}}
\End

\Data\ID{9000105404}
\Updated{2021-10-14 06:10:56}
\Level{B}
\SubArea{Conic Sections}
\LANGen{
\Question{
The parabola \(P\colon y^{2} + 4y + 6x - 5 = 0\) has
two \(y\)-intercepts.
Find their distance.}
{\(6\)}{\(8\)}{\(10\)}{\(4\)}{}{}}
\LANGpl{
\Question{
Parabola \(P\colon y^{2} + 4y + 6x - 5 = 0\) przecina oś
 \(y\) w dwóch punktach.
Wskaż odległość pomiędzy nimi.}
{\(6\)}{\(8\)}{\(10\)}{\(4\)}{}{}}
\LANGcs{
\Question{
Parabola \(P\colon y^{2} + 4y + 6x - 5 = 0\) 
protíná osu \(y\) 
ve dvou bodech. Jejich vzdálenost je:}
{\(6\)}{\(8\)}{\(10\)}{\(4\)}{}{}}
\LANGsk{
\Question{
Parabola \(P\colon y^{2} + 4y + 6x - 5 = 0\) 
pretína os \(y\) 
v dvoch bodoch. Ich vzdialenosť je:}
{\(6\)}{\(8\)}{\(10\)}{\(4\)}{}{}}
\LANGes{
\Question{
La  parábola \(P\colon y^{2} + 4y + 6x - 5 = 0\) corta el eje \(y\) en dos puntos. Calcula la distancia entre estos dos puntos.}
{\(6\)}{\(8\)}{\(10\)}{\(4\)}{}{}}
\End

\Data\ID{9000105405}
\Updated{2021-10-14 07:10:22}
\Level{B}
\SubArea{Conic Sections}
\LANGen{
\Question{
Parabola is a set of the points that are equidistant from the point (focus)
and the line (directrix). Find the equation of the directrix of the parabola
\(P\colon x^{2} - 4x - 6y - 17 = 0\).}
{\(y + 5 = 0\)}{\(y - 3 = 0\)}{\(x + 4 = 0\)}{\(x - 2 = 0\)}{}{}}
\LANGpl{
\Question{
Parabola jest zbiorem punktów jednakowo odległych od punktu zwanego ogniskiem
i prostej zwanej kierownicą.  Wskaż równanie określające kierownicę paraboli
\(P\colon x^{2} - 4x - 6y - 17 = 0\).}
{\(y + 5 = 0\)}{\(y - 3 = 0\)}{\(x + 4 = 0\)}{\(x - 2 = 0\)}{}{}}
\LANGcs{
\Question{
Je dána parabola \(P\colon x^{2} - 4x - 6y - 17 = 0\).
Rovnice řídící přímky této paraboly je:}
{\(d\colon y + 5 = 0\)}{\(d\colon y - 3 = 0\)}{\(d\colon x + 4 = 0\)}{\(d\colon x - 2 = 0\)}{}{}}
\LANGsk{
\Question{
Je daná parabola \(P\colon x^{2} - 4x - 6y - 17 = 0\).
Rovnica riadiacej priamky tejto paraboly je:}
{\(d\colon y + 5 = 0\)}{\(d\colon y - 3 = 0\)}{\(d\colon x + 4 = 0\)}{\(d\colon x - 2 = 0\)}{}{}}
\LANGes{
\Question{
Una parábola es el conjunto de puntos del plano equidistantes a un punto fijo (llamado foco) y a una recta fija (llamada directriz). Encuentra la ecuación de la directriz de la parábola
\(P\colon x^{2} - 4x - 6y - 17 = 0\).}
{\(y + 5 = 0\)}{\(y - 3 = 0\)}{\(x + 4 = 0\)}{\(x - 2 = 0\)}{}{}}
\End

\Data\ID{9000105406}
\Updated{2021-10-14 07:10:47}
\Level{B}
\SubArea{Conic Sections}
\LANGen{
\Question{
Parabola is a set of the points that are equidistant from the point (focus)
and the line (directrix). Find the equation of the directrix of the parabola
\(P\colon y^{2} + 4y + 4x - 4 = 0\).}
{\(x - 3 = 0\)}{\(x + 4 = 0\)}{\(y + 1 = 0\)}{\(y - 2 = 0\)}{}{}}
\LANGpl{
\Question{
Parabola jest zbiorem punktów jednakowo odległych od punktu zwanego ogniskiem
i prostej zwanej kierownicą.  Wskaż równanie określające kierownicę paraboli
\(P\colon y^{2} + 4y + 4x - 4 = 0\).}
{\(x - 3 = 0\)}{\(x + 4 = 0\)}{\(y + 1 = 0\)}{\(y - 2 = 0\)}{}{}}
\LANGcs{
\Question{
Je dána parabola \(P\colon y^{2} + 4y + 4x - 4 = 0\).
Rovnice řídící přímky této paraboly je:}
{\(d\colon x - 3 = 0\)}{\(d\colon x + 4 = 0\)}{\(d\colon y + 1 = 0\)}{\(d\colon y - 2 = 0\)}{}{}}
\LANGsk{
\Question{
Je daná parabola \(P\colon y^{2} + 4y + 4x - 4 = 0\).
Rovnica riadiacej priamky tejto paraboly je:}
{\(d\colon x - 3 = 0\)}{\(d\colon x + 4 = 0\)}{\(d\colon y + 1 = 0\)}{\(d\colon y - 2 = 0\)}{}{}}
\LANGes{
\Question{
Una parábola es el conjunto de puntos del plano equidistantes a un punto fijo (llamado foco) y a una recta fija (llamada directriz). Encuentra la ecuación de la directriz de la parábola
\(P\colon y^{2} + 4y + 4x - 4 = 0\).}
{\(x - 3 = 0\)}{\(x + 4 = 0\)}{\(y + 1 = 0\)}{\(y - 2 = 0\)}{}{}}
\End

\Data\ID{9000105407}
\Updated{2022-04-23 05:04:26}
\Level{B}
\SubArea{Conic Sections}
\LANGen{
\Question{
Parabola is a set of points that are equidistant from the point (focus)
and the line (directrix). Find the equation of the directrix of the parabola
\(y^{2} + 6y - 12x + 21 = 0\).}
{\(x + 2 = 0\)}{\(x - 5 = 0\)}{\(y + 3 = 0\)}{\(y - 1 = 0\)}{}{}}
\LANGpl{
\Question{
Parabola jest zbiorem punktów jednakowo odległych od punktu zwanego ogniskiem
i prostej zwanej kierownicą.  Wskaż równanie określające kierownicę paraboli
\(y^{2} + 6y - 12x + 21 = 0\).}
{\(x + 2 = 0\)}{\(x - 5 = 0\)}{\(y + 3 = 0\)}{\(y - 1 = 0\)}{}{}}
\LANGcs{
\Question{
Je dána parabola \(y^{2} + 6y - 12x + 21 = 0\).
Rovnice řídící přímky této paraboly je:}
{\(d\colon x + 2 = 0\)}{\(d\colon x - 5 = 0\)}{\(d\colon y + 3 = 0\)}{\(d\colon y - 1 = 0\)}{}{}}
\LANGsk{
\Question{
Je daná parabola \(y^{2} + 6y - 12x + 21 = 0\).
Rovnica riadiacej priamky tejto paraboly je:}
{\(d\colon x + 2 = 0\)}{\(d\colon x - 5 = 0\)}{\(d\colon y + 3 = 0\)}{\(d\colon y - 1 = 0\)}{}{}}
\LANGes{
\Question{
Una parábola es el conjunto de puntos del plano equidistantes a un punto fijo (llamado foco) y a una recta fija (llamada directriz). Encuentra la ecuación de la directriz de la parábola
\(P\colon y^{2} + 6y - 12x + 21 = 0\).}
{\(x + 2 = 0\)}{\(x - 5 = 0\)}{\(y + 3 = 0\)}{\(y - 1 = 0\)}{}{}}
\End

\Data\ID{9000105408}
\Updated{2021-10-14 07:10:23}
\Level{B}
\SubArea{Conic Sections}
\LANGen{
\Question{
Parabola is a set of the points that are equidistant from the point (focus)
and the line (directrix). Find the equation of the directrix of the parabola
\(P\colon x^{2} - 8x + 6y + 19 = 0\).}
{\(y - 1 = 0\)}{\(y + 2 = 0\)}{\(x + 4 = 0\)}{\(x - 3 = 0\)}{}{}}
\LANGpl{
\Question{
Parabola jest zbiorem punktów jednakowo odległych od punktu zwanego ogniskiem
i prostej zwanej kierownicą.  Wskaż równanie określające kierownicę paraboli
\(P\colon x^{2} - 8x + 6y + 19 = 0\).}
{\(y - 1 = 0\)}{\(y + 2 = 0\)}{\(x + 4 = 0\)}{\(x - 3 = 0\)}{}{}}
\LANGcs{
\Question{
Je dána parabola \(P\colon x^{2} - 8x + 6y + 19 = 0\).
Rovnice řídící přímky této paraboly je:}
{\(d\colon y - 1 = 0\)}{\(d\colon y + 2 = 0\)}{\(d\colon x + 4 = 0\)}{\(d\colon x - 3 = 0\)}{}{}}
\LANGsk{
\Question{
Je daná parabola \(P\colon x^{2} - 8x + 6y + 19 = 0\).
Rovnica riadiacej priamky tejto paraboly je:}
{\(d\colon y - 1 = 0\)}{\(d\colon y + 2 = 0\)}{\(d\colon x + 4 = 0\)}{\(d\colon x - 3 = 0\)}{}{}}
\LANGes{
\Question{
Una parábola es el conjunto de puntos del plano equidistantes a un punto fijo (llamado foco) y a una recta fija (llamada directriz). Encuentra la ecuación de la directriz de la parábola
\(P\colon x^{2} - 8x + 6y + 19 = 0\).}
{\(y - 1 = 0\)}{\(y + 2 = 0\)}{\(x + 4 = 0\)}{\(x - 3 = 0\)}{}{}}
\End

\Data\ID{9000105409}
\Updated{2022-04-23 05:04:26}
\Level{B}
\SubArea{Conic Sections}
\LANGen{
\Question{
Find the distance between the point \([7;3]\) 
and the focus of the parabola \(x^{2} - 8x + 8y + 8 = 0\).}
{\(5\)}{\(9\sqrt5\)}{\(3\)}{\(\sqrt{13}\)}{}{}}
\LANGpl{
\Question{
Wskaż odległość pomiędzy punktem \([7;3]\), 
a ogniskiem paraboli \(x^{2} - 8x + 8y + 8 = 0\).}
{\(5\)}{\(9\sqrt5\)}{\(3\)}{\(\sqrt{13}\)}{}{}}
\LANGcs{
\Question{
Je dána parabola \(x^{2} - 8x + 8y + 8 = 0\).
Vzdálenost ohniska této paraboly od bodu
\([7;3]\) je
rovna:}
{\(5\)}{\(9\sqrt5\)}{\(3\)}{\(\sqrt{13}\)}{}{}}
\LANGsk{
\Question{
Je daná parabola \(x^{2} - 8x + 8y + 8 = 0\).
Vzdialenosť ohniska tejto paraboly od bodu
\([7;3]\) je
rovná:}
{\(5\)}{\(9\sqrt5\)}{\(3\)}{\(\sqrt{13}\)}{}{}}
\LANGes{
\Question{
Encuentra la distancia entre el punto \([7;3]\) 
y el foco de la parábola \(x^{2} - 8x + 8y + 8 = 0\).}
{\(5\)}{\(9\sqrt5\)}{\(3\)}{\(\sqrt{13}\)}{}{}}
\End

\Data\ID{9000105410}
\Updated{2020-11-08 09:11:39}
\Level{B}
\SubArea{Conic Sections}
\LANGen{
\Question{
Find the distance between the point \(X = [-3;-6]\) 
and the focus of the parabola \(P\colon y^{2} + 2y - 24x + 73 = 0\).}
{\(13\)}{\(12\)}{\(5\)}{\(1\)}{}{}}
\LANGpl{
\Question{
Wskaż odległość pomiędzy punktem \(X = [-3;-6]\), 
a ogniskiem paraboli  \(P\colon y^{2} + 2y - 24x + 73 = 0\).}
{\(13\)}{\(12\)}{\(5\)}{\(1\)}{}{}}
\LANGcs{
\Question{
Je dána parabola \(P\colon y^{2} + 2y - 24x + 73 = 0\).
Vzdálenost ohniska této paraboly od bodu
\(X = [-3;-6]\) je
rovna:}
{\(13\)}{\(12\)}{\(5\)}{\(1\)}{}{}}
\LANGsk{
\Question{
Je daná parabola \(P\colon y^{2} + 2y - 24x + 73 = 0\).
Vzdialenosť ohniska tejto paraboly od bodu
\(X = [-3;-6]\) je
rovná:}
{\(13\)}{\(12\)}{\(5\)}{\(1\)}{}{}}
\LANGes{
\Question{
Encuentra la distancia entre el punto  \(X = [-3;-6]\) 
y el foco de la parábola \(P\colon y^{2} + 2y - 24x + 73 = 0\).}
{\(13\)}{\(12\)}{\(5\)}{\(1\)}{}{}}
\End

\Data\ID{9000105501}
\Updated{2022-04-23 05:04:20}
\Level{B}
\SubArea{Conic Sections}
\LANGen{
\Question{
Find the distance between the points where the
\(x\)-axis
intersects the following hyperbola. 
\[ 
 H\colon \frac{\left (x - 3\right )^{2}}
 {20} -\frac{\left (y - 2\right )^{2}} 
 {5} = 1
\]}
{\(12\)}{\(14\)}{\(10\)}{\(8\)}{}{}}
\LANGpl{
\Question{
Wskaż  odległość punktów  przecięcia hiperboli $H$ z osią \(x\).
\[ 
 H\colon \frac{\left (x - 3\right )^{2}}
 {20} -\frac{\left (y - 2\right )^{2}} 
 {5} = 1
\]}
{\(12\)}{\(14\)}{\(10\)}{\(8\)}{}{}}
\LANGcs{
\Question{
Je dána hyperbola \(H\colon \frac{\left (x-3\right )^{2}} 
 {20} -\frac{\left (y-2\right )^{2}} 
 {5} = 1\).
Vzdálenost průsečíků této hyperboly s osou
\(x\) je
rovna:}
{\(12\)}{\(14\)}{\(10\)}{\(8\)}{}{}}
\LANGsk{
\Question{
Je daná hyperbola \(H\colon \frac{\left (x-3\right )^{2}} 
 {20} -\frac{\left (y-2\right )^{2}} 
 {5} = 1\).
Vzdialenosť priesečníkov tejto hyperboly s osou
\(x\) je
rovná:}
{\(12\)}{\(14\)}{\(10\)}{\(8\)}{}{}}
\LANGes{
\Question{
Encuentra la distancia entre los puntos en los que el eje \(x\)
corta a la siguiente hipérbola.  
\[ 
 H\colon \frac{\left (x - 3\right )^{2}}
 {20} -\frac{\left (y - 2\right )^{2}} 
 {5} = 1
\]}
{\(12\)}{\(14\)}{\(10\)}{\(8\)}{}{}}
\End

\Data\ID{9000105502}
\Updated{2021-10-14 07:10:14}
\Level{B}
\SubArea{Conic Sections}
\LANGen{
\Question{
Find the distance between the points where the
\(x\)-axis
intersects the following hyperbola. 
\[ 
 H\colon \frac{\left (x - 1\right )^{2}}
 {10} -\frac{\left (y - 3\right )^{2}} 
 {6} = 1
\]}
{\(10\)}{\(14\)}{\(12\)}{\(8\)}{}{}}
\LANGpl{
\Question{
Wskaż  odległość punktów  przecięcia hiperboli $H$ z osią \(x\).
\[ 
 H\colon \frac{\left (x - 1\right )^{2}}
 {10} -\frac{\left (y - 3\right )^{2}} 
 {6} = 1
\]}
{\(10\)}{\(14\)}{\(12\)}{\(8\)}{}{}}
\LANGcs{
\Question{
Je dána hyperbola \(H\colon \frac{\left (x-1\right )^{2}} 
 {10} -\frac{\left (y-3\right )^{2}} 
 {6} = 1\).
Vzdálenost průsečíků této hyperboly s osou
\(x\) je
rovna:}
{\(10\)}{\(14\)}{\(12\)}{\(8\)}{}{}}
\LANGsk{
\Question{
Je daná hyperbola \(H\colon \frac{\left (x-1\right )^{2}} 
 {10} -\frac{\left (y-3\right )^{2}} 
 {6} = 1\).
Vzdialenosť priesečníkov tejto hyperboly s osou
\(x\) je
rovná:}
{\(10\)}{\(14\)}{\(12\)}{\(8\)}{}{}}
\LANGes{
\Question{
Encuentra la distancia entre los puntos en los que el eje \(x\) corta a la siguiente hipérbola. 
\[ 
 H\colon \frac{\left (x - 1\right )^{2}}
 {10} -\frac{\left (y - 3\right )^{2}} 
 {6} = 1
\]}
{\(10\)}{\(14\)}{\(12\)}{\(8\)}{}{}}
\End

\Data\ID{9000105503}
\Updated{2021-10-14 07:10:25}
\Level{B}
\SubArea{Conic Sections}
\LANGen{
\Question{
Find the distance between the points where the
\(y\)-axis
intersects the following hyperbola. 
\[ 
 H\colon \frac{\left (x - 4\right )^{2}}
 {8} -\frac{\left (y - 3\right )^{2}} 
 {1} = 1
\]}
{\(2\)}{\(4\)}{\(6\)}{\(8\)}{}{}}
\LANGpl{
\Question{
Wskaż  odległość punktów  przecięcia hiperboli $H$ z osią \(y\).
\[ 
 H\colon \frac{\left (x - 4\right )^{2}}
 {8} -\frac{\left (y - 3\right )^{2}} 
 {1} = 1
\]}
{\(2\)}{\(4\)}{\(6\)}{\(8\)}{}{}}
\LANGcs{
\Question{
Je dána hyperbola \(H\colon \frac{\left (x-4\right )^{2}} 
 {8} -\frac{\left (y-3\right )^{2}} 
 {1} = 1\).
Vzdálenost průsečíků této hyperboly s osou
\(y\) je
rovna:}
{\(2\)}{\(4\)}{\(6\)}{\(8\)}{}{}}
\LANGsk{
\Question{
Je daná hyperbola \(H\colon \frac{\left (x-4\right )^{2}} 
 {8} -\frac{\left (y-3\right )^{2}} 
 {1} = 1\).
Vzdialenosť priesečníkov tejto hyperboly s osou
\(y\) je
rovná:}
{\(2\)}{\(4\)}{\(6\)}{\(8\)}{}{}}
\LANGes{
\Question{
Encuentra la distancia entre los puntos en los que el eje \(y\)
corta a la siguiente hipérbola.  
\[ 
 H\colon \frac{\left (x - 4\right )^{2}}
 {8} -\frac{\left (y - 3\right )^{2}} 
 {1} = 1
\]}
{\(2\)}{\(4\)}{\(6\)}{\(8\)}{}{}}
\End

\Data\ID{9000105504}
\Updated{2021-10-11 11:10:33}
\Level{B}
\SubArea{Conic Sections}
\LANGen{
\Question{
Find the distance between the points where the
\(y\)-axis
intersects the following hyperbola. 
\[ 
 H\colon \frac{\left (x - 4\right )^{2}}
 {10} -\frac{\left (y - 5\right )^{2}} 
 {15} = 1
\]}
{\(6\)}{\(2\)}{\(4\)}{\(8\)}{}{}}
\LANGpl{
\Question{
Wskaż  odległość punktów  przecięcia hiperboli $H$ z osią \(y\).
\[ 
 H\colon \frac{\left (x - 4\right )^{2}}
 {10} -\frac{\left (y - 5\right )^{2}} 
 {15} = 1
\]}
{\(6\)}{\(2\)}{\(4\)}{\(8\)}{}{}}
\LANGcs{
\Question{
Je dána hyperbola \(H\colon \frac{\left (x-4\right )^{2}} 
 {10} -\frac{\left (y-5\right )^{2}} 
 {15} = 1\).
Vzdálenost průsečíků této hyperboly s osou
\(y\) je
rovna:}
{\(6\)}{\(2\)}{\(4\)}{\(8\)}{}{}}
\LANGsk{
\Question{
Je daná hyperbola \(H\colon \frac{\left (x-4\right )^{2}} 
 {10} -\frac{\left (y-5\right )^{2}} 
 {15} = 1\).
Vzdialenosť priesečníkov tejto hyperboly s osou
\(y\) je
rovná:}
{\(6\)}{\(2\)}{\(4\)}{\(8\)}{}{}}
\LANGes{
\Question{
Encuentra la distancia entre los puntos en los que el eje \(y\)
corta a la siguiente hipérbola.  
\[ 
 H\colon \frac{\left (x - 4\right )^{2}}
 {10} -\frac{\left (y - 5\right )^{2}} 
 {15} = 1
\]}
{\(6\)}{\(2\)}{\(4\)}{\(8\)}{}{}}
\End

\Data\ID{9000105505}
\Updated{2020-11-07 07:11:25}
\Level{B}
\SubArea{Conic Sections}
\LANGen{
\Question{
Find the distance between the vertices of the following hyperbola. 
\[ 
 H\colon \frac{\left (x + 1\right )^{2}}
 {25} -\frac{\left (y + 2\right )^{2}} 
 {16} = 1
\]}
{\(10\)}{\(12\)}{\(6\)}{\(8\)}{}{}}
\LANGpl{
\Question{
Wskaż odległość pomiędzy wierzchołkami podanej hiperboli.
\[ 
 H\colon \frac{\left (x + 1\right )^{2}}
 {25} -\frac{\left (y + 2\right )^{2}} 
 {16} = 1
\]}
{\(10\)}{\(12\)}{\(6\)}{\(8\)}{}{}}
\LANGcs{
\Question{
Je dána hyperbola \(H\colon \frac{\left (x+1\right )^{2}} 
 {25} -\frac{\left (y+2\right )^{2}} 
 {16} = 1\).
Vzdálenost hlavních vrcholů této hyperboly je rovna:}
{\(10\)}{\(12\)}{\(6\)}{\(8\)}{}{}}
\LANGsk{
\Question{
Je daná hyperbola \(H\colon \frac{\left (x+1\right )^{2}} 
 {25} -\frac{\left (y+2\right )^{2}} 
 {16} = 1\).
Vzdialenosť hlavných vrcholov tejto hyperboly je rovná:}
{\(10\)}{\(12\)}{\(6\)}{\(8\)}{}{}}
\LANGes{
\Question{
Encuentra la distancia entre los vértices de la siguiente hipérbola. 
\[ 
 H\colon \frac{\left (x + 1\right )^{2}}
 {25} -\frac{\left (y + 2\right )^{2}} 
 {16} = 1
\]}
{\(10\)}{\(12\)}{\(6\)}{\(8\)}{}{}}
\End

\Data\ID{9000105506}
\Updated{2020-11-07 08:11:28}
\Level{B}
\SubArea{Conic Sections}
\LANGen{
\Question{
Find the distance between the vertices of the following hyperbola. 
\[ 
 H\colon \frac{\left (x - 3\right )^{2}}
 {16} -\frac{\left (y + 2\right )^{2}} 
 {25} = 1
\]}
{\(8\)}{\(12\)}{\(14\)}{\(10\)}{}{}}
\LANGpl{
\Question{
Wskaż odległość pomiędzy wierzchołkami podanej hiperboli.
\[ 
 H\colon \frac{\left (x - 3\right )^{2}}
 {16} -\frac{\left (y + 2\right )^{2}} 
 {25} = 1
\]}
{\(8\)}{\(12\)}{\(14\)}{\(10\)}{}{}}
\LANGcs{
\Question{
Je dána hyperbola \(H\colon \frac{\left (x-3\right )^{2}} 
 {16} -\frac{\left (y+2\right )^{2}} 
 {25} = 1\).
Vzdálenost hlavních vrcholů této hyperboly je rovna:}
{\(8\)}{\(12\)}{\(14\)}{\(10\)}{}{}}
\LANGsk{
\Question{
Je daná hyperbola \(H\colon \frac{\left (x-3\right )^{2}} 
 {16} -\frac{\left (y+2\right )^{2}} 
 {25} = 1\).
Vzdialenosť hlavných vrcholov tejto hyperboly je rovná:}
{\(8\)}{\(12\)}{\(14\)}{\(10\)}{}{}}
\LANGes{
\Question{
Encuentra la distancia entre los vértices de la siguiente hipérbola.
\[ 
 H\colon \frac{\left (x - 3\right )^{2}}
 {16} -\frac{\left (y + 2\right )^{2}} 
 {25} = 1
\]}
{\(8\)}{\(12\)}{\(14\)}{\(10\)}{}{}}
\End

\Data\ID{9000105507}
\Updated{2020-11-07 08:11:19}
\Level{B}
\SubArea{Conic Sections}
\LANGen{
\Question{
Find the distance between the foci of the following hyperbola. 
\[ 
 H\colon \frac{\left (x + 1\right )^{2}}
 {16} -\frac{\left (y + 5\right )^{2}} 
 {9} = 1
\]}
{\(10\)}{\(12\)}{\(8\)}{\(6\)}{}{}}
\LANGpl{
\Question{
Wskaż odległość pomiędzy ogniskami podanej hiperboli.
\[ 
 H\colon \frac{\left (x + 1\right )^{2}}
 {16} -\frac{\left (y + 5\right )^{2}} 
 {9} = 1
\]}
{\(10\)}{\(12\)}{\(8\)}{\(6\)}{}{}}
\LANGcs{
\Question{
Je dána hyperbola \(H\colon \frac{\left (x+1\right )^{2}} 
 {16} -\frac{\left (y+5\right )^{2}} 
 {9} = 1\).
Vzdálenost ohnisek této hyperboly je rovna:}
{\(10\)}{\(12\)}{\(8\)}{\(6\)}{}{}}
\LANGsk{
\Question{
Je daná hyperbola \(H\colon \frac{\left (x+1\right )^{2}} 
 {16} -\frac{\left (y+5\right )^{2}} 
 {9} = 1\).
Vzdialenosť ohnísk tejto hyperboly je rovná:}
{\(10\)}{\(12\)}{\(8\)}{\(6\)}{}{}}
\LANGes{
\Question{
Encuentra la distancia entre los focos de la siguiente hipérbola.
\[ 
 H\colon \frac{\left (x + 1\right )^{2}}
 {16} -\frac{\left (y + 5\right )^{2}} 
 {9} = 1
\]}
{\(10\)}{\(12\)}{\(8\)}{\(6\)}{}{}}
\End

\Data\ID{9000105508}
\Updated{2020-11-07 08:11:24}
\Level{B}
\SubArea{Conic Sections}
\LANGen{
\Question{
Find the distance between the foci of the following hyperbola. 
\[ 
 H\colon \frac{\left (x + 3\right )^{2}}
 {9} -\frac{\left (y - 2\right )^{2}} 
 {27} = 1
\]}
{\(12\)}{\(14\)}{\(10\)}{\(8\)}{}{}}
\LANGpl{
\Question{
Wskaż odległość pomiędzy ogniskami podanej hiperboli.
\[ 
 H\colon \frac{\left (x + 3\right )^{2}}
 {9} -\frac{\left (y - 2\right )^{2}} 
 {27} = 1
\]}
{\(12\)}{\(14\)}{\(10\)}{\(8\)}{}{}}
\LANGcs{
\Question{
Je dána hyperbola \(H\colon \frac{\left (x+3\right )^{2}} 
 {9} -\frac{\left (y-2\right )^{2}} 
 {27} = 1\).
Vzdálenost ohnisek této hyperboly je rovna:}
{\(12\)}{\(14\)}{\(10\)}{\(8\)}{}{}}
\LANGsk{
\Question{
Je daná hyperbola \(H\colon \frac{\left (x+3\right )^{2}} 
 {9} -\frac{\left (y-2\right )^{2}} 
 {27} = 1\).
Vzdialenosť ohnísk tejto hyperboly je rovná:}
{\(12\)}{\(14\)}{\(10\)}{\(8\)}{}{}}
\LANGes{
\Question{
Encuentra la distancia entre los focos de la siguiente hipérbola.
\[ 
 H\colon \frac{\left (x + 3\right )^{2}}
 {9} -\frac{\left (y - 2\right )^{2}} 
 {27} = 1
\]}
{\(12\)}{\(14\)}{\(10\)}{\(8\)}{}{}}
\End

\Data\ID{9000105509}
\Updated{2022-04-23 05:04:20}
\Level{B}
\SubArea{Conic Sections}
\LANGen{
\Question{
Find the distance between the points of intersection of the given hyperbola with the given straight line.
\[ 
 H\colon \frac{\left (x - 6\right )^{2}}
 {10} -\frac{\left (y - 2\right )^{2}} 
 {6} = 1;\quad p\colon x - 11 = 0
\]}
{\(6\)}{\(12\)}{\(10\)}{\(8\)}{}{}}
\LANGpl{
\Question{
Wskaż odległość pomiędzy punktami  przecięcia hiperboli i prostej.
\[ 
 H\colon \frac{\left (x - 6\right )^{2}}
 {10} -\frac{\left (y - 2\right )^{2}} 
 {6} = 1;\quad p\colon x - 11 = 0
\]}
{\(6\)}{\(12\)}{\(10\)}{\(8\)}{}{}}
\LANGcs{
\Question{
Je dána hyperbola \(H\colon \frac{\left (x-6\right )^{2}} 
 {10} -\frac{\left (y-2\right )^{2}} 
 {6} = 1\) 
a přímka \(p\colon x - 11 = 0\).
Vzdálenost průsečíků této hyperboly s přímkou
\(p\) je
rovna:}
{\(6\)}{\(12\)}{\(10\)}{\(8\)}{}{}}
\LANGsk{
\Question{
Je daná hyperbola \(H\colon \frac{\left (x-6\right )^{2}} 
 {10} -\frac{\left (y-2\right )^{2}} 
 {6} = 1\) 
a priamka \(p\colon x - 11 = 0\).
Vzdialenosť priesečníkov tejto hyperboly s priamkou
\(p\) je
rovná:}
{\(6\)}{\(12\)}{\(10\)}{\(8\)}{}{}}
\LANGes{
\Question{
Encuentra la distancia entre las intresecciones de la siguiente hipérbola y la recta. 
\[ 
 H\colon \frac{\left (x - 6\right )^{2}}
 {10} -\frac{\left (y - 2\right )^{2}} 
 {6} = 1;\quad p\colon x - 11 = 0
\]}
{\(6\)}{\(12\)}{\(10\)}{\(8\)}{}{}}
\End

\Data\ID{9000105510}
\Updated{2020-11-07 08:11:24}
\Level{B}
\SubArea{Conic Sections}
\LANGen{
\Question{
Find the distance between the intersections of the following hyperbola and straight
line. 
\[ 
 H\colon \frac{\left (x - 2\right )^{2}}
 {10} -\frac{\left (y + 2\right )^{2}} 
 {6} = 1;\quad p\colon y + 5 = 0
\]}
{\(10\)}{\(12\)}{\(6\)}{\(8\)}{}{}}
\LANGpl{
\Question{
Odległość od przecięcia hiperboli do   prostej jest równa:
\[ 
 H\colon \frac{\left (x - 2\right )^{2}}
 {10} -\frac{\left (y + 2\right )^{2}} 
 {6} = 1;\quad p\colon y + 5 = 0
\]}
{\(10\)}{\(12\)}{\(6\)}{\(8\)}{}{}}
\LANGcs{
\Question{
Je dána hyperbola \(H\colon \frac{\left (x-2\right )^{2}} 
 {10} -\frac{\left (y+2\right )^{2}} 
 {6} = 1\) 
a přímka \(p\colon y + 5 = 0\).
Vzdálenost průsečíků této hyperboly s přímkou
\(p\) je
rovna:}
{\(10\)}{\(12\)}{\(6\)}{\(8\)}{}{}}
\LANGsk{
\Question{
Je daná hyperbola \(H\colon \frac{\left (x-2\right )^{2}} 
 {10} -\frac{\left (y+2\right )^{2}} 
 {6} = 1\) 
a priamka \(p\colon y + 5 = 0\).
Vzdialenosť priesečníkov tejto hyperboly s priamkou
\(p\) je
rovná:}
{\(10\)}{\(12\)}{\(6\)}{\(8\)}{}{}}
\LANGes{
\Question{
Encuentra la distancia entre las intresecciones de la siguiente hipérbola y la línea recta. 
\[ 
 H\colon \frac{\left (x - 2\right )^{2}}
 {10} -\frac{\left (y + 2\right )^{2}} 
 {6} = 1;\quad p\colon y + 5 = 0
\]}
{\(10\)}{\(12\)}{\(6\)}{\(8\)}{}{}}
\End

\Data\ID{9000120004}
\Updated{2021-10-28 09:10:23}
\Level{B}
\SubArea{Conic Sections}
\LANGen{
\Question{
Find a conic section which allows to draw a line through the center of this conic and
this line does not have any common point with the conic.}
{hyperbola}{circle}{ellipse}{parabola}{no solution}{}}
\LANGpl{
\Question{
Wskaż przekrój stożka tak, aby móc narysować  prostą przechodzącą przez środek  krzywej stożkowej  i tak,  aby ta  prosta nie miała żadnych punktów wspólnych z krzywą stożkową.}
{hiperbola}{okrąg}{elipsa}{parabola}{brak rozwiązania}{}}
\LANGcs{
\Question{
Vyberte takovou kuželosečku, aby bylo možno vést jejím středem
přímku, která s touto kuželosečkou nemá společný žádný bod.}
{Hyperbola}{Kružnice}{Elipsa}{Parabola}{Žádná z uvedených kuželoseček tuto vlastnost nemá.}{}}
\LANGsk{
\Question{
Vyberte takú kužeľosečku, aby bolo možné viesť jej stredom
priamku, ktorá s touto kužeľosečkou nemá spoločný žiadny bod.}
{Hyperbola}{Kružnica}{Elipsa}{Parabola}{Žiadna z uvedených kužeľosečiek túto vlastnosť nemá.}{}}
\LANGes{
\Question{
Halla una sección cónica por cuyo centro pase una recta y esta recta no tenga ningún punto en común con dicha sección cónica.}
{hipérbola}{circunferencia}{elipse}{parábola}{Esta sección cónica no existe.}{}}
\End

\Data\ID{9000120005}
\Updated{2021-10-28 10:10:49}
\Level{B}
\SubArea{Conic Sections}
\LANGen{
\Question{
The executives of a camp organize a holiday game. For this game it is important
that the direct distance kitchen - tent - fireplace is equal for all tents in the camp.
Is this information enough to determine the curve passing through all the
tents in the camp? Is this curve a conic? If yes, determine which conic.}
{Yes, all the tents are on an ellipse.}{Yes, all the tents are on a circle.}{Yes, all the tents are on a parabola.}{Yes, all the tents are on a hyperbola.}{No, we do not have enough information to draw any conclusion.}{}}
\LANGpl{
\Question{
Kierownik obozu zorganizował wakacyjną grę. Głównym  założeniem gry jest to, aby 
odległość:  kuchnia - namiot  - kominek była taka sama dla wszystkich namiotów w obozie.
Czy podana informacja jest wystarczająca, aby określić krzywą przechodzącą przez wszystkie namioty w obozie? Czy ta krzywa jest krzywą stożkową? Jeśli tak, to wskaż którą?}
{Tak, wszystkie namioty leżą na elipsie.}{Tak, wszystkie namioty leżą na okręgu.}{Tak, wszystkie namioty leżą na paraboli.}{Tak, wszystkie namioty leżą na hiperboli.}{Brak rozwiązania.}{}}
\LANGcs{
\Question{
Pro organizaci celotáborové hry je důležité, aby vzdušná vzdálenost
kuchyň - stan - ohniště byla pro všechny stany stejná. Lze z této informace
odvodit, že stany leží na jedné kuželosečce? Pokud ano, tak na které?}
{Ano, stany leží na elipse.}{Ano, stany leží na kružnici.}{Ano, stany leží na parabole.}{Ano, stany leží na hyperbole.}{Ne, z této informace nelze odvodit, že by stany ležely na nějaké konkrétní
kuželosečce.}{}}
\LANGsk{
\Question{
Pre organizáciu celotáborovej hry je dôležité, aby vzdušná vzdialenosť
kuchyňa - stan - ohnisko bola pre všetky stany rovnaká. Dá sa z tejto informácie
odvodiť, že stany ležia na jednej kužeľosečke? Pokiaľ áno, tak na ktorej?}
{Áno, stany ležia na elipse.}{Áno, stany ležia na kružnici.}{Áno, stany ležia na parabole.}{Áno, stany ležia na hyperbole.}{Ne, z tejto informácie nemožno odvodiť, že by stany ležali na nejakej konkrétnej
kužeľosečke.}{}}
\LANGes{
\Question{
Los organizadores de un campamento prepararon un juego. Para este juego es importante que la distancia directa entre la cocina, la tienda y la hoguera sea para todas las tiendas de campaña igual. ¿Basta esta información para determinar la curva que pasa por todas estas tiendas? ¿Es esta curva una sección cónica? Si es así en qué sección cónica están las tiendas?}
{Sí, todas la tiendas están en una elipse.}{Sí, todas las tiendas están en una circunferencia.}{Sí, todas las tiendas están en una parábola.}{Sí, todas la tiendas están en una hipérbola.}{No, no tenemos suficiente información para poder sacar conclusiones de esta sección cónica.}{}}
\End

\Data\ID{9000120007}
\Updated{2021-10-28 09:10:40}
\Level{B}
\SubArea{Conic Sections}
\LANGen{
\Question{
On a map of a city, the town hall is represented by a point and a river through the
city by a straight line. There are places in the city with the property that the direct
distance from each place to the town hall is equal to the direct distance to the river.
In the following list identify a curve which can be used to join all these places.}
{parabola}{circle}{ellipse}{hyperbola}{none of them}{}}
\LANGpl{
\Question{
Na mapie miasta ratusz został oznaczony punktem, natomiast rzeka prostą. W mieście znajdują się obiekty, których odległość do ratusza jest równa odległości do rzeki. Wskaż krzywą, która mogłaby połączyć wszystkie te obiekty.}
{parabola}{okrąg}{elipsa}{hiperbola}{Brak rozwiązania}{}}
\LANGcs{
\Question{
Ve městě je několik míst, která jsou stejně vzdálená od řeky i od
městské radnice. Vyberte křivku, kterou se dají všechna tato místa na
mapě propojit za předpokladu, že tok řeky je na mapě města a jeho okolí
zobrazen přímkou a radnice bodem.}
{Parabola}{Kružnice}{Elipsa}{Hyperbola}{Žádnou z výše uvedených kuželoseček nelze zmíněná místa propojit.}{}}
\LANGsk{
\Question{
V meste je niekoľko miest, ktoré sú rovnako vzdialené od rieky i od
mestskej radnice. Vyberte krivku, ktorou sa dajú všetky tieto miesta na
mape prepojiť za predpokladu, že tok rieky je na mape mesta a jeho okolia
zobrazený priamkou a radnica bodom.}
{Parabola}{Kružnica}{Elipsa}{Hyperbola}{Žiadnou z vyššie uvedených kužeľosečiek nemožno spomínané miesta prepojiť.}{}}
\LANGes{
\Question{
En un mapa de una ciudad el ayuntamiento está representado como un punto y el río que pasa por la ciudad como una recta. En la ciudad hay lugares a la misma distancia tanto del río como del ayuntamiento.  De la lista siguiente define la cónica que pueda conectar todos estos lugares.}
{parábola}{circunferencia}{elipse}{hipérbola}{ninguna de la lista}{}}
\End

\Data\ID{9000149706}
\Updated{2020-08-30 08:08:20}
\Level{B}
\SubArea{Conic Sections}
\LANGen{
\Question{
Find the center of the following hyperbola. 
\[ 
 4x^{2} - 3y^{2} + 8x - 30y - 49 = 0
\]}
{\([-1;-5]\)}{\([-1;5]\)}{\([1;-5]\)}{\([1;5]\)}{}{}}
\LANGpl{
\Question{
Wyznacz środek hiperboli.
\[ 
 4x^{2} - 3y^{2} + 8x - 30y - 49 = 0
\]}
{\([-1;-5]\)}{\([-1;5]\)}{\([1;-5]\)}{\([1;5]\)}{}{}}
\LANGcs{
\Question{
Hyperbola je dána rovnicí \(4x^{2} - 3y^{2} + 8x - 30y - 49 = 0\).
Její střed má souřadnice:}
{\([-1;-5]\)}{\([-1;5]\)}{\([1;-5]\)}{\([1;5]\)}{}{}}
\LANGsk{
\Question{
Hyperbola je daná rovnicou \(4x^{2} - 3y^{2} + 8x - 30y - 49 = 0\).
Jej stred má súradnice:}
{\([-1;-5]\)}{\([-1;5]\)}{\([1;-5]\)}{\([1;5]\)}{}{}}
\LANGes{
\Question{
Halla el centro de la siguiente hipérbola.  
\[ 
 4x^{2} - 3y^{2} + 8x - 30y - 49 = 0
\]}
{\([-1;-5]\)}{\([-1;5]\)}{\([1;-5]\)}{\([1;5]\)}{}{}}
\End

\Data\ID{9000149707}
\Updated{2020-08-30 08:08:54}
\Level{B}
\SubArea{Conic Sections}
\LANGen{
\Question{
Find the center of the following hyperbola. 
\[ 
 5x^{2} - 6y^{2} - 30x + 12y + 9 = 0
\]}
{\([3;1]\)}{\([3;-1]\)}{\([-3;1]\)}{\([-3;-1]\)}{}{}}
\LANGpl{
\Question{
Wyznacz środek hiperboli.
\[ 
 5x^{2} - 6y^{2} - 30x + 12y + 9 = 0
\]}
{\([3;1]\)}{\([3;-1]\)}{\([-3;1]\)}{\([-3;-1]\)}{}{}}
\LANGcs{
\Question{
Hyperbola je dána rovnicí \(5x^{2} - 6y^{2} - 30x + 12y + 9 = 0\).
Její střed má souřadnice:}
{\([3;1]\)}{\([3;-1]\)}{\([-3;1]\)}{\([-3;-1]\)}{}{}}
\LANGsk{
\Question{
Hyperbola je daná rovnicou \(5x^{2} - 6y^{2} - 30x + 12y + 9 = 0\).
Jej stred má súradnice:}
{\([3;1]\)}{\([3;-1]\)}{\([-3;1]\)}{\([-3;-1]\)}{}{}}
\LANGes{
\Question{
Halla el centro de la siguiente hipérbola.  
\[ 
 5x^{2} - 6y^{2} - 30x + 12y + 9 = 0
\]}
{\([3;1]\)}{\([3;-1]\)}{\([-3;1]\)}{\([-3;-1]\)}{}{}}
\End

\Data\ID{9000149708}
\Updated{2022-04-23 05:04:20}
\Level{B}
\SubArea{Conic Sections}
\LANGen{
\Question{
Find the vertex of the following parabola. 
\[ 
 x^{2} + 8x - 4y + 24 = 0
\]}
{\([-4;2]\)}{\([-4;-2]\)}{\([4;2]\)}{\([4;-2]\)}{}{}}
\LANGpl{
\Question{
Wskaż współrzędne wierzchołka paraboli.
\[ 
 x^{2} + 8x - 4y + 24 = 0
\]}
{\([-4;2]\)}{\([-4;-2]\)}{\([4;2]\)}{\([4;-2]\)}{}{}}
\LANGcs{
\Question{
Parabola je dána rovnicí \(x^{2} + 8x - 4y + 24 = 0\).
Její vrchol má souřadnice:}
{\([-4;2]\)}{\([-4;-2]\)}{\([4;2]\)}{\([4;-2]\)}{}{}}
\LANGsk{
\Question{
Parabola je daná rovnicou \(x^{2} + 8x - 4y + 24 = 0\).
Jej vrchol má súradnice:}
{\([-4;2]\)}{\([-4;-2]\)}{\([4;2]\)}{\([4;-2]\)}{}{}}
\LANGes{
\Question{
Halla el vértice de la siguiente parábola. 
\[ 
 x^{2} + 8x - 4y + 24 = 0
\]}
{\([-4;2]\)}{\([-4;-2]\)}{\([4;2]\)}{\([4;-2]\)}{}{}}
\End

\Data\ID{9000149709}
\Updated{2020-08-30 08:08:10}
\Level{B}
\SubArea{Conic Sections}
\LANGen{
\Question{
Find the vertex of the following parabola. 
\[ 
 y^{2} - 12x + 4y + 64 = 0
\]}
{\([5;-2]\)}{\([5;2]\)}{\([-5;2]\)}{\([-5;-2]\)}{}{}}
\LANGpl{
\Question{
Wskaż współrzędne wierzchołka paraboli.
\[ 
 y^{2} - 12x + 4y + 64 = 0
\]}
{\([5;-2]\)}{\([5;2]\)}{\([-5;2]\)}{\([-5;-2]\)}{}{}}
\LANGcs{
\Question{
Parabola je dána rovnicí \(y^{2} - 12x + 4y + 64 = 0\).
Její vrchol má souřadnice:}
{\([5;-2]\)}{\([5;2]\)}{\([-5;2]\)}{\([-5;-2]\)}{}{}}
\LANGsk{
\Question{
Parabola je daná rovnicou \(y^{2} - 12x + 4y + 64 = 0\).
Jej vrchol má súradnice:}
{\([5;-2]\)}{\([5;2]\)}{\([-5;2]\)}{\([-5;-2]\)}{}{}}
\LANGes{
\Question{
Halla el vértice de la siguiente parábola.
\[ 
 y^{2} - 12x + 4y + 64 = 0
\]}
{\([5;-2]\)}{\([5;2]\)}{\([-5;2]\)}{\([-5;-2]\)}{}{}}
\End

\Data\ID{9000149710}
\Updated{2020-08-30 08:08:43}
\Level{B}
\SubArea{Conic Sections}
\LANGen{
\Question{
Find the vertex of the following parabola. 
\[ 
 x^{2} - 6x - 12y - 3 = 0
\]}
{\([3;-1]\)}{\([3;1]\)}{\([-3;1]\)}{\([-3;-1]\)}{}{}}
\LANGpl{
\Question{
Wskaż współrzędne wierzchołka paraboli.
\[ 
 x^{2} - 6x - 12y - 3 = 0
\]}
{\([3;-1]\)}{\([3;1]\)}{\([-3;1]\)}{\([-3;-1]\)}{}{}}
\LANGcs{
\Question{
Parabola je dána rovnicí \(x^{2} - 6x - 12y - 3 = 0\).
Její vrchol má souřadnice:}
{\([3;-1]\)}{\([3;1]\)}{\([-3;1]\)}{\([-3;-1]\)}{}{}}
\LANGsk{
\Question{
Parabola je daná rovnicou \(x^{2} - 6x - 12y - 3 = 0\).
Jej vrchol má súradnice:}
{\([3;-1]\)}{\([3;1]\)}{\([-3;1]\)}{\([-3;-1]\)}{}{}}
\LANGes{
\Question{
Halla el vértice de la siguiente parábola.
\[ 
 x^{2} - 6x - 12y - 3 = 0
\]}
{\([3;-1]\)}{\([3;1]\)}{\([-3;1]\)}{\([-3;-1]\)}{}{}}
\End

\Data\ID{9000168706}
\Updated{2022-04-23 05:04:20}
\Level{B}
\SubArea{Conic Sections}
\LANGen{
\Question{
Which of the given points is one of the vertices of the hyperbola  \(16x^{2} - 9y^{2} + 64x + 18y - 89 = 0\)?}
{\([-5;1]\)}{\([-6;1]\)}{\([-2;4]\)}{\([-2;5]\)}{}{}}
\LANGpl{
\Question{
Wskaż współrzędne jednego z wierzchołków hiperboli  \(16x^{2} - 9y^{2} + 64x + 18y - 89 = 0\).}
{\([-5;1]\)}{\([-6;1]\)}{\([-2;4]\)}{\([-2;5]\)}{}{}}
\LANGcs{
\Question{
Hyperbola je dána rovnicí \(16x^{2} - 9y^{2} + 64x + 18y - 89 = 0\).
Jeden z jejich vrcholů vrchol má souřadnice:}
{\([-5;1]\)}{\([-6;1]\)}{\([-2;4]\)}{\([-2;5]\)}{}{}}
\LANGsk{
\Question{
Hyperbola je daná rovnicou \(16x^{2} - 9y^{2} + 64x + 18y - 89 = 0\).
Jeden z ich vrcholov má súradnice:}
{\([-5;1]\)}{\([-6;1]\)}{\([-2;4]\)}{\([-2;5]\)}{}{}}
\LANGes{
\Question{
Halla el vértice de la hipérbola \(16x^{2} - 9y^{2} + 64x + 18y - 89 = 0\).}
{\([-5;1]\)}{\([-6;1]\)}{\([-2;4]\)}{\([-2;5]\)}{}{}}
\End

\Data\ID{9000168707}
\Updated{2020-11-07 04:11:16}
\Level{B}
\SubArea{Conic Sections}
\LANGen{
\Question{
Find the vertex of the hyperbola \(9x^{2} - 4y^{2} - 54x - 16y + 29 = 0\).}
{\([5;-2]\)}{\([6;-2]\)}{\([3;-4]\)}{\([3;-5]\)}{}{}}
\LANGpl{
\Question{
Wskaż współrzędne wierzchołka hiperboli \(9x^{2} - 4y^{2} - 54x - 16y + 29 = 0\).}
{\([5;-2]\)}{\([6;-2]\)}{\([3;-4]\)}{\([3;-5]\)}{}{}}
\LANGcs{
\Question{
Hyperbola je dána rovnicí \(9x^{2} - 4y^{2} - 54x - 16y + 29 = 0\).
Jeden z jejích vrcholů má souřadnice:}
{\([5;-2]\)}{\([6;-2]\)}{\([3;-4]\)}{\([3;-5]\)}{}{}}
\LANGsk{
\Question{
Hyperbola je daná rovnicou \(9x^{2} - 4y^{2} - 54x - 16y + 29 = 0\).
Jej vrchol má súradnice:}
{\([5;-2]\)}{\([6;-2]\)}{\([3;-4]\)}{\([3;-5]\)}{}{}}
\LANGes{
\Question{
Halla el vértice de la hipérbola \(9x^{2} - 4y^{2} - 54x - 16y + 29 = 0\).}
{\([5;-2]\)}{\([6;-2]\)}{\([3;-4]\)}{\([3;-5]\)}{}{}}
\End

\Data\ID{9000168708}
\Updated{2020-11-07 04:11:47}
\Level{B}
\SubArea{Conic Sections}
\LANGen{
\Question{
Find the vertex of the hyperbola \(4x^{2} - 9y^{2} + 24x + 18y + 63 = 0\).}
{\([-3;3]\)}{\([-3;4]\)}{\([-6;1]\)}{\([-5;1]\)}{}{}}
\LANGpl{
\Question{
Wskaż współrzędne wierzchołka hiperboli \(4x^{2} - 9y^{2} + 24x + 18y + 63 = 0\).}
{\([-3;3]\)}{\([-3;4]\)}{\([-6;1]\)}{\([-5;1]\)}{}{}}
\LANGcs{
\Question{
Hyperbola je dána rovnicí \(4x^{2} - 9y^{2} + 24x + 18y + 63 = 0\).
Její vrchol má souřadnice:}
{\([-3;3]\)}{\([-3;4]\)}{\([-6;1]\)}{\([-5;1]\)}{}{}}
\LANGsk{
\Question{
Hyperbola je daná rovnicou \(4x^{2} - 9y^{2} + 24x + 18y + 63 = 0\).
Jej vrchol má súradnice:}
{\([-3;3]\)}{\([-3;4]\)}{\([-6;1]\)}{\([-5;1]\)}{}{}}
\LANGes{
\Question{
Halla el vértice de la hipérbola \(4x^{2} - 9y^{2} + 24x + 18y + 63 = 0\).}
{\([-3;3]\)}{\([-3;4]\)}{\([-6;1]\)}{\([-5;1]\)}{}{}}
\End

\Data\ID{9000168709}
\Updated{2020-08-30 07:08:52}
\Level{B}
\SubArea{Conic Sections}
\LANGen{
\Question{
Find the vertex of the hyperbola \(9x^{2} - 16y^{2} - 36x - 96y + 36 = 0\).}
{\([2;-6]\)}{\([2;-7]\)}{\([5;-3]\)}{\([6;-3]\)}{}{}}
\LANGpl{
\Question{
Wskaż współrzędne wierzchołka hiperboli \(9x^{2} - 16y^{2} - 36x - 96y + 36 = 0\).}
{\([2;-6]\)}{\([2;-7]\)}{\([5;-3]\)}{\([6;-3]\)}{}{}}
\LANGcs{
\Question{
Hyperbola je dána rovnicí \(9x^{2} - 16y^{2} - 36x - 96y + 36 = 0\).
Její vrchol má souřadnice:}
{\([2;-6]\)}{\([2;-7]\)}{\([5;-3]\)}{\([6;-3]\)}{}{}}
\LANGsk{
\Question{
Hyperbola je daná rovnicou \(9x^{2} - 16y^{2} - 36x - 96y + 36 = 0\).
Jej vrchol má súradnice:}
{\([2;-6]\)}{\([2;-7]\)}{\([5;-3]\)}{\([6;-3]\)}{}{}}
\LANGes{
\Question{
Halla el vértice de la hipérbola \(9x^{2} - 16y^{2} - 36x - 96y + 36 = 0\).}
{\([2;-6]\)}{\([2;-7]\)}{\([5;-3]\)}{\([6;-3]\)}{}{}}
\End

\Data\ID{9000168710}
\Updated{2020-08-24 09:08:11}
\Level{B}
\SubArea{Conic Sections}
\LANGen{
\Question{
Find the vertex of the hyperbola \(9x^{2} - 25y^{2} + 18x + 100y - 316 = 0\).}
{\([4;2]\)}{\([2;2]\)}{\([-1;-3]\)}{\([-1;-1]\)}{}{}}
\LANGpl{
\Question{
Wskaż współrzędne wierzchołka hiperboli \(9x^{2} - 25y^{2} + 18x + 100y - 316 = 0\).}
{\([4;2]\)}{\([2;2]\)}{\([-1;-3]\)}{\([-1;-1]\)}{}{}}
\LANGcs{
\Question{
Hyperbola je dána rovnicí \(9x^{2} - 25y^{2} + 18x + 100y - 316 = 0\).
Její vrchol má souřadnice:}
{\([4;2]\)}{\([2;2]\)}{\([-1;-3]\)}{\([-1;-1]\)}{}{}}
\LANGsk{
\Question{
Hyperbola je daná rovnicou \(9x^{2} - 25y^{2} + 18x + 100y - 316 = 0\).
Jej vrchol má súradnice:}
{\([4;2]\)}{\([2;2]\)}{\([-1;-3]\)}{\([-1;-1]\)}{}{}}
\LANGes{
\Question{
Dada la ecuación de la hipérbola. \(9x^{2} - 25y^{2} + 18x + 100y - 316 = 0\). Halla las coordenadas del vértice de la hipérbola.}
{\([4;2]\)}{\([2;2]\)}{\([-1;-3]\)}{\([-1;-1]\)}{}{}}
\End

\Data\ID{1003024103}
\Updated{2022-08-25 10:08:40}
\Level{C}
\SubArea{Conic Sections}
\LANGen{
\Question{
Find the equation of the  circle, which passes through the points \( A=[-3;2] \), \( B=[-1;4] \) and \( C=[3;0] \).}
{\( x^2 + (y-1)^2 = 10 \)}{\( (x-1)^2 + y^2 = 10 \)}{\( x^2 + (y-1)^2 = \sqrt{10} \)}{\( x^2 + y^2 = 10 \)}{\( x^2 + (y+1)^2 = 10 \)}{}}
\LANGpl{
\Question{
Wskaż równanie okręgu przechodzącego przez punkty \( A=[-3;2] \), \( B=[-1;4] \) i \( C=[3;0] \).}
{\( x^2 + (y-1)^2 = 10 \)}{\( (x-1)^2 + y^2 = 10 \)}{\( x^2 + (y-1)^2 = \sqrt{10} \)}{\( x^2 + y^2 = 10 \)}{\( x^2 + (y+1)^2 = 10 \)}{}}
\LANGcs{
\Question{
Rovnice kružnice, která prochází body \( A=[-3;2] \), \( B=[-1;4] \) a \( C=[3;0] \), je:}
{\( x^2 + (y-1)^2 = 10 \)}{\( (x-1)^2 + y^2 = 10 \)}{\( x^2 + (y-1)^2 = \sqrt{10} \)}{\( x^2 + y^2 = 10 \)}{\( x^2 + (y+1)^2 = 10 \)}{}}
\LANGsk{
\Question{
Rovnica kružnice, ktorá prechádza bodmi \( A=[-3;2] \), \( B=[-1;4] \) a \( C=[3;0] \), je:}
{\( x^2 + (y-1)^2 = 10 \)}{\( (x-1)^2 + y^2 = 10 \)}{\( x^2 + (y-1)^2 = \sqrt{10} \)}{\( x^2 + y^2 = 10 \)}{\( x^2 + (y+1)^2 = 10 \)}{}}
\LANGes{
\Question{
Encuentra la ecuación de la circunferencia que pasa por los puntos \( A=(-3;2) \), \( B=(-1;4) \) y \( C=(3;0) \).}
{\( x^2 + (y-1)^2 = 10 \)}{\( (x-1)^2 + y^2 = 10 \)}{\( x^2 + (y-1)^2 = \sqrt{10} \)}{\( x^2 + y^2 = 10 \)}{\( x^2 + (y+1)^2 = 10 \)}{}}
\End

\Data\ID{1003024104}
\Updated{2022-08-25 10:08:40}
\Level{C}
\SubArea{Conic Sections}
\LANGen{
\Question{
Find the points of intersection of the circle \( x^2+y^2=4 \) and the line \( x+y-2=0 \).}
{\( [0;2] \), \( [2;0] \)}{\( [0;-2] \), \(  [-2;0] \)}{\( [0;-2] \), \(  [2;0] \)}{\( [0;2] \), \(  [-2;0] \)}{\( [0;-2] \), \(  [0;2] \)}{}}
\LANGpl{
\Question{
Wskaż punkty przecięcia okręgu \( x^2+y^2=4 \) i prostej \( x+y-2=0 \).}
{\( [0;2] \), \( [2;0] \)}{\( [0;-2] \), \(  [-2;0] \)}{\( [0;-2] \), \(  [2;0] \)}{\( [0;2] \), \(  [-2;0] \)}{\( [0;-2] \), \(  [0;2] \)}{}}
\LANGcs{
\Question{
Společné body kružnice \( x^2+y^2=4 \) a přímky \( x+y-2=0 \) jsou:}
{\( [0;2] \), \( [2;0] \)}{\( [0;-2] \), \(  [-2;0] \)}{\( [0;-2] \), \(  [2;0] \)}{\( [0;2] \), \(  [-2;0] \)}{\( [0;-2] \), \(  [0;2] \)}{}}
\LANGsk{
\Question{
Spoločné body kružnice \( x^2+y^2=4 \) a priamky \( x+y-2=0 \) sú:}
{\( [0;2] \), \( [2;0] \)}{\( [0;-2] \), \(  [-2;0] \)}{\( [0;-2] \), \(  [2;0] \)}{\( [0;2] \), \(  [-2;0] \)}{\( [0;-2] \), \(  [0;2] \)}{}}
\LANGes{
\Question{
Halla los puntos de la intersección de la circunferencia \( x^2+y^2=4 \) y la recta \( x+y-2=0 \).}
{\( [0;2] \), \( [2;0] \)}{\( [0;-2] \), \(  [-2;0] \)}{\( [0;-2] \), \(  [2;0] \)}{\( [0;2] \), \(  [-2;0] \)}{\( [0;-2] \), \(  [0;2] \)}{}}
\End

\Data\ID{1003024105}
\Updated{2022-08-25 10:08:40}
\Level{C}
\SubArea{Conic Sections}
\LANGen{
\Question{
Find the value of the parameter \( q\in\mathbb{R} \) for which the line \( x+2y-1=0 \) is a tangent of the hyperbola  \( x^2-2y^2=q \).}
{\( -1 \)}{\( 1 \)}{\( 2 \)}{\( -2 \)}{\( \frac12 \)}{}}
\LANGpl{
\Question{
Wyznacz wartość parametru \( q\in\mathbb{R} \) tak, aby prosta \( x+2y-1=0 \) była styczną do hiperboli \( x^2-2y^2=q \).}
{\( -1 \)}{\( 1 \)}{\( 2 \)}{\( -2 \)}{\( \frac12 \)}{}}
\LANGcs{
\Question{
Pro kterou hodnotu parametru \( q\in\mathbb{R} \) je přímka \( x+2y-1=0 \) tečnou hyperboly  \( x^2-2y^2=q \) ?}
{\( -1 \)}{\( 1 \)}{\( 2 \)}{\( -2 \)}{\( \frac12 \)}{}}
\LANGsk{
\Question{
Pre ktorú hodnotu parametra \( q\in\mathbb{R} \) je priamka \( x+2y-1=0 \) dotyčnica hyperboly  \( x^2-2y^2=q \) ?}
{\( -1 \)}{\( 1 \)}{\( 2 \)}{\( -2 \)}{\( \frac12 \)}{}}
\LANGes{
\Question{
Encuentra el valor del parámetro \( q\in\mathbb{R} \) para el cual la recta \( x+2y-1=0 \) es tangente a la hipérbola \( x^2-2y^2=q \).}
{\( -1 \)}{\( 1 \)}{\( 2 \)}{\( -2 \)}{\( \frac12 \)}{}}
\End

\Data\ID{1003024106}
\Updated{2022-08-25 10:08:44}
\Level{C}
\SubArea{Conic Sections}
\LANGen{
\Question{
Find the value of the parameter \( q\in\mathbb{R} \) for which the line \( x+2y-1=0 \) is a tangent of the ellipse  \( x^2+4y^2=q \).}
{\( \frac12 \)}{\( \frac14 \)}{\( 2 \)}{\( 4 \)}{\( 1 \)}{}}
\LANGpl{
\Question{
Wyznacz wartość parametru \( q\in\mathbb{R} \) tak, aby prosta \( x+2y-1=0 \) była styczną do elipsy \( x^2+4y^2=q \).}
{\( \frac12 \)}{\( \frac14 \)}{\( 2 \)}{\( 4 \)}{\( 1 \)}{}}
\LANGcs{
\Question{
Pro kterou hodnotu parametru \( q\in\mathbb{R} \) je přímka \( x+2y-1=0 \) tečnou elipsy  \( x^2+4y^2=q \) ?}
{\( \frac12 \)}{\( \frac14 \)}{\( 2 \)}{\( 4 \)}{\( 1 \)}{}}
\LANGsk{
\Question{
Pre ktorú hodnotu parametra \( q\in\mathbb{R} \) je priamka \( x+2y-1=0 \) dotyčnicou elipsy  \( x^2+4y^2=q \) ?}
{\( \frac12 \)}{\( \frac14 \)}{\( 2 \)}{\( 4 \)}{\( 1 \)}{}}
\LANGes{
\Question{
Halla el valor del parámetro \( q\in\mathbb{R} \) para el cual la recta \( x+2y-1=0 \) es tangente a la elipse \( x^2+4y^2=q \).}
{\( \frac12 \)}{\( \frac14 \)}{\( 2 \)}{\( 4 \)}{\( 1 \)}{}}
\End

\Data\ID{1003024107}
\Updated{2022-08-25 10:08:44}
\Level{C}
\SubArea{Conic Sections}
\LANGen{
\Question{
Find the value of the parameter \( p\in\mathbb{R} \) for which the line \( x+2y-1=0 \) is a tangent of the  parabola \( y^2=2px \).}
{\( -\frac12 \)}{\( \frac12 \)}{\( 2 \)}{\( -2 \)}{\( 1 \)}{}}
\LANGpl{
\Question{
Wyznacz wartość parametru \( p\in\mathbb{R} \) tak, aby prosta \( x+2y-1=0 \) była styczną do  paraboli \( y^2=2px \).}
{\( -\frac12 \)}{\( \frac12 \)}{\( 2 \)}{\( -2 \)}{\( 1 \)}{}}
\LANGcs{
\Question{
Pro kterou hodnotu parametru \( p\in\mathbb{R} \) je přímka \( x+2y-1=0 \) tečnou paraboly \( y^2=2px \) ?}
{\( -\frac12 \)}{\( \frac12 \)}{\( 2 \)}{\( -2 \)}{\( 1 \)}{}}
\LANGsk{
\Question{
Pre ktorú hodnotu parametra \( p\in\mathbb{R} \) je priamka \( x+2y-1=0 \) dotyčnicou paraboly \( y^2=2px \) ?}
{\( -\frac12 \)}{\( \frac12 \)}{\( 2 \)}{\( -2 \)}{\( 1 \)}{}}
\LANGes{
\Question{
Halla el valor del parámetro \( p\in\mathbb{R} \) para el cual la recta \( x+2y-1=0 \) es tangente a la parábola \( y^2=2px \).}
{\( -\frac12 \)}{\( \frac12 \)}{\( 2 \)}{\( -2 \)}{\( 1 \)}{}}
\End

\Data\ID{1003024108}
\Updated{2022-08-25 10:08:44}
\Level{C}
\SubArea{Conic Sections}
\LANGen{
\Question{
Find the value of the parameter \( q\in\mathbb{R} \) for which the line \( x+2y-1=0 \) is a tangent of the circle \( x^2+y^2=q^2 \).}
{\( \frac{\sqrt5}5 \)}{\( 5 \)}{\( \sqrt5 \)}{\( \frac15 \)}{\( 1 \)}{}}
\LANGpl{
\Question{
Wyznacz wartość parametru \( q\in\mathbb{R} \) tak, aby prosta \( x+2y-1=0 \) była styczną do okręgu \( x^2+y^2=q^2 \).}
{\( \frac{\sqrt5}5 \)}{\( 5 \)}{\( \sqrt5 \)}{\( \frac15 \)}{\( 1 \)}{}}
\LANGcs{
\Question{
Pro kterou hodnotu parametru \( q\in\mathbb{R} \) je přímka \( x+2y-1=0 \) tečnou kružnice \( x^2+y^2=q^2 \) ?}
{\( \frac{\sqrt5}5 \)}{\( 5 \)}{\( \sqrt5 \)}{\( \frac15 \)}{\( 1 \)}{}}
\LANGsk{
\Question{
Pre ktorú hodnotu parametra \( q\in\mathbb{R} \) je priamka \( x+2y-1=0 \) dotyčnicou kružnice \( x^2+y^2=q^2 \) ?}
{\( \frac{\sqrt5}5 \)}{\( 5 \)}{\( \sqrt5 \)}{\( \frac15 \)}{\( 1 \)}{}}
\LANGes{
\Question{
Halla el valor del parámetro \( q\in\mathbb{R} \) para el que la recta \( x+2y-1=0 \) es tangente a la circunferencia \( x^2+y^2=q^2 \).}
{\( \frac{\sqrt5}5 \)}{\( 5 \)}{\( \sqrt5 \)}{\( \frac15 \)}{\( 1 \)}{}}
\End

\Data\ID{1003024109}
\Updated{2022-08-25 10:08:44}
\Level{C}
\SubArea{Conic Sections}
\LANGen{
\Question{
Find the set of values of the parameter \( c\in\mathbb{R} \) for which the line \( 3x+5y-c=0 \) is a secant of the ellipse \( 16x^2+25y^2=400 \).}
{\( (-25;25) \)}{\( (-\infty;-25)\cup(25;\infty) \)}{\( \{25\} \)}{\( \{-25\} \)}{\( (25;\infty) \)}{}}
\LANGpl{
\Question{
Wyznacz wartości parametru \( c\in\mathbb{R} \) tak, aby prosta \( 3x+5y-c=0 \) była sieczną elipsy \( 16x^2+25y^2=400 \).}
{\( (-25;25) \)}{\( (-\infty;-25)\cup(25;\infty) \)}{\( \{25\} \)}{\( \{-25\} \)}{\( (25;\infty) \)}{}}
\LANGcs{
\Question{
Množina všech hodnot parametru \( c\in\mathbb{R} \), pro které je přímka \( 3x+5y-c=0 \) sečnou elipsy \( 16x^2+25y^2=400 \), je:}
{\( (-25;25) \)}{\( (-\infty;-25)\cup(25;\infty) \)}{\( \{25\} \)}{\( \{-25\} \)}{\( (25;\infty) \)}{}}
\LANGsk{
\Question{
Množina všetkých hodnôt parametra \( c\in\mathbb{R} \), pre ktoré je priamka \( 3x+5y-c=0 \) sečnicou elipsy \( 16x^2+25y^2=400 \), je:}
{\( (-25;25) \)}{\( (-\infty;-25)\cup(25;\infty) \)}{\( \{25\} \)}{\( \{-25\} \)}{\( (25;\infty) \)}{}}
\LANGes{
\Question{
Halla el conjunto de valores del parámetro \( c\in\mathbb{R} \) para los cuales la recta \( 3x+5y-c=0 \) es secante a la elipse \( 16x^2+25y^2=400 \).}
{\( (-25;25) \)}{\( (-\infty;-25)\cup(25;\infty) \)}{\( \{25\} \)}{\( \{-25\} \)}{\( (25;\infty) \)}{}}
\End

\Data\ID{1003024110}
\Updated{2022-08-25 10:08:44}
\Level{C}
\SubArea{Conic Sections}
\LANGen{
\Question{
Find the value of the parameter \( q\in\mathbb{R} \) for which the line \( y=x+q \) is a tangent of the  parabola \( y^2=6x \).}
{\( \frac32 \)}{\( \frac23 \)}{\( 0 \)}{\( -\frac32 \)}{\( -\frac23 \)}{}}
\LANGpl{
\Question{
Wyznacz wartość parametru \( q\in\mathbb{R} \) tak, aby prosta \( y=x+q \) była styczną do paraboli \( y^2=6x \).}
{\( \frac32 \)}{\( \frac23 \)}{\( 0 \)}{\( -\frac32 \)}{\( -\frac23 \)}{}}
\LANGcs{
\Question{
Pro kterou hodnotu parametru \( q\in\mathbb{R} \) je přímka \( y=x+q \) tečnou paraboly \( y^2=6x \)?}
{\( \frac32 \)}{\( \frac23 \)}{\( 0 \)}{\( -\frac32 \)}{\( -\frac23 \)}{}}
\LANGsk{
\Question{
Pre ktorú hodnotu parametra\( q\in\mathbb{R} \) je priamka \( y=x+q \) dotyčnicou paraboly \( y^2=6x \)?}
{\( \frac32 \)}{\( \frac23 \)}{\( 0 \)}{\( -\frac32 \)}{\( -\frac23 \)}{}}
\LANGes{
\Question{
Encuentra el valor del parámetro \( q\in\mathbb{R} \) para el que la recta \( y=x+q \) es tangente a la parábola \( y^2=6x \).}
{\( \frac32 \)}{\( \frac23 \)}{\( 0 \)}{\( -\frac32 \)}{\( -\frac23 \)}{}}
\End

\Data\ID{1003024111}
\Updated{2023-06-05 09:06:34}
\Level{C}
\SubArea{Conic Sections}
\LANGen{
\Question{
The equation of a parabola is given by \( y^2 -12x -6y+57=0 \). Find the equation of  a line, which  passes through the vertex of this parabola and is  parallel to the line \( 5x-3y-2=0 \).}
{\( 5x-3y-11 = 0 \)}{\( -5x+3y-11 = 0 \)}{\( 5x-3y-3 = 0 \)}{\( 5x-3y+11 = 0 \)}{\(-5x+3y-3 = 0 \)}{}}
\LANGpl{
\Question{
Parabola jest opisana równaniem \( y^2 -12x -6y+57=0 \). Wyznacz równanie prostej przechodzącej przez wierzchołek tej  paraboli i równoległej do prostej \( 5x-3y-2=0 \).}
{\( 5x-3y-11 = 0 \)}{\( -5x+3y-11 = 0 \)}{\( 5x-3y-3 = 0 \)}{\( 5x-3y+11 = 0 \)}{\(-5x+3y-3 = 0 \)}{}}
\LANGcs{
\Question{
Parabola je dána rovnicí \( y^2 -12x -6y+57=0 \). Rovnice přímky, která prochází vrcholem paraboly a je rovnoběžná s přímkou \( 5x-3y-2=0 \), je:}
{\( 5x-3y-11 = 0 \)}{\( -5x+3y-11 = 0 \)}{\( 5x-3y-3 = 0 \)}{\( 5x-3y+11 = 0 \)}{\(-5x+3y-3 = 0 \)}{}}
\LANGsk{
\Question{
Parabola je daná rovnicou \( y^2 -12x -6y+57=0 \). Rovnice priamky, ktorá prechádza vrcholom paraboly a je rovnobežná s priamkou \( 5x-3y-2=0 \), je:}
{\( 5x-3y-11 = 0 \)}{\( -5x+3y-11 = 0 \)}{\( 5x-3y-3 = 0 \)}{\( 5x-3y+11 = 0 \)}{\(-5x+3y-3 = 0 \)}{}}
\LANGes{
\Question{
Dada la ecuación de la parábola \( y^2 -12x -6y+57=0 \). Halla la ecuación de la recta que pasa por el vértice de esta parábola y es paralela a la recta \( 5x-3y-2=0 \).}
{\( 5x-3y-11 = 0 \)}{\( -5x+3y-11 = 0 \)}{\( 5x-3y-3 = 0 \)}{\( 5x-3y+11 = 0 \)}{\(-5x+3y-3 = 0 \)}{}}
\End

\Data\ID{1003024201}
\Updated{2022-08-25 10:08:44}
\Level{C}
\SubArea{Conic Sections}
\LANGen{
\Question{
Find the tangent equations from the point \( K=[5;0] \)  to the circle given by \( (x-1)^2+(y-2)^2=4 \).}
{\( 4x+3y-20=0 \), \( y=0 \)}{\( 3x+4y-20=0 \), \( y=0 \)}{\( 4x+3y+20=0 \), \( y=0 \)}{\( -4x+3y-20=0 \), \( y=0 \)}{\( -3x+4y-20=0 \), \( y=0 \)}{}}
\LANGpl{
\Question{
Wyznacz równanie stycznej do okręgu opisanego równaniem \( (x-1)^2+(y-2)^2=4 \) przechodzącej przez punkt  \( K=[5;0] \)}
{\( 4x+3y-20=0 \), \( y=0 \)}{\( 3x+4y-20=0 \), \( y=0 \)}{\( 4x+3y+20=0 \), \( y=0 \)}{\( -4x+3y-20=0 \), \( y=0 \)}{\( -3x+4y-20=0 \), \( y=0 \)}{}}
\LANGcs{
\Question{
Rovnice tečen vedených z bodu \( K=[5;0] \)  ke kružnici o rovnici \( (x-1)^2+(y-2)^2=4 \)  jsou:}
{\( 4x+3y-20=0 \), \( y=0 \)}{\( 3x+4y-20=0 \), \( y=0 \)}{\( 4x+3y+20=0 \), \( y=0 \)}{\( -4x+3y-20=0 \), \( y=0 \)}{\( -3x+4y-20=0 \), \( y=0 \)}{}}
\LANGsk{
\Question{
Rovnice dotyčníc vedených z bodu \( K=[5;0] \)  ku kružnici s rovnicou \( (x-1)^2+(y-2)^2=4 \)  sú:}
{\( 4x+3y-20=0 \), \( y=0 \)}{\( 3x+4y-20=0 \), \( y=0 \)}{\( 4x+3y+20=0 \), \( y=0 \)}{\( -4x+3y-20=0 \), \( y=0 \)}{\( -3x+4y-20=0 \), \( y=0 \)}{}}
\LANGes{
\Question{
Elige las ecuaciones de las tangentes por punto \( K=(5;0) \) a la circunferencia dada por la ecuación \( (x-1)^2+(y-2)^2=4 \).}
{\( 4x+3y-20=0 \), \( y=0 \)}{\( 3x+4y-20=0 \), \( y=0 \)}{\( 4x+3y+20=0 \), \( y=0 \)}{\( -4x+3y-20=0 \), \( y=0 \)}{\( -3x+4y-20=0 \), \( y=0 \)}{}}
\End

\Data\ID{1003024202}
\Updated{2022-08-25 10:08:44}
\Level{C}
\SubArea{Conic Sections}
\LANGen{
\Question{
Find the tangent equations from the point \( M=[0;0] \)   to the ellipse given by \( x^2+2y^2-8x+4y+12=0\).}
{\( x-5y=0 \), \( x+y=0 \)}{\( 5x-y=0 \), \( x+y=0 \)}{\( x+5y=0 \), \( x-y=0 \)}{\( -x-5y=0 \), \( -x+y=0 \)}{\( 5x+y=0 \), \( x-y=0 \)}{}}
\LANGpl{
\Question{
Wyznacz równania stycznej do elipsy opisanej równaniem  \( x^2+2y^2-8x+4y+12=0\) przechodzącej przez punkt \( M=[0;0] \).}
{\( x-5y=0 \), \( x+y=0 \)}{\( 5x-y=0 \), \( x+y=0 \)}{\( x+5y=0 \), \( x-y=0 \)}{\( -x-5y=0 \), \( -x+y=0 \)}{\( 5x+y=0 \), \( x-y=0 \)}{}}
\LANGcs{
\Question{
Rovnice tečen vedených z bodu  \( M=[0;0] \)   k elipse o rovnici  \( x^2+2y^2-8x+4y+12=0\)  jsou:}
{\( x-5y=0 \), \( x+y=0 \)}{\( 5x-y=0 \), \( x+y=0 \)}{\( x+5y=0 \), \( x-y=0 \)}{\( -x-5y=0 \), \( -x+y=0 \)}{\( 5x+y=0 \), \( x-y=0 \)}{}}
\LANGsk{
\Question{
Rovnice dotyčníc vedených z bodu  \( M=[0;0] \)   k elipse s rovnicou  \( x^2+2y^2-8x+4y+12=0\)  sú:}
{\( x-5y=0 \), \( x+y=0 \)}{\( 5x-y=0 \), \( x+y=0 \)}{\( x+5y=0 \), \( x-y=0 \)}{\( -x-5y=0 \), \( -x+y=0 \)}{\( 5x+y=0 \), \( x-y=0 \)}{}}
\LANGes{
\Question{
Dada la ecuación de la elipse \( x^2+2y^2-8x+4y+12=0\), halla las ecuaciones de las tangentes por el punto \( M=(0;0) \) a la elipse.}
{\( x-5y=0 \), \( x+y=0 \)}{\( 5x-y=0 \), \( x+y=0 \)}{\( x+5y=0 \), \( x-y=0 \)}{\( -x-5y=0 \), \( -x+y=0 \)}{\( 5x+y=0 \), \( x-y=0 \)}{}}
\End

\Data\ID{1003024203}
\Updated{2022-08-25 10:08:44}
\Level{C}
\SubArea{Conic Sections}
\LANGen{
\Question{
Find the tangent equations from the point \( L=[1;-1] \)  to the parabola given by \( y^2-4x+2y+9=0 \).}
{\( x-y-2=0 \), \( x+y=0 \)}{\( x+y-2=0 \), \( x+y=0 \)}{\( x+y+2=0 \), \( x-y=0 \)}{\( x-y+2=0 \), \( -x+y=0 \)}{\( -x+y-2=0 \), \( x-y=0 \)}{}}
\LANGpl{
\Question{
Wyznacz równania stycznej do paraboli opisanej równaniem  \( y^2-4x+2y+9=0 \) przechodzącej przez punkt \( L=[1;-1] \)}
{\( x-y-2=0 \), \( x+y=0 \)}{\( x+y-2=0 \), \( x+y=0 \)}{\( x+y+2=0 \), \( x-y=0 \)}{\( x-y+2=0 \), \( -x+y=0 \)}{\( -x+y-2=0 \), \( x-y=0 \)}{}}
\LANGcs{
\Question{
Rovnice tečen vedených z bodu  \( L=[1;-1] \)  k parabole o rovnici  \( y^2-4x+2y+9=0 \)  jsou:}
{\( x-y-2=0 \), \( x+y=0 \)}{\( x+y-2=0 \), \( x+y=0 \)}{\( x+y+2=0 \), \( x-y=0 \)}{\( x-y+2=0 \), \( -x+y=0 \)}{\( -x+y-2=0 \), \( x-y=0 \)}{}}
\LANGsk{
\Question{
Rovnice dotyčníc vedených z bodu  \( L=[1;-1] \)  k parabole s rovnicou  \( y^2-4x+2y+9=0 \)  sú:}
{\( x-y-2=0 \), \( x+y=0 \)}{\( x+y-2=0 \), \( x+y=0 \)}{\( x+y+2=0 \), \( x-y=0 \)}{\( x-y+2=0 \), \( -x+y=0 \)}{\( -x+y-2=0 \), \( x-y=0 \)}{}}
\LANGes{
\Question{
Dada la ecuación de la parábola \( y^2-4x+2y+9=0 \), halla las ecuaciones de las tangentes por el punto \( L=(1;-1) \).}
{\( x-y-2=0 \), \( x+y=0 \)}{\( x+y-2=0 \), \( x+y=0 \)}{\( x+y+2=0 \), \( x-y=0 \)}{\( x-y+2=0 \), \( -x+y=0 \)}{\( -x+y-2=0 \), \( x-y=0 \)}{}}
\End

\Data\ID{2010005903}
\Updated{2022-05-16 11:05:31}
\Level{C}
\SubArea{Conic Sections}
\LANGen{
\Question{
Find the values of the real parameter \(r\) 
which ensure that the line \(x = r\) 
is a tangent to the circle 
\[ 
 x^{2} + y^{2} - 4x + 10y + 4 = 0.
\]}
{\(\{-3;7\}\)}{\(\{ - 7;3\}\)}{\(\{ - 10;0\}\)}{\(\{ 0;10\}\)}{}{}}
\LANGpl{
\Question{
Znajdź wartości rzeczywistego parametru \(r\)
które sprawiają, że prosta \(x = r\)
jest styczną do okręgu
\[ 
 x^{2} + y^{2} - 4x + 10y + 4 = 0.
\]}
{\(\{-3;7\}\)}{\(\{ - 7;3\}\)}{\(\{ - 10;0\}\)}{\(\{ 0;10\}\)}{}{}}
\LANGcs{
\Question{
Určete množinu všech hodnot reálného parametru \(r\), pro které je přímka
 \(x = r\) 
tečnou ke kružnici
\[ 
 x^{2} + y^{2} - 4x + 10y + 4 = 0.
\]}
{\(\{-3;7\}\)}{\(\{ - 7;3\}\)}{\(\{ - 10;0\}\)}{\(\{ 0;10\}\)}{}{}}
\LANGsk{
\Question{
Množina všetkých hodnôt reálneho parametra \(r,\) pre ktoré je priamka \(x = r\) dotyčnicou kružnice \( x^{2} + y^{2} - 4x + 10y + 4 = 0, \) je:}
{\(\{-3;7\}\)}{\(\{ - 7;3\}\)}{\(\{ - 10;0\}\)}{\(\{ 0;10\}\)}{}{}}
\LANGes{
\Question{
Encuentra todos los valores del parámetro real \(r\) para que la recta \(x = r\) sea tangente a la circunferencia  
\[ 
 x^{2} + y^{2} - 4x + 10y + 4 = 0.
\]}
{\(\{-3;7\}\)}{\(\{ - 7;3\}\)}{\(\{ - 10;0\}\)}{\(\{ 0;10\}\)}{}{}}
\End

\Data\ID{2010005904}
\Updated{2022-05-23 02:05:40}
\Level{C}
\SubArea{Conic Sections}
\LANGen{
\Question{
Find a true statement about the ellipse 
\[ 
9 x^{2} + y^{2} + 18x = 0.
\]}
{The tangent to the ellipse can pass through any point on the line
\(x = 1\).}{The tangent to the ellipse can pass through any point on the line
\(y = 1\).}{The tangent to the ellipse can pass through the point
\([-1;1]\).}{The tangent to the ellipse can pass through any point on the line
\(y = -1\).}{}{}}
\LANGpl{
\Question{
Wybierz prawdziwe stwierdzenie dotyczące elipsy:
\[ 
9 x^{2} + y^{2} + 18x = 0.
\]}
{Styczna do elipsy może przechodzić przez dowolny punkt na prostej
\(x = 1\).}{Styczna do elipsy może przechodzić przez dowolny punkt na prostej
\(y = 1\).}{Styczna do elipsy może przechodzić przez punkt
\([-1;1]\).}{Styczna do elipsy może przechodzić przez dowolny punkt na prostej
\(y = -1\).}{}{}}
\LANGcs{
\Question{
Vyberte pravdivé tvrzení o elipse
\[ 
9 x^{2} + y^{2} + 18x = 0.
\]}
{Tečna elipsy může procházet libovolným bodem přímky \(x = 1\).}{Tečna elipsy může procházet libovolným bodem přímky \(y = 1\).}{Tečna elipsy může procházet bodem \([-1;1]\).}{Tečna elipsy může procházet libovolným bodem přímky \(y = -1\).}{}{}}
\LANGsk{
\Question{
Vytvorte pravdivé tvrdenie: Dotyčnicu k elipse \( 9 x^{2} + y^{2} + 18x = 0 \) možno viesť:}
{ľubovolným bodom priamky
\(x = 1\).}{ľubovolným bodom priamky
\(y = 1\).}{bodom
\([-1;1]\).}{ľubovolným bodom priamky
\(y = -1\).}{}{}}
\LANGes{
\Question{
Elige el enunciado correcto relacionado con la elipse  
\[ 
9 x^{2} + y^{2} + 18x = 0.
\]}
{La tangente a la elipse puede pasar por cualquier punto de la recta 
\(x = 1\).}{La tangente a la elipse puede pasar por cualquier punto de la recta 
\(y = 1\).}{La tangente a la elipse puede pasar por el punto 
\((-1;1)\).}{La tangente a la elipse puede pasar por cualquier punto de la recta 
\(y = -1\).}{}{}}
\End

\Data\ID{2010005905}
\Updated{2022-05-23 02:05:58}
\Level{C}
\SubArea{Conic Sections}
\LANGen{
\Question{
Find the tangent line \(r\) to
the parabola \(6(x+1) = (y-3)^{2}\), so that the tangent \(r\)
is parallel to the line \(q\colon 3x - 2y + 7 = 0.\)}
{\(r\colon 3x - 2y + 11 = 0\)}{\(r\colon 3x - 2y - 7 = 0\)}{\(r\colon 3x - 2y - 13 = 0\)}{\(r\colon 3x -2y + 13 = 0\)}{}{}}
\LANGpl{
\Question{
Wyznacz styczną \(r\) do
paraboli \(6(x+1) = (y-3)^{2}\), tak aby styczna \(r\)
była  równoległa do prostej \(q\colon 3x - 2y + 7 = 0.\)}
{\(r\colon 3x - 2y + 11 = 0\)}{\(r\colon 3x - 2y - 7 = 0\)}{\(r\colon 3x - 2y - 13 = 0\)}{\(r\colon 3x -2y + 13 = 0\)}{}{}}
\LANGcs{
\Question{
Parabole \(6(x+1) = (y-3)^{2}\) určete tečnu, která je rovnoběžná k přímce \(q\colon 3x - 2y + 7 = 0.\)}
{\(3x - 2y + 11 = 0\)}{\(3x - 2y - 7 = 0\)}{\(3x - 2y - 13 = 0\)}{\(3x -2y + 13 = 0\)}{}{}}
\LANGsk{
\Question{
Rovnice dotyčnice paraboly \(6(x+1) = (y-3)^{2},\)  ktorá je rovnobežná s priamkou \(q\colon 3x - 2y + 7 = 0,\) má tvar:}
{\(3x - 2y + 11 = 0\)}{\(3x - 2y - 7 = 0\)}{\(3x - 2y - 13 = 0\)}{\(3x -2y + 13 = 0\)}{}{}}
\LANGes{
\Question{
Encuentra la tangente \(r\) a la parábola \(6(x+1) = (y-3)^{2}\) que sea paralela a la recta  \(q\colon 3x - 2y + 7 = 0.\)}
{\(r\colon 3x - 2y + 11 = 0\)}{\(r\colon 3x - 2y - 7 = 0\)}{\(r\colon 3x - 2y - 13 = 0\)}{\(r\colon 3x -2y + 13 = 0\)}{}{}}
\End

\Data\ID{2010005906}
\Updated{2022-05-16 11:05:31}
\Level{C}
\SubArea{Conic Sections}
\LANGen{
\Question{
Find all the tangents to the hyperbola
\(y^{2} - 4x^{2} = 12\) 
such that the angle between each tangent and the
\(x\)-axis
is \(45^{\circ }\).}
{\(y = x + 3,\ y = x - 3,\ y = -x + 3,\ y = -x - 3\)}{\(y = x + 3,\ y =- x - 3\)}{\(y = x + 3,\ y = x - 3\)}{\(y = x + 3\)}{}{}}
\LANGpl{
\Question{
Znajdź wszystkie styczne do hiperboli
\(y^{2} - 4x^{2} = 12\)
takie, że kąt między każdą styczną a
osiou \(x\) to \(45^{\circ }\).}
{\(y = x + 3,\ y = x - 3,\ y = -x + 3,\ y = -x - 3\)}{\(y = x + 3,\ y =- x - 3\)}{\(y = x + 3,\ y = x - 3\)}{\(y = x + 3\)}{}{}}
\LANGcs{
\Question{
Najděte všechny tečny hyperboly
\(y^{2} - 4x^{2} = 12\), které mají s osou \(x\) odchylku \(45^{\circ }\).}
{\(y = x + 3,\ y = x - 3,\ y = -x + 3,\ y = -x - 3\)}{\(y = x + 3,\ y =- x - 3\)}{\(y = x + 3,\ y = x - 3\)}{\(y = x + 3\)}{}{}}
\LANGsk{
\Question{
Všetky dotyčnice hyperboly \(y^{2} - 4x^{2} = 12,\) ktorých odchýlka s osou \(x\) je \(45^{\circ }, \) majú rovnice:}
{\(y = x + 3,\ y = x - 3,\ y = -x + 3,\ y = -x - 3\)}{\(y = x + 3,\ y =- x - 3\)}{\(y = x + 3,\ y = x - 3\)}{\(y = x + 3\)}{}{}}
\LANGes{
\Question{
Halla todas las tangentes a la hipérbola \(y^{2} - 4x^{2} = 12\) 
para las que el ángulo entre cada tangente y el eje \(x\) es \(45^{\circ }\).}
{\(y = x + 3,\ y = x - 3,\ y = -x + 3,\ y = -x - 3\)}{\(y = x + 3,\ y =- x - 3\)}{\(y = x + 3,\ y = x - 3\)}{\(y = x + 3\)}{}{}}
\End

\Data\ID{2010006001}
\Updated{2022-05-16 11:05:35}
\Level{C}
\SubArea{Conic Sections}
\LANGen{
\Question{
Find the value of the parameter \( r\in\mathbb{R} \) for which the line \( 2x-y-1=0 \) is a tangent of the hyperbola  \( 2x^2-y^2=r \).}
{\( 1 \)}{\( -1 \)}{\( 2 \)}{\( -3 \)}{\( 3 \)}{}}
\LANGpl{
\Question{
Znajdź wartość parametru \( r\in\mathbb{R} \), dla którego prosta \( 2x-y-1=0 \) jest styczną do hiperboli \( 2x^2-y^2=r \).}
{\( 1 \)}{\( -1 \)}{\( 2 \)}{\( -3 \)}{\( 3 \)}{}}
\LANGcs{
\Question{
Určete hodnotu parametru \( r\in\mathbb{R} \), pro kterou je přímka \( 2x-y-1=0 \) tečnou hyperboly  \( 2x^2-y^2=r \).}
{\( 1 \)}{\( -1 \)}{\( 2 \)}{\( -3 \)}{\( 3 \)}{}}
\LANGsk{
\Question{
Pre ktorú hodnotu parametra \( r\in\mathbb{R} \) je priamka \( 2x-y-1=0 \) dotyčnica hyperboly  \( 2x^2-y^2=r \) ?}
{\( 1 \)}{\( -1 \)}{\( 2 \)}{\( -3 \)}{\( 3 \)}{}}
\LANGes{
\Question{
Encuentra el valor del parámetro \( r\in\mathbb{R} \)  para el cual la recta \( 2x-y-1=0 \) es tangente a la hipérbola \( 2x^2-y^2=r \).}
{\( 1 \)}{\( -1 \)}{\( 2 \)}{\( -3 \)}{\( 3 \)}{}}
\End

\Data\ID{2010006002}
\Updated{2022-05-16 11:05:35}
\Level{C}
\SubArea{Conic Sections}
\LANGen{
\Question{
Find the value of the parameter \( p\in\mathbb{R} \) for which the line \( 2x-y-1=0 \) is a tangent of the  parabola \( x^2=2py \).}
{\( \frac12 \)}{\(- \frac12 \)}{\( 2 \)}{\( -2 \)}{\( 1 \)}{}}
\LANGpl{
\Question{
Znajdź wartość parametru \( p\in\mathbb{R} \), dla którego prosta \( 2x-y-1=0 \) jest styczną do paraboli \( x^2=2py \).}
{\( \frac12 \)}{\(- \frac12 \)}{\( 2 \)}{\( -2 \)}{\( 1 \)}{}}
\LANGcs{
\Question{
Určete hodnotu parametru \( p\in\mathbb{R} \), pro kterou je přímka \( 2x-y-1=0 \) tečnou paraboly \( x^2=2py \).}
{\( \frac12 \)}{\(- \frac12 \)}{\( 2 \)}{\( -2 \)}{\( 1 \)}{}}
\LANGsk{
\Question{
Pre ktorú hodnotu parametra \( p\in\mathbb{R} \) je priamka \( 2x-y-1=0 \) dotyčnicou paraboly \( x^2=2py \)?}
{\( \frac12 \)}{\(- \frac12 \)}{\( 2 \)}{\( -2 \)}{\( 1 \)}{}}
\LANGes{
\Question{
Halla el valor del parámetro \( p\in\mathbb{R} \) para el cual la recta \( 2x-y-1=0 \) es tangente a la parábola \( x^2=2py \).}
{\( \frac12 \)}{\(- \frac12 \)}{\( 2 \)}{\( -2 \)}{\( 1 \)}{}}
\End

\Data\ID{2010006003}
\Updated{2022-05-23 02:05:34}
\Level{C}
\SubArea{Conic Sections}
\LANGen{
\Question{
Find the set of values of the parameter \( c\in\mathbb{R} \) for which the line \( 5x-3y-c=0 \)  isn't a secant of the ellipse \( 25x^2+16y^2=400 \).}
{\( (-\infty;-25 ] \cup [ 25;\infty) \)}{\( (-25;25) \)}{\( \{-25,25\} \)}{\( (-\infty;-25) \cup (25;\infty) \)}{\( [ 25;\infty) \)}{}}
\LANGpl{
\Question{
Znajdź zbiór wartości parametru \( c\in\mathbb{R} \), dla którego prosta \( 5x-3y-c=0 \)  nie jest sieczną elipsy \( 25x^2+16y^2=400 \).}
{\( (-\infty;-25 \rangle \cup \langle 25;\infty) \)}{\( (-25;25) \)}{\( \{-25,25\} \)}{\( (-\infty;-25) \cup (25;\infty) \)}{\( \langle 25;\infty) \)}{}}
\LANGcs{
\Question{
Určete množinu všech hodnot parametru \( c\in\mathbb{R} \), pro které přímka \( 5x-3y-c=0 \) není sečnou elipsy \( 25x^2+16y^2=400 \).}
{\( (-\infty;-25 \rangle \cup \langle 25;\infty) \)}{\( (-25;25) \)}{\( \{-25,25\} \)}{\( (-\infty;-25) \cup (25;\infty) \)}{\( \langle 25;\infty) \)}{}}
\LANGsk{
\Question{
Množina všetkých hodnôt parametra \( c\in\mathbb{R}, \) pre ktoré priamka \( 5x-3y-c=0 \)  nie je sečnicou elipsy  \( 25x^2+16y^2=400, \) je:}
{\( (-\infty;-25 \rangle \cup \langle 25;\infty) \)}{\( (-25;25) \)}{\( \{-25,25\} \)}{\( (-\infty;-25) \cup (25;\infty) \)}{\( \langle 25;\infty) \)}{}}
\LANGes{
\Question{
Halla el conjunto de valores del parámetro \( c\in\mathbb{R} \) para los cuales la recta \( 5x-3y-c=0 \)  no es secante a la elipse \( 25x^2+16y^2=400 \).}
{\( (-\infty;-25 ] \cup [ 25;\infty) \)}{\( (-25;25) \)}{\( \{-25,25\} \)}{\( (-\infty;-25) \cup (25;\infty) \)}{\( [ 25;\infty) \)}{}}
\End

\Data\ID{2010006004}
\Updated{2023-06-05 09:06:04}
\Level{C}
\SubArea{Conic Sections}
\LANGen{
\Question{
The equation of a parabola is given by \( x^2 -8x +3y-2=0 \). Find the equation of  a line, which  passes through the vertex of this parabola and is  parallel to the line \( 2x-5y+8=0 \).}
{\( -2x+5y-22 = 0 \)}{\( 2x-5y-22 = 0 \)}{\( 2x-5y-38 = 0 \)}{\( 2x-5y+38 = 0 \)}{\( -2x+5y+22 = 0 \)}{}}
\LANGpl{
\Question{
Równanie paraboli jest podane przez \( x^2 -8x +3y-2=0 \). Znajdź równanie prostej, która przechodzi przez wierzchołek tej paraboli i jest równoległa do prostej \( 2x-5y+8=0 \).}
{\( -2x+5y-22 = 0 \)}{\( 2x-5y-22 = 0 \)}{\( 2x-5y-38 = 0 \)}{\( 2x-5y+38 = 0 \)}{\( -2x+5y+22 = 0 \)}{}}
\LANGcs{
\Question{
Parabola je dána rovnicí \( x^2 -8x +3y-2=0 \). Určete rovnici přímky, která prochází vrcholem paraboly a je rovnoběžná s přímkou \( 2x-5y+8=0 \).}
{\( -2x+5y-22 = 0 \)}{\( 2x-5y-22 = 0 \)}{\( 2x-5y-38 = 0 \)}{\( 2x-5y+38 = 0 \)}{\( -2x+5y+22 = 0 \)}{}}
\LANGsk{
\Question{
Parabola je daná rovnicou \( x^2 -8x +3y-2=0 \). Rovnica priamky, ktorá prechádza vrcholom paraboly a je rovnobežná s priamkou \( 2x-5y+8=0, \) je:}
{\( -2x+5y-22 = 0 \)}{\( 2x-5y-22 = 0 \)}{\( 2x-5y-38 = 0 \)}{\( 2x-5y+38 = 0 \)}{\( -2x+5y+22 = 0 \)}{}}
\LANGes{
\Question{
Dada la ecuación de la parábola \( x^2 -8x +3y-2=0 \). Halla la ecuación de la recta que pasa por el vértice de esta parábola y es paralela a la recta \( 2x-5y+8=0 \).}
{\( -2x+5y-22 = 0 \)}{\( 2x-5y-22 = 0 \)}{\( 2x-5y-38 = 0 \)}{\( 2x-5y+38 = 0 \)}{\( -2x+5y+22 = 0 \)}{}}
\End

\Data\ID{2010006005}
\Updated{2022-05-23 02:05:01}
\Level{C}
\SubArea{Conic Sections}
\LANGen{
\Question{
Find the tangent equations from the point \( N=[0;0] \)   to the ellipse given by \( 2x^2+y^2-4x-8y+12=0\).}
{\( 5x+y=0 \), \( x-y=0 \)}{\( 5x-y=0 \), \( x+y=0 \)}{\( 5x-y=0 \), \( x-y=0 \)}{\( x+5y=0 \), \( x-y=0 \)}{\( x-5y=0 \), \( x+y=0 \)}{}}
\LANGpl{
\Question{
Znajdź równania styczne od punktu \( N=[0;0] \) do elipsy podanej przez \( 2x^2+y^2-4x-8y+12=0\).}
{\( 5x+y=0 \), \( x-y=0 \)}{\( 5x-y=0 \), \( x+y=0 \)}{\( 5x-y=0 \), \( x-y=0 \)}{\( x+5y=0 \), \( x-y=0 \)}{\( x-5y=0 \), \( x+y=0 \)}{}}
\LANGcs{
\Question{
Určete rovnice tečen vedených z bodu \( N=[0;0] \) k elipse \( 2x^2+y^2-4x-8y+12=0\).}
{\( 5x+y=0 \), \( x-y=0 \)}{\( 5x-y=0 \), \( x+y=0 \)}{\( 5x-y=0 \), \( x-y=0 \)}{\( x+5y=0 \), \( x-y=0 \)}{\( x-5y=0 \), \( x+y=0 \)}{}}
\LANGsk{
\Question{
Rovnice dotyčníc, vedených z bodu \( N=[0;0] \) k elipse s rovnicou \( 2x^2+y^2-4x-8y+12=0\), sú:}
{\( 5x+y=0 \), \( x-y=0 \)}{\( 5x-y=0 \), \( x+y=0 \)}{\( 5x-y=0 \), \( x-y=0 \)}{\( x+5y=0 \), \( x-y=0 \)}{\( x-5y=0 \), \( x+y=0 \)}{}}
\LANGes{
\Question{
Dada la ecuación de la elipse \( 2x^2+y^2-4x-8y+12=0\), halla las ecuaciones de las tangentes en el punto \( N=[0;0] \).}
{\( 5x+y=0 \), \( x-y=0 \)}{\( 5x-y=0 \), \( x+y=0 \)}{\( 5x-y=0 \), \( x-y=0 \)}{\( x+5y=0 \), \( x-y=0 \)}{\( x-5y=0 \), \( x+y=0 \)}{}}
\End

\Data\ID{2010006006}
\Updated{2022-05-16 11:05:35}
\Level{C}
\SubArea{Conic Sections}
\LANGen{
\Question{
Find the tangent equations from the point \( M=[-2;2] \)  to the parabola given by \( x^2+4x-4y+16=0 \).}
{\( x-y+4=0 \), \( x+y=0 \)}{\( x+y-4=0 \), \( x+y=0 \)}{\( x+y+4=0 \), \( x-y=0 \)}{\( x-y-4=0 \), \( x-y=0 \)}{\( x+y-4=0 \), \( x-y=0 \)}{}}
\LANGpl{
\Question{
Znajdź równania styczne od punktu \( M=[-2;2] \) do paraboli podanej przez \( x^2+4x-4y+16=0 \).}
{\( x-y+4=0 \), \( x+y=0 \)}{\( x+y-4=0 \), \( x+y=0 \)}{\( x+y+4=0 \), \( x-y=0 \)}{\( x-y-4=0 \), \( x-y=0 \)}{\( x+y-4=0 \), \( x-y=0 \)}{}}
\LANGcs{
\Question{
Určete rovnice tečen z bodu \( M=[-2;2] \) k parabole \( x^2+4x-4y+16=0 \).}
{\( x-y+4=0 \), \( x+y=0 \)}{\( x+y-4=0 \), \( x+y=0 \)}{\( x+y+4=0 \), \( x-y=0 \)}{\( x-y-4=0 \), \( x-y=0 \)}{\( x+y-4=0 \), \( x-y=0 \)}{}}
\LANGsk{
\Question{
Rovnice dotyčníc, vedených z bodu \( M=[-2;2] \)  k parabole s rovnicou \( x^2+4x-4y+16=0, \) sú:}
{\( x-y+4=0 \), \( x+y=0 \)}{\( x+y-4=0 \), \( x+y=0 \)}{\( x+y+4=0 \), \( x-y=0 \)}{\( x-y-4=0 \), \( x-y=0 \)}{\( x+y-4=0 \), \( x-y=0 \)}{}}
\LANGes{
\Question{
Dada la ecuación de la parábola \( x^2+4x-4y+16=0 \), halla las ecuaciones de las tangentes por el punto \( M=[-2;2] \).}
{\( x-y+4=0 \), \( x+y=0 \)}{\( x+y-4=0 \), \( x+y=0 \)}{\( x+y+4=0 \), \( x-y=0 \)}{\( x-y-4=0 \), \( x-y=0 \)}{\( x+y-4=0 \), \( x-y=0 \)}{}}
\End

\Data\ID{2010006308}
\Updated{2022-05-16 11:05:29}
\Level{C}
\SubArea{Conic Sections}
\LANGen{
\Question{
Find the equation of the  circle, which passes through the points \( A=[-4;2] \), \( B=[1;3] \) and \( C=[2;-2] \).}
{\( (x+1)^2+y^2 = 13 \)}{\( (x-1)^2+y^2 = 13 \)}{\( x^2 + (y-1)^2 = 13 \)}{\( x^2 + y^2 = 13x \)}{\( x^2 + (y+1)^2 = 13 \)}{}}
\LANGpl{
\Question{
Znajdź równanie okręgu, który przechodzi przez punkty \( A=[-4;2] \), \( B=[1;3] \) i \( C=[2;-2] \).}
{\( (x+1)^2+y^2 = 13 \)}{\( (x-1)^2+y^2 = 13 \)}{\( x^2 + (y-1)^2 = 13 \)}{\( x^2 + y^2 = 13x \)}{\( x^2 + (y+1)^2 = 13 \)}{}}
\LANGcs{
\Question{
Určete rovnici kružnice, která prochází body \( A=[-4;2] \), \( B=[1;3] \) a \( C=[2;-2] \).}
{\( (x+1)^2+y^2 = 13 \)}{\( (x-1)^2+y^2 = 13 \)}{\( x^2 + (y-1)^2 = 13 \)}{\( x^2 + y^2 = 13x \)}{\( x^2 + (y+1)^2 = 13 \)}{}}
\LANGsk{
\Question{
Rovnica kružnice, ktorá prechádza bodmi \( A=[-4;2] \), \( B=[1;3] \) a \( C=[2;-2] \) je:}
{\( (x+1)^2+y^2 = 13 \)}{\( (x-1)^2+y^2 = 13 \)}{\( x^2 + (y-1)^2 = 13 \)}{\( x^2 + y^2 = 13x \)}{\( x^2 + (y+1)^2 = 13 \)}{}}
\LANGes{
\Question{
Encuentra la ecuación de la circunferencia que pasa por los puntos \( A=[-4;2] \), \( B=[1;3] \) y \( C=[2;-2] \).}
{\( (x+1)^2+y^2 = 13 \)}{\( (x-1)^2+y^2 = 13 \)}{\( x^2 + (y-1)^2 = 13 \)}{\( x^2 + y^2 = 13x \)}{\( x^2 + (y+1)^2 = 13 \)}{}}
\End

\Data\ID{2010006309}
\Updated{2022-05-16 11:05:29}
\Level{C}
\SubArea{Conic Sections}
\LANGen{
\Question{
Find the points of intersection of the circle \( x^2+y^2=16 \) and the line \( x-y+4=0 \).}
{\( [-4;0] \), \( [0;4] \)}{\( [0;-4] \), \( [-4;0] \)}{\( [4;0] \), \( [0;4] \)}{\( [-4;0] \), \( [4;0] \)}{\( [4;0] \), \( [0;-4] \)}{}}
\LANGpl{
\Question{
Znajdź punkty przecięcia okręgu \( x^2+y^2=16 \) i prostej \( x-y+4=0 \).}
{\( [-4;0] \), \( [0;4] \)}{\( [0;-4] \), \( [-4;0] \)}{\( [4;0] \), \( [0;4] \)}{\( [-4;0] \), \( [4;0] \)}{\( [4;0] \), \( [0;-4] \)}{}}
\LANGcs{
\Question{
Určete průsečíky kružnice \( x^2+y^2=16 \) a přímky \( x-y+4=0 \).}
{\( [-4;0] \), \( [0;4] \)}{\( [0;-4] \), \( [-4;0] \)}{\( [4;0] \), \( [0;4] \)}{\( [-4;0] \), \( [4;0] \)}{\( [4;0] \), \( [0;-4] \)}{}}
\LANGsk{
\Question{
Spoločné body kružnice \( x^2+y^2=16 \) a priamky \( x-y+4=0 \) sú:}
{\( [-4;0] \), \( [0;4] \)}{\( [0;-4] \), \( [-4;0] \)}{\( [4;0] \), \( [0;4] \)}{\( [-4;0] \), \( [4;0] \)}{\( [4;0] \), \( [0;-4] \)}{}}
\LANGes{
\Question{
Halla los puntos de la intersección de la circunferencia \( x^2+y^2=16 \) y la recta \( x-y+4=0 \).}
{\( [-4;0] \), \( [0;4] \)}{\( [0;-4] \), \( [-4;0] \)}{\( [4;0] \), \( [0;4] \)}{\( [-4;0] \), \( [4;0] \)}{\( [4;0] \), \( [0;-4] \)}{}}
\End

\Data\ID{9000104801}
\Updated{2021-10-14 06:10:09}
\Level{C}
\SubArea{Conic Sections}
\LANGen{
\Question{
Consider the hyperbola 
\[ 
 xy = -1
\]
and a line \(p\) 
parallel to one of the axes but not identical with this axis. Find the true statement.}
{The line \(p\) 
has a unique intersection with the hyperbola.}{The line \(p\) 
has two intersections with the hyperbola.}{The line \(p\) 
does not have any intersection with the hyperbola.}{We cannot draw any conclusion on the number of intersections of the line
\(p\) with
the hyperbola.}{}{}}
\LANGpl{
\Question{
Dana jest hiperbola
\[ 
 xy = -1
\]
i prosta  \(p\) równoległa do jednej z osi, ale nie pokrywa się z tą osią. Oznacz zdanie prawdziwe.}
{Prosta \(p\) 
ma dokładnie jeden punkt wspólny z hiperbolą.}{Prosta \(p\) 
ma dwa punkty wspólne z hiperbolą.}{Prosta \(p\) 
nie ma punktów wspólnych z hiperbolą.}{Na podstawie podanych informacji nie można stwierdzić ile punktów wspólnych ma prosta 
\(p\) z hiperbolą.}{}{}}
\LANGcs{
\Question{
Je dána hyperbola \(xy = -1\) 
a přímka \(p\),
která je rovnoběžná s některou ze souřadnicových os. Zároveň víme, že
přímka \(p\) 
není s žádnou ze souřadnicových os totožná. Pak lze tvrdit, že:}
{Přímka \(p\) 
má s danou hyperbolou společný právě jeden bod.}{Přímka \(p\) 
má s danou hyperbolou společné právě dva body.}{Přímka \(p\) 
nemá s danou hyperbolou společný žádný bod.}{Z daných informací není možné jednoznačně určit
počet společných bodů dané hyperboly a přímky
\(p\).}{}{}}
\LANGsk{
\Question{
Je daná hyperbola \(xy = -1\) 
a priamka \(p\),
ktorá je rovnobežná s niektorou zo súradnicových osí. Zároveň vieme, že
priamka \(p\) 
nie je so žiadnou zo súradnicových osí totožná. Potom môžme tvrdiť, že:}
{Priamka \(p\) 
má s danou hyperbolou spoločný práve jeden bod.}{Priamka \(p\) 
má s danou hyperbolou spoločné práve dva body.}{Priamka \(p\) 
nemá s danou hyperbolou spoločný žiadny bod.}{Z daných informácií nie je možné jednoznačne určiť
počet spoločných bodov danej hyperboly a priamky
\(p\).}{}{}}
\LANGes{
\Question{
Consideremos la hipérbola 
\[ 
 xy = -1
\]
y la recta \(p\) 
paralela a uno de los ejes pero no idéntica a él. Encuentra el enunciado verdadero.}
{La recta \(p\) 
tiene una única intersección con la hipérbola.}{La recta \(p\) 
tiene dos intersecciones con la hipérbola.}{La recta \(p\) 
no tiene ninguna intersección con la hipérbola.}{Con la información dada no es posible averiguar el número de intersecciones de la recta
\(p\) con
la hipérbola.}{}{}}
\End

\Data\ID{9000104803}
\Updated{2021-10-14 06:10:06}
\Level{C}
\SubArea{Conic Sections}
\LANGen{
\Question{
Consider the hyperbola 
\[ 
 \frac{x^{2}} 
{16} -\frac{y^{2}} 
 {4} = 1
\]
and a line \(p\) 
parallel to one of the axes. Find the true statement.}
{We cannot draw any conclusion on the number of intersections of the line
\(p\) with
the hyperbola.}{The line \(p\) 
has two intersections with the hyperbola.}{The line \(p\) 
has a unique intersection with the hyperbola.}{The line \(p\) 
does not have any intersection with the hyperbola.}{}{}}
\LANGpl{
\Question{
Dana jest hiperbola
\[ 
 \frac{x^{2}} 
{16} -\frac{y^{2}} 
 {4} = 1
\]
i prosta \(p\) 
równoległa do jednej z osi. Oznacz zdanie prawdziwe.}
{Na podstawie podanych informacji nie można stwierdzić ile punktów wspólnych ma prosta
\(p\) z
hiperbolą.}{Prosta  \(p\) 
ma dwa punkty wspólne z hiperbolą.}{Prosta \(p\) 
ma dokładnie jeden punkt wspólny z hiperbolą.}{Prosta \(p\) 
nie ma punktów wspólnych z hiperbolą.}{}{}}
\LANGcs{
\Question{
Je dána hyperbola \(\frac{x^{2}} 
{16} -\frac{y^{2}} 
 {4} = 1\) 
a přímka \(p\),
která je rovnoběžná s některou ze souřadnicových os. Pak lze tvrdit, že:}
{Z daných informací není možné jednoznačně určit
počet společných bodů dané hyperboly a přímky
\(p\).}{Přímka \(p\) 
má s danou hyperbolou společné právě dva body.}{Přímka \(p\) 
má s danou hyperbolou společný právě jeden bod.}{Přímka \(p\) 
nemá s danou hyperbolou společný žádný bod.}{}{}}
\LANGsk{
\Question{
Je daná hyperbola \(\frac{x^{2}} 
{16} -\frac{y^{2}} 
 {4} = 1\) 
a priamka \(p\),
ktorá je rovnobežná s niektorou zo súradnicových osí. Potom možno tvrdiť, že:}
{Z daných informácií nie je možné jednoznačne určiť
počet spoločných bodov danej hyperboly a priamky
\(p\).}{Priamka \(p\) 
má s danou hyperbolou spoločné práve dva body.}{Priamka \(p\) 
má s danou hyperbolou spoločný práve jeden bod.}{Priamka \(p\) 
nemá s danou hyperbolou spoločný žiadny bod.}{}{}}
\LANGes{
\Question{
Consideremos la hipérbola  
\[ 
 \frac{x^{2}} 
{16} -\frac{y^{2}} 
 {4} = 1
\]
y la recta \(p\) 
paralela a uno de los ejes. Identifica el enunciado verdadero.}
{No podemos averiguar el número de intersecciones de la recta \(p\) con la hipérbola.}{La recta \(p\) 
tiene dos intersecciones con la hipérbola.}{La recta \(p\) 
tiene una única intersección con la hipérbola.}{La recta \(p\) 
no tiene ninguna intersección con la hipérbola.}{}{}}
\End

\Data\ID{9000104805}
\Updated{2021-10-14 06:10:29}
\Level{C}
\SubArea{Conic Sections}
\LANGen{
\Question{
Find the slope of a line through the center of the hyperbola 
\[ 
 \frac{(x - 2)^{2}} 
 {4} -\frac{(y + 3)^{2}} 
 {9} = 1
\]
which has a unique intersection with this hyperbola.}
{There is no solution, the line with these properties does not exist.}{\(\frac{3}
{2}\)}{\(-\frac{3} 
{2}\)}{\(\frac{2}
{3}\)}{\(1\)}{\(0\)}}
\LANGpl{
\Question{
Wskaż współczynnik kierunkowy prostej przechodzącej przez środek hiperboli
\[ 
 \frac{(x - 2)^{2}} 
 {4} -\frac{(y + 3)^{2}} 
 {9} = 1
\]
tak, aby miała dokładnie jeden punkt wspólny z hiperbolą.}
{Rozwiązanie nie istnieje.}{\(\frac{3}
{2}\)}{\(-\frac{3} 
{2}\)}{\(\frac{2}
{3}\)}{\(1\)}{\(0\)}}
\LANGcs{
\Question{
Co platí pro přímku, která prochází středem hyperboly
\(\frac{(x-2)^{2}} 
 {4} -\frac{(y+3)^{2}} 
 {9} = 1\) a
má s ní společný právě jeden bod?}
{Taková přímka neexistuje.}{Směrnice přímky je \(\frac{3}
{2}\).}{Směrnice přímky je \(-\frac{3} 
{2}\).}{Směrnice přímky je \(\frac{2}
{3}\).}{Směrnice přímky je \(1\).}{Směrnice přímky je \(0\).}}
\LANGsk{
\Question{
Čo platí pre priamku, ktorá prechádza stredom hyperboly
\(\frac{(x-2)^{2}} 
 {4} -\frac{(y+3)^{2}} 
 {9} = 1\) a
má s ňou spoločný práve jeden bod?}
{Taká priamka neexistuje.}{Smernica priamky je \(\frac{3}
{2}\).}{Smernica priamky je \(-\frac{3} 
{2}\).}{Smernica priamky je \(\frac{2}
{3}\).}{Smernica priamky je \(1\).}{Smernice priamky je \(0\).}}
\LANGes{
\Question{
Encuentra la pendiente de la recta que pasa por el centro de la hipérbola 
\[ 
 \frac{(x - 2)^{2}} 
 {4} -\frac{(y + 3)^{2}} 
 {9} = 1
\]
y tiene una única intersección con esta hipérbola.}
{No hay solución, la recta dada por estas propiedades no existe.}{\(\frac{3}
{2}\).}{\(-\frac{3} 
{2}\).}{\(\frac{2}
{3}\).}{\(1\).}{\(0\).}}
\End

\Data\ID{9000104809}
\Updated{2021-10-14 06:10:09}
\Level{C}
\SubArea{Conic Sections}
\LANGen{
\Question{
Among the following lines (which all pass through the point
\([-1;3]\))
identify a line which is tangent to the following hyperbola. 
\[ 
 (x + 2)\cdot (y - 2) = 1
\]}
{\(k\colon \ y = -x + 2\)}{\(p\colon \ y = 3\)}{\(q\colon \ x = -1\)}{\(r\colon \ y = x + 4\)}{None of the given answers is correct.}{}}
\LANGpl{
\Question{
Wszystkie proste przechodzą przez punkt 
\([-1;3]\))
wskaż prostą, która jest styczną do hiperboli.
\[ 
 (x + 2)\cdot (y - 2) = 1
\]}
{\(k\colon \ y = -x + 2\)}{\(p\colon \ y = 3\)}{\(q\colon \ x = -1\)}{\(r\colon \ y = x + 4\)}{Brak poprawnej odpowiedzi.}{}}
\LANGcs{
\Question{
Všechny uvedené přímky procházejí bodem
\([-1;3]\). Která z nich je
tečnou hyperboly \((x + 2)\cdot (y - 2) = 1\)?}
{\(k\colon \ y = -x + 2\)}{\(p\colon \ y = 3\)}{\(q\colon \ x = -1\)}{\(r\colon \ y = x + 4\)}{Žádná z předcházejících odpovědí není správná.}{}}
\LANGsk{
\Question{
Všetky uvedené priamky prechádzajú bodom
\([-1;3]\). Ktorá z nich je
dotyčnicou hyperboly \((x + 2)\cdot (y - 2) = 1\)?}
{\(k\colon \ y = -x + 2\)}{\(p\colon \ y = 3\)}{\(q\colon \ x = -1\)}{\(r\colon \ y = x + 4\)}{Žiadna z uvedených priamok nie je dotyčnicou danej hyperboly.}{}}
\LANGes{
\Question{
Entre las siguientes rectas que pasan todas por el punto \([-1;3]\)
Identifica la que es tangente a la siguiente hipérbola.  
\[ 
 (x + 2)\cdot (y - 2) = 1
\]}
{\(k\colon \ y = -x + 2\)}{\(p\colon \ y = 3\)}{\(q\colon \ x = -1\)}{\(r\colon \ y = x + 4\)}{Ninguna de las respuestas dadas es correcta.}{}}
\End

\Data\ID{9000106901}
\Updated{2021-10-16 11:10:10}
\Level{C}
\SubArea{Conic Sections}
\LANGen{
\Question{
A body is thrown at the initial angle \(\alpha = 45^{\circ }\) 
and the initial velocity \(v_{0} = 10\, \mathrm{m}/\mathrm{s}\).
The trajectory of the body is a parabola. Find the equation of this parabola. Hint: The coordinates of the moving body as functions of time are 
\[ 
\begin{aligned}x& = v_{0}t\cdot \cos \alpha , &
\\y& = v_{0}t\cdot \sin \alpha -\frac{1} 
{2}gt^{2}.
\\ \end{aligned}
\]
Consider the standard acceleration due to gravity
\(g = 10\, \mathrm{m}/\mathrm{s}^{2}\).}
{\((x - 5)^{2} = -10\cdot (y - 2.5)\)}{\((x - 5)^{2} = 10\cdot (y + 2.5)\)}{\(x^{2} = -10\cdot (y - 5)\)}{\((x - 5)^{2} = -10\cdot (y + 2.5)\)}{}{}}
\LANGpl{
\Question{
Ciało rzucone pod kątem \(\alpha = 45^{\circ }\) względem powierzchni Ziemi
z prędkością początkową \(v_{0} = 10\, \mathrm{m}/\mathrm{s}\) porusza się po torze parabolicznym
opisanym równaniami parametrycznymi:
\[ 
\begin{aligned}x& = v_{0}t\cdot \cos \alpha , &
\\y& = v_{0}t\cdot \sin \alpha -\frac{1} 
{2}gt^{2}.
\\ \end{aligned}
\]
Standardowe przyspieszenie ziemskie wynosi
\(g = 10\, \mathrm{m}/\mathrm{s}^{2}\).
Wskaż równanie paraboli.}
{\((x - 5)^{2} = -10\cdot (y - 2.5)\)}{\((x - 5)^{2} = 10\cdot (y + 2.5)\)}{\(x^{2} = -10\cdot (y - 5)\)}{\((x - 5)^{2} = -10\cdot (y + 2.5)\)}{}{}}
\LANGcs{
\Question{
Okamžitá poloha šikmo vzhůru vrženého tělesa je v homogenním
gravitačním poli Země popsána rovnicemi:
\[\begin{aligned}
 x & = v_{0}t\cdot \cos \alpha , & &
 \\y & = v_{0}t\cdot \sin \alpha -\frac{1} 
{2}gt^{2}. & &
\end{aligned}\]
V případě, že pohyb není brzděn odporovými silami, je
jeho trajektorií část paraboly. Určete rovnici paraboly, po
jejíž části se pohybuje těleso, které je vrženo pod úhlem
\(\alpha = 45^{\circ }\) počáteční rychlostí
\(v_{0} = 10\, \mathrm{m}/\mathrm{s}\). Tíhové zrychlení
zaokrouhlete na hodnotu \(g = 10\, \mathrm{m}/\mathrm{s}^{2}\).}
{\((x - 5)^{2} = -10\cdot (y - 2{,}5)\)}{\((x - 5)^{2} = 10\cdot (y + 2{,}5)\)}{\(x^{2} = -10\cdot (y - 5)\)}{\((x - 5)^{2} = -10\cdot (y + 2{,}5)\)}{}{}}
\LANGsk{
\Question{
Okamžitá poloha šikmo hore vrhnutého telesa je v homogénnom
gravitačnom poli Zeme opísaná rovnicami:
\[\begin{aligned}
 x & = v_{0}t\cdot \cos \alpha , & &
 \\y & = v_{0}t\cdot \sin \alpha -\frac{1} 
{2}gt^{2}. & &
\end{aligned}\]
V prípade, že pohyb nieje brzdený odporovými silami, je
jeho trajektóriou časť paraboly. Určte rovnicu paraboly, po
ktorých častiach sa pohybuje teleso, ktoré je vrhnuté pod uhlom
\(\alpha = 45^{\circ }\) počiatočnou rýchlosťou
\(v_{0} = 10\, \mathrm{m}/\mathrm{s}\). Ťahové zrýchlenie
zaokrúhlite na hodnotu \(g = 10\, \mathrm{m}/\mathrm{s}^{2}\).}
{\((x - 5)^{2} = -10\cdot (y - 2{,}5)\)}{\((x - 5)^{2} = 10\cdot (y + 2{,}5)\)}{\(x^{2} = -10\cdot (y - 5)\)}{\((x - 5)^{2} = -10\cdot (y + 2{,}5)\)}{}{}}
\LANGes{
\Question{
Un cuerpo lanzado con un movimiento parabólico tiene un ángulo inicial de \(\alpha = 45^{\circ }\) 
y una velocidad inicial de \(v_{0} = 10\, \mathrm{m}/\mathrm{s}\).
Encuentra la ecuación de la parábola que describe su movimiento. Pista: Las coordenadas de un cuerpo que se mueve en el campo gravitatorio son:  
\[ 
\begin{aligned}x& = v_{0}t\cdot \cos \alpha , &
\\y& = v_{0}t\cdot \sin \alpha -\frac{1} 
{2}gt^{2}.
\\ \end{aligned}
\]
Consideremos la gravedad de la Tierra \(g = 10\, \mathrm{m}/\mathrm{s}^{2}\).}
{\((x - 5)^{2} = -10\cdot (y - 2.5)\)}{\((x - 5)^{2} = 10\cdot (y + 2.5)\)}{\(x^{2} = -10\cdot (y - 5)\)}{\((x - 5)^{2} = -10\cdot (y + 2.5)\)}{}{}}
\End

\Data\ID{9000106902}
\Updated{2021-10-28 09:10:19}
\Level{C}
\SubArea{Conic Sections}
\LANGen{
\Question{
Consider a planet traveling around the Sun on an elliptic trajectory. In the perihelion (the
point where the planet is nearest to the Sun) is the distance from the planet to the Sun
\(4.5\, \mathrm{AU}\). The excentricity
of the ellipse is \(0.5\, \mathrm{AU}\).
Find the equation for the trajectory of the planet. Use the coordinate system with center in the
Sun and \(x\)-axis
along the major axis of the ellipse.}
{\(\frac{(x-0.5)^{2}} 
 {25} + \frac{y^{2}} 
{24.75} = 1\)}{\(\frac{x^{2}} 
{25} + \frac{(y-0.5)^{2}} 
 {24.75} = 1\)}{\(\frac{x^{2}} 
{25} + \frac{y^{2}} 
{24.75} = 1\)}{\(\frac{(x-0.5)^{2}} 
 {24.75} + \frac{y^{2}} 
{25} = 1\)}{}{}}
\LANGpl{
\Question{
Planeta krąży wokół Słońca po orbicie eliptycznej.  W peryhelium 
(punkt, w którym planeta znajduje się najbliżej Słońca) odległość do Słońca wynosi
\(4.5\, \mathrm{AU}\). Mimośród elipsy
jest równy \(0.5\, \mathrm{AU}\).
Wskaż równanie przedstawiające tor planety.  Centrum układu współrzędnych stanowi Słońce a oś \(x\),   leży wzdłuż głównej  osi elipsy.}
{\(\frac{(x-0.5)^{2}} 
 {25} + \frac{y^{2}} 
{24.75} = 1\)}{\(\frac{x^{2}} 
{25} + \frac{(y-0.5)^{2}} 
 {24.75} = 1\)}{\(\frac{x^{2}} 
{25} + \frac{y^{2}} 
{24.75} = 1\)}{\(\frac{(x-0.5)^{2}} 
 {24.75} + \frac{y^{2}} 
{25} = 1\)}{}{}}
\LANGcs{
\Question{
Planetka obíhá kolem Slunce po eliptické trajektorii, přičemž vzdálenost v
perihéliu je \(4{,}5\) 
AU (AU je tzv. astronomická jednotka, perihélium je místo, v němž
má planetka minimální vzdálenost od Slunce) a excentricita elipsy je
\(0{,}5\) 
AU. Určete, která z nabídnutých rovnic vyjadřuje tuto elipsu
v soustavě souřadnic, v jejímž středu bude Slunce a osa
„\(x\) ” bude
určena hlavní osou elipsy.}
{\(\frac{(x-0{,}5)^{2}} 
 {25} + \frac{y^{2}} 
{24{,}75} = 1\)}{\(\frac{x^{2}} 
{25} + \frac{(y-0{,}5)^{2}} 
 {24{,}75} = 1\)}{\(\frac{x^{2}} 
{25} + \frac{y^{2}} 
{24{,}75} = 1\)}{\(\frac{(x-0{,}5)^{2}} 
 {24{,}75} + \frac{y^{2}} 
{25} = 1\)}{}{}}
\LANGsk{
\Question{
Planétka obieha okolo Slnka po eliptické trajektórii, pričom vzdialenosť v
perihéliu je \(4{,}5\) 
AU (AU je tzv. astronomická jednotka, perihélium je miesto, v ktorom
má planétka minimálnu vzdialenosť od Slnka) a excentricita elipsy je
\(0{,}5\) 
AU. Určte, ktorá z ponúknutých rovníc vyjadruje túto elipsu
v sústave súradníc, v jej strede bude Slnko a os
„\(x\) ” bude
určená hlavnou osou elipsy.}
{\(\frac{(x-0{,}5)^{2}} 
 {25} + \frac{y^{2}} 
{24{,}75} = 1\)}{\(\frac{x^{2}} 
{25} + \frac{(y-0{,}5)^{2}} 
 {24{,}75} = 1\)}{\(\frac{x^{2}} 
{25} + \frac{y^{2}} 
{24{,}75} = 1\)}{\(\frac{(x-0{,}5)^{2}} 
 {24{,}75} + \frac{y^{2}} 
{25} = 1\)}{}{}}
\LANGes{
\Question{
Consideremos un planeta moviéndose alrededor del Sol en una órbita elíptica. En el perihelio (punto de la órbita más próximo al Sol) la distancia del planeta al Sol es \(4.5\, \mathrm{AU}\). La excentricidad de la elipse es \(0.5\, \mathrm{AU}\).
Halla la ecuación de la trayectoria de este planeta. Usa el sistema de coordenadas cuyo centro es el Sol y el eje \(x\) es el eje mayor de la elipse.}
{\(\frac{(x-0.5)^{2}} 
 {25} + \frac{y^{2}} 
{24.75} = 1\)}{\(\frac{x^{2}} 
{25} + \frac{(y-0.5)^{2}} 
 {24.75} = 1\)}{\(\frac{x^{2}} 
{25} + \frac{y^{2}} 
{24.75} = 1\)}{\(\frac{(x-0.5)^{2}} 
 {24.75} + \frac{y^{2}} 
{25} = 1\)}{}{}}
\End

\Data\ID{9000106903}
\Updated{2021-10-11 11:10:44}
\Level{C}
\SubArea{Conic Sections}
\LANGen{
\Question{
The motion with a constant acceleration is described by the relation
\(s = \frac{1} 
{2}at^{2}\).
Consequently, the graph which shows the distance as a function of
time is part of a parabola. Find the directrix of this parabola, if
\(a = 4\, \mathrm{m}/\mathrm{s}^{2}\).}
{\(s = -\frac{1}
{8}\)}{\(s = -1\)}{\(s = \frac{1} 
{8}\)}{\(s = 1\)}{}{}}
\LANGpl{
\Question{
Ruch o stałym przyspieszeniu jest określony zależnością
\(s = \frac{1} 
{2}at^{2}\).
Wykres wskazuje drogę w funkcji czasu, która jest częścią paraboli.  Wyznacz kierownicę paraboli, jeśli
\(a = 4\, \mathrm{m}/\mathrm{s}^{2}\).}
{\(s = -\frac{1}
{8}\)}{\(s = -1\)}{\(s = \frac{1} 
{8}\)}{\(s = 1\)}{}{}}
\LANGcs{
\Question{
Grafem funkční závislosti dráhy na čase rovnoměrně
zrychleného pohybu je část paraboly. Funkce je určena rovnicí
\(s = \frac{1} 
{2}at^{2}\). Určete
rovnici řídící přímky paraboly, jestliže se těleso začalo pohybovat v čase
\(t = 0\, \mathrm{s}\) a pohybuje se
se zrychlením \(a = 4\, \mathrm{m}/\mathrm{s}^{2}\).}
{\(s = -\frac{1}
{8}\)}{\(s = -1\)}{\(s = \frac{1} 
{8}\)}{\(s = 1\)}{}{}}
\LANGsk{
\Question{
Grafom funkčnej závislosti dráhy na čase rovnomerne
zrýchleného pohybu je časť paraboly. Funkcia je určená rovnicou
\(s = \frac{1} 
{2}at^{2}\). Určte
rovnicu riadiacej priamky paraboly, ak sa teleso začalo pohybovať v čase
\(t = 0\, \mathrm{s}\) a pohybuje sa
so zrýchlením \(a = 4\, \mathrm{m}/\mathrm{s}^{2}\).}
{\(s = -\frac{1}
{8}\)}{\(s = -1\)}{\(s = \frac{1} 
{8}\)}{\(s = 1\)}{}{}}
\LANGes{
\Question{
El Movimiento Uniformemente Acelerado se describe con la siguiente ecuación \(s = \frac{1} 
{2}at^{2}\) y el gráfico que representa esta relación es una parábola. Encuentra la directriz de esta parábola, si la aceleración es
\(a = 4\, \mathrm{m}/\mathrm{s}^{2}\) y el tiempo inicial \(t = 0\, \mathrm{s}\).}
{\(s = -\frac{1}
{8}\)}{\(s = -1\)}{\(s = \frac{1} 
{8}\)}{\(s = 1\)}{}{}}
\End

\Data\ID{9000106904}
\Updated{2021-10-28 09:10:02}
\Level{C}
\SubArea{Conic Sections}
\LANGen{
\Question{
The motion with a constant deceleration is described by the relation 
\[ 
 s = v_{0}t -\frac{1} 
{2}at^{2}.
\]
Consequently, the graph which shows the distance as a function
of time is part of a parabola. Find the focus of this parabola, if
\(v_{0} = 16\, \mathrm{m}/\mathrm{s}\) and
\(a = 4\, \mathrm{m}/\mathrm{s}^{2}\).}
{\([4;\ 31.875]\)}{\([8;\ 31.875]\)}{\([4;\ 63.5]\)}{\([8;\ 63.5]\)}{}{}}
\LANGpl{
\Question{
Ruch o stałym opóźnieniu  jest określony zależnością 
\[ 
 s = v_{0}t -\frac{1} 
{2}at^{2}.
\]
Wykres wskazuje drogę w funkcji czasu, która jest częścią paraboli.  Wyznacz ognisko paraboli, jeśli
\(v_{0} = 16\, \mathrm{m}/\mathrm{s}\) and
\(a = 4\, \mathrm{m}/\mathrm{s}^{2}\).}
{\([4;\ 31.875]\)}{\([8;\ 31.875]\)}{\([4;\ 63.5]\)}{\([8;\ 63.5]\)}{}{}}
\LANGcs{
\Question{
Grafem funkční závislosti dráhy na čase rovnoměrně
zpomaleného pohybu je část paraboly. Funkce je určena rovnicí
\(s = v_{0}t -\frac{1} 
{2}at^{2}\). Určete
souřadnice ohniska této paraboly, jestliže těleso začalo zpomalovat v čase
\(t = 0\, \mathrm{s}\) a počáteční
rychlost tělesa byla \(v_{0} = 16\, \mathrm{m}/\mathrm{s}\).
Zpomalení má hodnotu \(a = 4\, \mathrm{m}/\mathrm{s}^{2}\).}
{\([4;\ 31{,}875]\)}{\([8;\ 31{,}875]\)}{\([4;\ 63{,}5]\)}{\([8;\ 63{,}5]\)}{}{}}
\LANGsk{
\Question{
Grafom funkčnej závislosti dráhy na čase rovnomerne
spomaleného pohybu je časť paraboly. Funkcia je určená rovnicou
\(s = v_{0}t -\frac{1} 
{2}at^{2}\). Určte
súradnice ohniska tejto paraboly, ak teleso začalo spomaľovať v čase
\(t = 0\, \mathrm{s}\) a počiatočná
rýchlosť telesa bola \(v_{0} = 16\, \mathrm{m}/\mathrm{s}\).
Spomalenie má hodnotu \(a = 4\, \mathrm{m}/\mathrm{s}^{2}\).}
{\([4;\ 31{,}875]\)}{\([8;\ 31{,}875]\)}{\([4;\ 63{,}5]\)}{\([8;\ 63{,}5]\)}{}{}}
\LANGes{
\Question{
El Movimiento Rectilíneo Uniformemente Acelerado se describe por la siguiente ecuación \[ 
 s = v_{0}t -\frac{1} 
{2}at^{2}.
\]
La gráfica que representa la distancia en función del tiempo es parte de una parábola. Halla el foco de esta parábola, si \(v_{0} = 16\, \mathrm{m}/\mathrm{s}\) y
\(a = 4\, \mathrm{m}/\mathrm{s}^{2}\).}
{\([4;\ 31.875]\)}{\([8;\ 31.875]\)}{\([4;\ 63.5]\)}{\([8;\ 63.5]\)}{}{}}
\End

\Data\ID{9000106905}
\Updated{2021-10-28 09:10:28}
\Level{C}
\SubArea{Conic Sections}
\LANGen{
\Question{
The motion with a constant deceleration is described by the relation 
\[ 
 s = v_{0}t -\frac{1} 
{2}at^{2}.
\]
Consequently, the graph which shows the distance as a function of time
is part of a parabola. Find the vertex equation of this parabola, if
\(v_{0} = 8\, \mathrm{m}/\mathrm{s}\) and
\(a = 4\, \mathrm{m}/\mathrm{s}^{2}\).}
{\(-\frac{1} 
{2}(s - 8) = (t - 2)^{2}\)}{\(\frac{1}
{2}(s + 4) = (t + 2)^{2}\)}{\(2(s + 8) = (t + 2)^{2}\)}{\(- 2(s + 4) = (t + 2)^{2}\)}{}{}}
\LANGpl{
\Question{
Ruch o stałym opóźnieniu jest określony zależnością 
\[ 
 s = v_{0}t -\frac{1} 
{2}at^{2}.
\]
Wykres wskazuje drogę w funkcji czasu, która jest częścią paraboli.  Wyznacz równanie wierzchołka paraboli, jeśli
\(v_{0} = 8\, \mathrm{m}/\mathrm{s}\) i
\(a = 4\, \mathrm{m}/\mathrm{s}^{2}\).}
{\(-\frac{1} 
{2}(s - 8) = (t - 2)^{2}\)}{\(\frac{1}
{2}(s + 4) = (t + 2)^{2}\)}{\(2(s + 8) = (t + 2)^{2}\)}{\(- 2(s + 4) = (t + 2)^{2}\)}{}{}}
\LANGcs{
\Question{
Grafem funkční závislosti dráhy na čase rovnoměrně
zpomaleného pohybu je část paraboly. Funkce je určena rovnicí
\(s = v_{0}t -\frac{1} 
{2}at^{2}\).
Určete vrcholovou rovnici této paraboly, jestliže je v čase
\(t = 0\, \mathrm{s}\) počáteční
rychlost tělesa \(v_{0} = 8\, \mathrm{m}/\mathrm{s}\) 
a zrychlení \(a = 4\, \mathrm{m}/\mathrm{s}^{2}\).}
{\(-\frac{1} 
{2}(s - 8) = (t - 2)^{2}\)}{\(\frac{1}
{2}(s + 4) = (t + 2)^{2}\)}{\(2(s + 8) = (t + 2)^{2}\)}{\(- 2(s + 4) = (t + 2)^{2}\)}{}{}}
\LANGsk{
\Question{
Grafom funkčnej závislosti dráhy na čase rovnomerne
spomaleného pohybu je časť paraboly. Funkcia je určená rovnicou
\(s = v_{0}t -\frac{1} 
{2}at^{2}\).
Určte vrcholovú rovnicu tejto paraboly, ak je v čase
\(t = 0\, \mathrm{s}\) počiatočná
rýchlosť telesa \(v_{0} = 8\, \mathrm{m}/\mathrm{s}\) 
a zrýchlenie \(a = 4\, \mathrm{m}/\mathrm{s}^{2}\).}
{\(-\frac{1} 
{2}(s - 8) = (t - 2)^{2}\)}{\(\frac{1}
{2}(s + 4) = (t + 2)^{2}\)}{\(2(s + 8) = (t + 2)^{2}\)}{\(- 2(s + 4) = (t + 2)^{2}\)}{}{}}
\LANGes{
\Question{
El Movimiento Rectilíneo Uniformemente Acelerado se describe por la siguiente ecuación
\[ 
 s = v_{0}t -\frac{1} 
{2}at^{2}.
\]
El gráfico que representa la distancia en función del tiempo es parte de una parábola. Halla la ecuación del vértice de esta parábola, si \(v_{0} = 8\, \mathrm{m}/\mathrm{s}\) y
\(a = 4\, \mathrm{m}/\mathrm{s}^{2}\).}
{\(-\frac{1} 
{2}(s - 8) = (t - 2)^{2}\)}{\(\frac{1}
{2}(s + 4) = (t + 2)^{2}\)}{\(2(s + 8) = (t + 2)^{2}\)}{\(- 2(s + 4) = (t + 2)^{2}\)}{}{}}
\End

\Data\ID{9000117701}
\Updated{2021-10-28 09:10:45}
\Level{C}
\SubArea{Conic Sections}
\LANGen{
\Question{
A body is thrown at the initial angle \(\alpha = 30^{\circ }\) 
and the initial velocity \(v_{0} = 20\, \mathrm{m}/\mathrm{s}\).
The trajectory of the body is a part of parabola. Find the directrix of this
parabola. Hint: The coordinates of the moving body as functions of time are 
\[ 
\begin{aligned}x& = v_{0}t\cdot \cos \alpha , &
\\y& = v_{0}t\cdot \sin \alpha -\frac{1} 
{2}gt^{2}.
\\ \end{aligned}
\]
Consider the standard acceleration due to gravity
\(g = 10\, \mathrm{m}/\mathrm{s}^{2}\).}
{\(y = 20\)}{\(y = 5\)}{\(y = 15\)}{\(y = 10\)}{}{}}
\LANGpl{
\Question{
Ciało rzucone pod kątem \(\alpha = 30^{\circ }\) 
 względem powierzchni ziemi z prędkością początkową  \(v_{0} = 20\, \mathrm{m}/\mathrm{s}\).
porusza się w próżni po torze parabolicznym, opisanym równaniami parametrycznym. 
\[ 
\begin{aligned}x& = v_{0}t\cdot \cos \alpha , &
\\y& = v_{0}t\cdot \sin \alpha -\frac{1} 
{2}gt^{2}.
\\ \end{aligned}
\]
Wskaż kierownicę paraboli.
Standardowe przyspieszenie ziemskie wynosi
\(g = 10\, \mathrm{m}/\mathrm{s}^{2}\).}
{\(y = 20\)}{\(y = 5\)}{\(y = 15\)}{\(y = 10\)}{}{}}
\LANGcs{
\Question{
Těleso vržené šikmo vzhůru pod úhlem
\(\alpha = 30^{\circ }\) počáteční
rychlostí o velikosti \(v_{0} = 20\, \mathrm{m}/\mathrm{s}\) 
opisuje při svém pohybu část paraboly. Určete rovnici řídící
přímky této paraboly. (Okamžitá poloha šikmo vzhůru vrženého
tělesa je v homogenním gravitačním poli Země popsána rovnicemi:
\(x = v_{0}t\cdot \cos \alpha \),
\(y = v_{0}t\cdot \sin \alpha -\frac{1} 
{2}gt^{2}\). Tíhové zrychlení
zaokrouhlete na hodnotu \(g = 10\, \mathrm{m}/\mathrm{s}^{2}\)).}
{\(y = 20\)}{\(y = 5\)}{\(y = 15\)}{\(y = 10\)}{}{}}
\LANGsk{
\Question{
Teleso vrhnuté šikmo hore pod uhlom
\(\alpha = 30^{\circ }\) začiatočná
rýchlosť o veľkosti \(v_{0} = 20\, \mathrm{m}/\mathrm{s}\) 
opisuje pri svojom pohybe časť paraboly. Určte rovnicu riadiacej
priamky tejto paraboly. (Okamžitá poloha šikmo hore hodeného
telesa je v homogénnom gravitačnom poli Zeme popísaná rovnicami:
\(x = v_{0}t\cdot \cos \alpha \),
\(y = v_{0}t\cdot \sin \alpha -\frac{1} 
{2}gt^{2}\). Tiažové zrýchlenie
zaokrúhlite na hodnotu \(g = 10\, \mathrm{m}/\mathrm{s}^{2}\)).}
{\(y = 20\)}{\(y = 5\)}{\(y = 15\)}{\(y = 10\)}{}{}}
\LANGes{
\Question{
Un cuerpo se lanza en movimiento parabólico con un ángulo inicial de \(\alpha = 30^{\circ }\) y una velocidad inicial de \(v_{0} = 20\, \mathrm{m}/\mathrm{s}\). Encuentra la directriz de esta parábola. Pista: Las coordenadas del cuerpo que se mueve en el campo gravitatorio son:
\[ 
\begin{aligned}x& = v_{0}t\cdot \cos \alpha , &
\\y& = v_{0}t\cdot \sin \alpha -\frac{1} 
{2}gt^{2}.
\\ \end{aligned}
\]
Consideremos la gravedad de la Tierra \(g = 10\, \mathrm{m}/\mathrm{s}^{2}\).}
{\(y = 20\)}{\(y = 5\)}{\(y = 15\)}{\(y = 10\)}{}{}}
\End

\Data\ID{9000117702}
\Updated{2021-10-28 09:10:52}
\Level{C}
\SubArea{Conic Sections}
\LANGen{
\Question{
The Earth travels around the Sun on an elliptical orbit. The Sun is in the
focus of this ellipse. The maximal distance from Earth to the Sun is
\(152.1\cdot 10^{6}\, \mathrm{km}\), the minimal distance
from Earth to the Sun is \(147.1\cdot 10^{6}\, \mathrm{km}\).
Find the length of the semi-minor axis (one half of the length
of the shorter axis) and round your answer to the nearest
\(10^{4}\, \mathrm{km}\).}
{\(149.58\cdot 10^{6}\, \mathrm{km}\)}{\(2.58\cdot 10^{6}\, \mathrm{km}\)}{\(299.21\cdot 10^{6}\, \mathrm{km}\)}{\(149.61\cdot 10^{6}\, \mathrm{km}\)}{}{}}
\LANGpl{
\Question{
Ziemia porusza się wokół Słońca po orbicie eliptycznej. Słońce stanowi ognisko tej elipsy.
Maksymalna odległość z Ziemi do Słońca to
\(152.1\cdot 10^{6}\, \mathrm{km}\), minimalna odległość
z Ziemi do Słońca to \(147.1\cdot 10^{6}\, \mathrm{km}\).
Wskaż długość małej półosi  (połowa długości krótszej osi
) i zaokrągli odpowiedź do pełnych
\(10^{4}\, \mathrm{km}\).}
{\(149.58\cdot 10^{6}\, \mathrm{km}\)}{\(2.58\cdot 10^{6}\, \mathrm{km}\)}{\(299.21\cdot 10^{6}\, \mathrm{km}\)}{\(149.61\cdot 10^{6}\, \mathrm{km}\)}{}{}}
\LANGcs{
\Question{
Země se pohybuje kolem Slunce po eliptické trajektorii, přičemž Slunce
je v ohnisku této elipsy. Jaká je velikost vedlejší poloosy, jestliže
víme, že maximální vzdálenost Země od Slunce (tzv. afélium) je
\(152{,}1\cdot 10^{6}\, \mathrm{km}\) 
a minimální vzdálenost Země od Slunce (tzv. perihélium) je
\(147{,}1\cdot 10^{6}\, \mathrm{km}\).
(Výsledek zaokrouhlete na desetitisíce km.)}
{\(149{,}58\cdot 10^{6}\, \mathrm{km}\)}{\(2{,}58\cdot 10^{6}\, \mathrm{km}\)}{\(299{,}21\cdot 10^{6}\, \mathrm{km}\)}{\(149{,}61\cdot 10^{6}\, \mathrm{km}\)}{}{}}
\LANGsk{
\Question{
Zem sa pohybuje okolo Slnka po eliptickej trajektórií, pričom Slnko
je v ohnisku tejto elipsy. Aká je veľkosť vedľajšej polosi, ak
vieme, že maximálna vzdialenosť Zeme od Slnka (tzv. afélium) je
\(152{,}1\cdot 10^{6}\, \mathrm{km}\) 
a minimálna vzdialenosť Zeme od Slnka (tzv. perihélium) je
\(147{,}1\cdot 10^{6}\, \mathrm{km}\).
(Výsledok zaokrúhlite na desaťtisíce km.)}
{\(149{,}58\cdot 10^{6}\, \mathrm{km}\)}{\(2{,}58\cdot 10^{6}\, \mathrm{km}\)}{\(299{,}21\cdot 10^{6}\, \mathrm{km}\)}{\(149{,}61\cdot 10^{6}\, \mathrm{km}\)}{}{}}
\LANGes{
\Question{
La Tierra se mueve alrededor del Sol en una órbita elíptica. El Sol está en uno de los  focos de la elipse. La distancia máxima de la Tierra al Sol es \(152.1\cdot 10^{6}\, \mathrm{km}\), la distancia mínima de la Tierra al Sol es \(147.1\cdot 10^{6}\, \mathrm{km}\).
Encuentra la longitud del semieje menor (la mitad de la longitud del eje más corto) y redondea tu respuesta a \(10^{4}\, \mathrm{km}\).}
{\(149.58\cdot 10^{6}\, \mathrm{km}\)}{\(2.58\cdot 10^{6}\, \mathrm{km}\)}{\(299.21\cdot 10^{6}\, \mathrm{km}\)}{\(149.61\cdot 10^{6}\, \mathrm{km}\)}{}{}}
\End

\Data\ID{9000117703}
\Updated{2021-10-28 09:10:30}
\Level{C}
\SubArea{Conic Sections}
\LANGen{
\Question{
For an isothermal process in an ideal gas the product
\(pV \) 
is constant (Boyle's law). In a pressure-volume diagram which shows
\(p\) as a
function of \(V \) 
this law describes a hyperbola (called isotherm). Do we have enough information to
identify the asymptotes? If so, find these asymptotes.}
{\(p = 0\),
\(V = 0\)}{\(p = V \),
\(p = -V \)}{\(p = 0\),
\(p = V \)}{It is not possible to draw any conclusion.}{}{}}
\LANGpl{
\Question{
W przemianie izotermicznej gazu doskonałego iloczyn 
\(pV \) 
jest wartością stałą ( prawo Boylea). Hiperbola przedstawia zależność ciśnienia \(p\) od objętości gazu
  \(V \) 
 (zwana izotermą ). Czy mamy wystarczającą ilość danych, aby wskazać asymptoty hiperboli ? Jeśli tak, to wskaż te asymptoty?}
{\(p = 0\),
\(V = 0\)}{\(p = V \),
\(p = -V \)}{\(p = 0\),
\(p = V \)}{Brak rozwiązania.}{}{}}
\LANGcs{
\Question{
Tzv. „izotermický děj” s ideálním plynem můžeme popsat rovnicí
\(pV = \mathrm{konst.}\), kde
\(p\) je tlak ideálního
plynu, \(V \) je
jeho objem. Graf funkční závislosti tlaku ideálního plynu stálé hmotnosti
na jeho objemu při konstantní teplotě se nazývá izoterma. Izoterma je část
hyperboly. Je-li to na základě výše uvedených informací možné, napište
rovnice asymptot této hyperboly. V opačném případě označte, že
asymptoty nelze určit.}
{\(p = 0,\ V = 0\)}{\(p = V,\ p = -V \)}{\(p = 0,\ p = V \)}{Rovnice asymptot jsou závislé na číselném
určení „konstanty”, takže asymptoty není možné
určit rovnicemi.}{}{}}
\LANGsk{
\Question{
Tzv. „izotermický dej” s ideálnym plynom môžeme popísať rovnicou
\(pV = \mathrm{konst.}\), kde
\(p\) tlak ideálneho
plynu, \(V \) je
jeho objem. Graf funkčnej závislosti tlaku ideálneho plynu stálej hmotnosti
na jeho objemu pri konštantnej teplote sa nazýva izoterma. Izoterma je časť
hyperboly. Ak je to na základe vyššie uvedených informácií možné, napíšte
rovnice asymptot tejto hyperboly. V opačnom prípade označte, že
asymptoty nie je možné určiť.}
{\(p = 0,\ V = 0\)}{\(p = V,\ p = -V \)}{\(p = 0,\ p = V \)}{Rovnice asymptot sú závislé na číselnom
určení „konštanty”, takže asymptoty nie je možné
určiť rovnicami.}{}{}}
\LANGes{
\Question{
Un proceso isotérmico con  gas ideal se puede describir por la siguiente ecuación \(pV = \mathrm{konst.}\) donde \(p\) es la presión del gas ideal y \(V \) su volumen. Una curva isoterma es una línea que sobre un diagrama representa los valores sucesivos de las diversas variables de un sistema en un proceso isotermo. Las isotermas forman parte de una hipérbola. Decide si tenemos suficiente información para poder definir las asíntontas de esta hipérbola. Si es así, elige la respuesta correcta.}
{\(p = 0\),
\(V = 0\)}{\(p = V \),
\(p = -V \)}{\(p = 0\),
\(p = V \)}{No es posible encontrar las asíntotas.}{}{}}
\End

\Data\ID{9000117704}
\Updated{2021-10-28 09:10:41}
\Level{C}
\SubArea{Conic Sections}
\LANGen{
\Question{
Given physical quantities and laws relating these quantities, identify an answer
where the graph which relates these quantities is a part of a hyperbola. (The other
quantities are supposed to be constant.)}
{The pressure (\(p\)) and the
area (\(S\)) over which the
pressure is distributed, if \(F = p\cdot S\).}{The mass (\(m\)) and
the kinetic energy (\(E_{k}\))
of a moving body, if \(E_{k} = \frac{1} 
{2}\cdot m\cdot v^{2}\).}{The velocity (\(v\)) and
the kinetic energy (\(E_{k}\))
of a moving body, if \(E_{k} = \frac{1} 
{2}\cdot m\cdot v^{2}\).}{The mass (\(m\)) and the potential
energy (\(E_{p}\)) in a homogeneous
gravitational field, if \(E_{p} = m\cdot g\cdot h\).}{}{}}
\LANGpl{
\Question{
Biorąc pod uwagę dane wielkości fizyczne i prawa określające te wielkości, wskaż odpowiedź 
tak, aby wykres przedstawiający podane wielkości był częścią hiperboli. (Pozostałe wielkości przyjęto jako stałe.)}
{Ciśnienie  (\(p\)) działające na powierzchnię
 (\(S\)), jeśli  \(F = p\cdot S\).}{Masa  (\(m\)) i
energia kinetyczna  (\(E_{k}\))
poruszającego się ciała, jeśli \(E_{k} = \frac{1} 
{2}\cdot m\cdot v^{2}\).}{Prędkość  (\(v\)) i
energia kinetyczna (\(E_{k}\))
poruszającego się ciała, jeśli \(E_{k} = \frac{1} 
{2}\cdot m\cdot v^{2}\).}{Masa  (\(m\)) i energia potencjalna
 (\(E_{p}\)) w jednorodnym polu grawitacyjnym   \(E_{p} = m\cdot g\cdot h\).}{}{}}
\LANGcs{
\Question{
Z nabídnutých možností vyberte tu dvojici fyzikálních veličin, jejichž
graficky vyjádřená závislost tvoří část hyperboly. (Zbývající veličiny
ve vztazích považujeme za konstantní).}
{Tlak (\(p\)) a plocha
(\(S\)), na kterou působí
tlaková síla, je-li \(F = p\cdot S\).}{Hmotnost (\(m\)) a
kinetická energie (\(E_{k}\))
tělesa, je-li \(E_{k} = \frac{1} 
{2}\cdot m\cdot v^{2}\).}{Rychlost (\(v\)) a
kinetická energie (\(E_{k}\))
tělesa, je-li \(E_{k} = \frac{1} 
{2}\cdot m\cdot v^{2}\).}{Hmotnost (\(m\)) a
polohová energie (\(E_{p}\)),
je-li \(E_{p} = mgh\).}{}{}}
\LANGsk{
\Question{
Z ponúknutých možností vyberte tú dvojicu fyzikálnych veličín, ktorých
graficky vyjadrená závislosť tvorí časť hyperboly. (Ostávajúce veličiny
vo vzťahoch považujeme za konštantné).}
{Tlak (\(p\)) a plocha
(\(S\)), na ktorú pôsobí
tlaková sila, ak \(F = p\cdot S\).}{Hmotnosť (\(m\)) a
kinetická energia (\(E_{k}\))
telesa, ak \(E_{k} = \frac{1} 
{2}\cdot m\cdot v^{2}\).}{Rýchlosť (\(v\)) a
kinetická energia (\(E_{k}\))
telesa, ak \(E_{k} = \frac{1} 
{2}\cdot m\cdot v^{2}\).}{Hmotnosť (\(m\)) a
polohová energie (\(E_{p}\)),
ak \(E_{p} = mgh\).}{}{}}
\LANGes{
\Question{
Dadas dos magnitudes físicas identifica cuál de las respuestas representa una relación entre las magnitudes como parte de una hipérbola. (Se supone que las otras magnitudes son constantes.)}
{La presión (\(p\)) y la superficie (\(S\)) en la que se ejerce la fuerza de la presión, si \(F = p\cdot S\).}{La masa (\(m\)) y
la energía cinética (\(E_{k}\))
de un sólido, si \(E_{k} = \frac{1} 
{2}\cdot m\cdot v^{2}\).}{La velocidad (\(v\)) y la energía cinética (\(E_{k}\)) de un sólido, si \(E_{k} = \frac{1} 
{2}\cdot m\cdot v^{2}\).}{La masa (\(m\)) y la energía potencial (\(E_{p}\)) en un campo de gravedad homogéneo, si \(E_{p} = m\cdot g\cdot h\).}{}{}}
\End

\Data\ID{9000117705}
\Updated{2021-10-28 09:10:07}
\Level{C}
\SubArea{Conic Sections}
\LANGen{
\Question{
Given physical quantities and laws relating these quantities, identify an answer
where the graph which relates these quantities is a part of a parabola. (The other
quantities are supposed to be constant.)}
{The electrical work (\(W\))
and the current (\(I\)),
if \(W = R\cdot I^{2}\cdot t\).}{The mass (\(m\)) and the
acceleration (\(a\)) of a
moving body, if \(F = m\cdot a\).}{The height (\(h\)) and the
potential energy (\(E_{p}\)),
if \(E_{p} = m\cdot g\cdot h\).}{The electrical work (\(W\))
and the time (\(t\)),
if \(W = R\cdot I^{2}\cdot t\).}{}{}}
\LANGpl{
\Question{
Biorąc pod uwagę dane wielkości fizyczne i prawa określające te wielkości, wskaż odpowiedź tak, aby wykres przedstawiający podane wielkości był częścią paraboli. (Pozostałe wielkości przyjęto jako stałe.)}
{Praca prądu  (\(W\))
i natężenie  (\(I\)),
jeśli \(W = R\cdot I^{2}\cdot t\).}{Masa  (\(m\)) i przyspieszenie
 (\(a\)) poruszającego się ciała, jeśli \(F = m\cdot a\).}{Wysokość  (\(h\)) i energia potencjalna 
 (\(E_{p}\)),
jeśli \(E_{p} = m\cdot g\cdot h\).}{Praca prądu  (\(W\))
i czas  (\(t\)),
jeśli \(W = R\cdot I^{2}\cdot t\).}{}{}}
\LANGcs{
\Question{
Z nabídnutých možností vyberte tu dvojici fyzikálních veličin, jejichž
graficky vyjádřená závislost tvoří část paraboly. (Zbývající veličiny
ve vztahu považujeme za konstantní).}
{Práce elektrických sil (\(W\)) a
velikost elektrického proudu (\(I\)),
je-li \(W = R\cdot I^{2}\cdot t\).}{Hmotnost (\(m\))
a zrychlení (\(a\))
tělesa, je-li \(F = m\cdot a\).}{Výška nad podložkou (\(h\))
a polohová energie (\(E_{p}\)),
je-li \(E_{p} = mgh\).}{Práce elektrických sil (\(W\))
a doba (\(t\)), po kterou
protéká proud, je-li \(W = R\cdot I^{2}\cdot t\).}{}{}}
\LANGsk{
\Question{
Z ponúknutých možností vyberte tú dvojicu fyzikálnych veličín, ktorých
graficky vyjadrená závislosť tvorí časť paraboly. (Ostávajúce veličiny
vo vzťahu považujeme za konštantné).}
{Práca elektrických síl (\(W\)) a
veľkosť elektrického prúdu (\(I\)),
ak \(W = R\cdot I^{2}\cdot t\).}{Hmotnosť (\(m\))
a zrýchlenie (\(a\))
telesa, ak \(F = m\cdot a\).}{Výška nad podložkou (\(h\))
a polohová energia (\(E_{p}\)),
ak \(E_{p} = mgh\).}{Práca elektrických síl (\(W\))
a doba (\(t\)), počas ktorej
preteká prúd, ak \(W = R\cdot I^{2}\cdot t\).}{}{}}
\LANGes{
\Question{
Dadas las magnitudes físicas, identifica cuál de las respuestas representa una relación entre magnitudes como parte de una parábola. (Se supone que las otras magnitudes son constantes.)}
{El trabajo eléctrico (\(W\))
y la corriente eléctrica (\(I\)),
si \(W = R\cdot I^{2}\cdot t\).}{La masa (\(m\)) y la aceleración (\(a\)) de un sólido, si \(F = m\cdot a\).}{La altitud (\(h\)) y la energía potencial (\(E_{p}\)),
si \(E_{p} = m\cdot g\cdot h\).}{El trabajo eléctrico (\(W\))
y el tiempo (\(t\)),
si \(W = R\cdot I^{2}\cdot t\).}{}{}}
\End

\Data\ID{9000117706}
\Updated{2021-10-28 10:10:53}
\Level{C}
\SubArea{Conic Sections}
\LANGen{
\Question{
Satellites travel along approximately circular paths. Consider a satellite in the height
\(h\) 
measured from the Earth surface. Further, consider the coordinate system
with origin on the Earth surface directly below the satellite and the
\(y\)-axis oriented up (to
the satellite). The \(x\)-axis
is perpendicular to \(y\)-axis
and it is in the plane defined by the trajectory of the satellite. Neglect the Earth's
rotation and find the equation which describes the path of the satellite. The Earth
radius is \(R\).}
{\(x^{2} + (y + R)^{2} = (R + h)^{2}\)}{\(x^{2} + y^{2} = (R + h)^{2}\)}{\(x^{2} + (y + R)^{2} = h^{2}\)}{\(x^{2} + y^{2} = h^{2}\)}{}{}}
\LANGpl{
\Question{
Satelity poruszają się w przybliżeniu po torach kołowych. Rozważ  satelitę o wysokości
\(h\) 
mierzonej od powierzchni Ziemi, układ współrzędnych z początkiem na powierzchni 
Ziemi znajdujący się bezpośrednio pod satelitą  oraz oś
\(y\)  ukierunkowaną w stronę satelity. Oś \(x\) jest 
prostopadła do osi \(y\)
oraz znajduje się na płaszczyźnie wyznaczonej przez tor satelity. Wskaż równanie określające tor satelity pomijając krążenie Ziemi. Promień Ziemi jest równy
 \(R\).}
{\(x^{2} + (y + R)^{2} = (R + h)^{2}\)}{\(x^{2} + y^{2} = (R + h)^{2}\)}{\(x^{2} + (y + R)^{2} = h^{2}\)}{\(x^{2} + y^{2} = h^{2}\)}{}{}}
\LANGcs{
\Question{
Pro pohyb těles (družic) v blízkém okolí Země je důležitá tzv. kruhová
rychlost. Tělesa s touto rychlostí se pohybují po kruhové trajektorii,
přičemž Země je ve středu této trajektorie. V blízkosti povrchu
Země se této rychlosti říká „1. kosmická rychlost” a její hodnota je
\(7{,}9\, \mathrm{km}/\mathrm{s}\). Hodnotu kruhové rychlosti
ve výšce \(h\) nad zemským
povrchem určuje vztah: \(v = \sqrt{ \frac{\kappa \cdot M_{Z } }
{R_{Z}+h}}\),
kde \(M_{Z}\) je hmotnost
Země, \(R_{Z}\) je
poloměr Země a \(\kappa \) 
je gravitační konstanta. Vyberte správnou rovnici kruhové
trajektorie družice, která se v okamžiku startu nachází ve výšce
\(h\) nad zemským povrchem
v soustavě, kde osa \(y\) 
spojuje střed Země s místem startu družice a počátek soustavy je na
povrchu Země.}
{\(x^{2} + (y + R_{Z})^{2} = (R_{Z} + h)^{2}\)}{\(x^{2} + y^{2} = (R_{Z} + h)^{2}\)}{\(x^{2} + (y + R_{Z})^{2} = h^{2}\)}{\(x^{2} + y^{2} = h^{2}\)}{}{}}
\LANGsk{
\Question{
Pre pohyb telies (družíc) v blízkom okolí Zeme je dôležitá tzv. kruhová
rýchlosť. Telesá s touto rýchlosťou sa pohybujú po kruhovej trajektórii,
pričom Zem je v strede tejto trajektórie. V blízkosti povrchu
Zeme sa tejto rýchlosti hovorí „1. kozmická rýchlosť” a jej hodnota je
\(7{,}9\, \mathrm{km}/\mathrm{s}\). Hodnotu kruhovej rýchlosti
vo výške \(h\) nad zemským
povrchom určuje vzťah: \(v = \sqrt{ \frac{\kappa \cdot M_{Z } }
{R_{Z}+h}}\),
kde \(M_{Z}\) je hmotnosť
Zeme, \(R_{Z}\) je
polomer Zeme a \(\kappa \) 
je gravitačná konštanta. Vyberte správnu rovnicu kruhovej
trajektórie družice, ktorá sa v okamžiku štartu nachádza vo výške
\(h\) nad zemským povrchom
v sústave, kde os \(y\) 
spojuje stred Zeme s miestom štartu družice a počiatok sústavy je na
povrchu Zeme.}
{\(x^{2} + (y + R_{Z})^{2} = (R_{Z} + h)^{2}\)}{\(x^{2} + y^{2} = (R_{Z} + h)^{2}\)}{\(x^{2} + (y + R_{Z})^{2} = h^{2}\)}{\(x^{2} + y^{2} = h^{2}\)}{}{}}
\LANGes{
\Question{
Los satélites se mueven alrededor de la Tierra con una trayectoria orbital. Consideremos un satélite a una altura \(h\) medida desde la Tierra. También consideremos el sistema de coordenadas cuyo origen es la Tierra directamente debajo del satélite y el eje \(y\) está orientado hacia arriba (al satélite). El eje \(x\)
es perpendicular al eje \(y\) y está en el plano definido por la trayectoria del satélite. Omitamos la rotación de la Tierra. Encuentra la ecuación de la trayectoria del satélite. El radio de la Tierra es \(R\).}
{\(x^{2} + (y + R)^{2} = (R + h)^{2}\)}{\(x^{2} + y^{2} = (R + h)^{2}\)}{\(x^{2} + (y + R)^{2} = h^{2}\)}{\(x^{2} + y^{2} = h^{2}\)}{}{}}
\End

\Data\ID{9000123101}
\Updated{2022-04-23 05:04:20}
\Level{C}
\SubArea{Conic Sections}
\LANGen{
\Question{
Find the values of the real parameter \(q\) 
which ensure that the line \(y = q\) 
is a tangent to the circle 
\[ 
 x^{2} + y^{2} + 4x - 8y + 4 = 0.
\]}
{\(\{0;8\}\)}{\(\{ - 6;2\}\)}{\(\{ - 8;0\}\)}{\(\{ - 2;6\}\)}{}{}}
\LANGpl{
\Question{
Wyznacz wartość rzeczywistą parametru  \(q\) 
tak, aby prosta \(y = q\) 
była styczną okręgu.
\[ 
 x^{2} + y^{2} + 4x - 8y + 4 = 0.
\]}
{\(\{0;8\}\)}{\(\{ - 6;2\}\)}{\(\{ - 8;0\}\)}{\(\{ - 2;6\}\)}{}{}}
\LANGcs{
\Question{
Množina všech hodnot reálného parametru
\(q\), pro které
je přímka \(y = q\) 
tečnou kružnice \(x^{2} + y^{2} + 4x - 8y + 4 = 0\),
je rovna:}
{\(\{0;8\}\)}{\(\{ - 6;2\}\)}{\(\{ - 8;0\}\)}{\(\{ - 2;6\}\)}{}{}}
\LANGsk{
\Question{
Množina všetkých hodnôt reálneho parametra
\(q\), pre ktoré
je priamka \(y = q\) 
dotyčnicou kružnice \(x^{2} + y^{2} + 4x - 8y + 4 = 0\),
je rovná:}
{\(\{0;8\}\)}{\(\{ - 6;2\}\)}{\(\{ - 8;0\}\)}{\(\{ - 2;6\}\)}{}{}}
\LANGes{
\Question{
Encuentra todos los valores del parámetro real \(q\) 
para que la recta \(y = q\) 
sea tangente a la circunferencia \[ 
 x^{2} + y^{2} + 4x - 8y + 4 = 0.
\]}
{\(\{0;8\}\)}{\(\{ - 6;2\}\)}{\(\{ - 8;0\}\)}{\(\{ - 2;6\}\)}{}{}}
\End

\Data\ID{9000123102}
\Updated{2022-04-23 05:04:20}
\Level{C}
\SubArea{Conic Sections}
\LANGen{
\Question{
Find a true statement about the ellipse 
\[ 
 x^{2} + 4y^{2} - 8y = 0.
\]}
{The tangent to the ellipse can pass through any point on the line
\(y = -1\).}{The tangent to the ellipse can pass through any point on the line
\(x = 1\).}{The tangent to the ellipse can pass through the point
\([-1;1]\).}{The tangent to the ellipse can pass through any point on the line
\(y = 1\).}{}{}}
\LANGpl{
\Question{
Oznacz zdanie prawdziwe określające elipsę.
\[ 
 x^{2} + 4y^{2} - 8y = 0.
\]}
{Styczna do elipsy może przechodzić przez każdy punkt na prostej
\(y = -1\).}{Styczna do elipsy może przechodzić przez każdy punkt na prostej
\(x = 1\).}{Styczna do elipsy może przechodzić przez punkt 
\([-1;1]\).}{Styczna do elipsy może przechodzić przez każdy punkt na prostej
\(y = 1\).}{}{}}
\LANGcs{
\Question{
Vytvořte pravdivé tvrzení: Tečnu k elipse \(x^{2} + 4y^{2} - 8y = 0\) 
lze vést}
{libovolným bodem přímky \(y = -1\).}{libovolným bodem přímky \(x = 1\).}{bodem \([-1;1]\).}{libovolným bodem přímky \(y = 1\).}{}{}}
\LANGsk{
\Question{
Vytvorte pravdivé tvrdenie: Dotyčnicu k elipse \(x^{2} + 4y^{2} - 8y = 0\) 
možno viesť}
{ľubovolným bodom priamky \(y = -1\).}{ľubovolným bodom priamky \(x = 1\).}{bodom \([-1;1]\).}{ľubovolným bodom priamky \(y = 1\).}{}{}}
\LANGes{
\Question{
Elige el enunciado correcto relacionado con la elipse  
\[ 
 x^{2} + 4y^{2} - 8y = 0.
\]}
{La tangente a la elipse puede pasar por cualquier punto de la recta \(y = -1\).}{La tangente a la elipse puede pasar por cualquier punto de la recta 
\(x = 1\).}{La tangente a la elipse puede pasar por el punto 
\([-1;1]\).}{La tangente a la elipse puede pasar por cualquier punto de la recta 
\(y = 1\).}{}{}}
\End

\Data\ID{9000123103}
\Updated{2021-10-28 09:10:24}
\Level{C}
\SubArea{Conic Sections}
\LANGen{
\Question{
The ellipse 
\[ 
 5x^{2} + 9y^{2} = 45
\]
has tangent \(2x + 3y = 9\). Find the
values of the real parameter \(k\) 
which ensure that the line \(y = kx + 3\) 
is a secant for the ellipse.}
{\(k\in \left (-\infty ;-\frac{2}
{3}\right )\cup \left (\frac{2}
{3};\infty \right )\)}{\(k\in \left [ -\frac{2}
{3}; \frac{2} 
{3}\right ] \)}{\(k\in \left (-\frac{2}
{3}; \frac{2} 
{3}\right )\)}{\(k\in \left (-\infty ;-\frac{2}
{3}\right ] \cup \left [ \frac{2}
{3};\infty \right )\)}{}{}}
\LANGpl{
\Question{
Dana elipsa
\[ 
 5x^{2} + 9y^{2} = 45
\]
ma styczną t \(2x + 3y = 9\). Wskaż wartość rzeczywistą parametru 
 \(k\) 
tak, aby prosta  \(y = kx + 3\) 
była sieczną elipsy.}
{\(k\in \left (-\infty ;-\frac{2}
{3}\right )\cup \left (\frac{2}
{3};\infty \right )\)}{\(k\in \left [ -\frac{2}
{3}; \frac{2} 
{3}\right ] \)}{\(k\in \left (-\frac{2}
{3}; \frac{2} 
{3}\right )\)}{\(k\in \left (-\infty ;-\frac{2}
{3}\right ] \cup \left [ \frac{2}
{3};\infty \right )\)}{}{}}
\LANGcs{
\Question{
Je dána elipsa \(5x^{2} + 9y^{2} = 45\) a její
tečna \(2x + 3y = 9\). Určete všechny
hodnoty parametru \(k\in \mathbb{R}\) 
tak, aby přímka \(y = kx + 3\) 
byla sečnou zadané elipsy.}
{\(k\in \left (-\infty ;-\frac{2}
{3}\right )\cup \left (\frac{2}
{3};\infty \right )\)}{\(k\in \left \langle -\frac{2}
{3}; \frac{2} 
{3}\right \rangle \)}{\(k\in \left (-\frac{2}
{3}; \frac{2} 
{3}\right )\)}{\(k\in \left (-\infty ;-\frac{2}
{3}\right \rangle \cup \left \langle \frac{2}
{3};\infty \right )\)}{}{}}
\LANGsk{
\Question{
Je daná elipsa \(5x^{2} + 9y^{2} = 45\) a jej
dotyčnica \(2x + 3y = 9\). Určte všetky
hodnoty parametra \(k\in \mathbb{R}\) 
tak, aby priamka \(y = kx + 3\) 
bola sečnica zadanej elipsy.}
{\(k\in \left (-\infty ;-\frac{2}
{3}\right )\cup \left (\frac{2}
{3};\infty \right )\)}{\(k\in \left \langle -\frac{2}
{3}; \frac{2} 
{3}\right \rangle \)}{\(k\in \left (-\frac{2}
{3}; \frac{2} 
{3}\right )\)}{\(k\in \left (-\infty ;-\frac{2}
{3}\right \rangle \cup \left \langle \frac{2}
{3};\infty \right )\)}{}{}}
\LANGes{
\Question{
La elipse
\[ 
 5x^{2} + 9y^{2} = 45
\]
tiene como recta  tangente \(2x + 3y = 9\). Encuentra todos los valores del parámetro real \(k\) de modo que la recta \(y = kx + 3\) 
sea secante a la elipse.}
{\(k\in \left (-\infty ;-\frac{2}
{3}\right )\cup \left (\frac{2}
{3};\infty \right )\)}{\(k\in \left [ -\frac{2}
{3}; \frac{2} 
{3}\right ] \)}{\(k\in \left (-\frac{2}
{3}; \frac{2} 
{3}\right )\)}{\(k\in \left (-\infty ;-\frac{2}
{3}\right ] \cup \left [ \frac{2}
{3};\infty \right )\)}{}{}}
\End

\Data\ID{9000123104}
\Updated{2021-10-14 07:10:02}
\Level{C}
\SubArea{Conic Sections}
\LANGen{
\Question{
In the following list identify a line which is a tangent to the ellipse 
\[ 
 (x - 2)^{2} + \frac{y^{2}} 
 {9} = 1
\]}
{\(\begin{aligned}[t] p\colon x& = 3 + t, &
\\y & = 3;\ t\in \mathbb{R}
\\ \end{aligned}\)}{\(p\colon x = 2\)}{\(p\colon y = 3x\)}{\(p\colon y = -x - 2\)}{}{}}
\LANGpl{
\Question{
Wskaż prostą, która jest styczną do elipsy.
\[ 
 (x - 2)^{2} + \frac{y^{2}} 
 {9} = 1
\]}
{\(\begin{aligned}[t] p\colon x& = 3 + t, &
\\y & = 3;\ t\in \mathbb{R}
\\ \end{aligned}\)}{\(p\colon x = 2\)}{\(p\colon y = 3x\)}{\(p\colon y = -x - 2\)}{}{}}
\LANGcs{
\Question{
Která z uvedených přímek je tečna elipsy
\((x - 2)^{2} + \frac{y^{2}} 
 {9} = 1\)?}
{\(\begin{aligned}p\colon x& = 3 + t,&
\\y & = 3;\ t\in \mathbb{R}
\\ \end{aligned}\)}{\(p\colon x = 2\)}{\(p\colon y = 3x\)}{\(p\colon y = -x - 2\)}{}{}}
\LANGsk{
\Question{
Ktorá z uvedených priamok je dotyčnica elipsy
\((x - 2)^{2} + \frac{y^{2}} 
 {9} = 1\)?}
{\(\begin{aligned}p\colon x& = 3 + t,&
\\y & = 3;\ t\in \mathbb{R}
\\ \end{aligned}\)}{\(p\colon x = 2\)}{\(p\colon y = 3x\)}{\(p\colon y = -x - 2\)}{}{}}
\LANGes{
\Question{
En la siguiente lista identifica una recta que sea  tangente a la elipse  
\[ 
 (x - 2)^{2} + \frac{y^{2}} 
 {9} = 1
\]}
{\(\begin{aligned}[t] p\colon x& = 3 + t, &
\\y & = 3;\ t\in \mathbb{R}
\\ \end{aligned}\)}{\(p\colon x = 2\)}{\(p\colon y = 3x\)}{\(p\colon y = -x - 2\)}{}{}}
\End

\Data\ID{9000123105}
\Updated{2021-10-14 07:10:23}
\Level{C}
\SubArea{Conic Sections}
\LANGen{
\Question{
Find the all values of the real parameter
\(p\) which ensure
that the line \(q\colon y = x - 1\) 
is a tangent to the parabola 
\[ 
 x^{2} = 2py.
\]}
{\(p = 2\)}{\(p\in \{0;2\}\)}{\(p = -2\)}{\(p\in \{ - 2;0\}\)}{}{}}
\LANGpl{
\Question{
Wskaż wszystkie wartości rzeczywistego parametru 
\(p\) tak, aby
prosta  \(q\colon y = x - 1\) 
była styczną do  paraboli.
\[ 
 x^{2} = 2py.
\]}
{\(p = 2\)}{\(p\in \{0;2\}\)}{\(p = -2\)}{\(p\in \{ - 2;0\}\)}{}{}}
\LANGcs{
\Question{
Určete všechny hodnoty parametru
\(p\in \mathbb{R}\) tak, aby přímka \(q\colon y = x - 1\) byla tečnou paraboly \(x^{2} = 2py\).}
{\(p = 2\)}{\(p\in \{0;2\}\)}{\(p = -2\)}{\(p\in \{ - 2;0\}\)}{}{}}
\LANGsk{
\Question{
Určte všetky hodnoty parametra
\(p\in \mathbb{R}\) tak, aby priamka \(q\colon y = x - 1\) bola dotyčnicou paraboly \(x^{2} = 2py\).}
{\(p = 2\)}{\(p\in \{0;2\}\)}{\(p = -2\)}{\(p\in \{ - 2;0\}\)}{}{}}
\LANGes{
\Question{
Halla todos los valores del parámetro \(p\in \mathbb{R}\) de manera de que la recta \(q\colon y = x - 1\) 
sea  tangente a la parábola \[ 
 x^{2} = 2py.
\]}
{\(p = 2\)}{\(p\in \{0;2\}\)}{\(p = -2\)}{\(p\in \{ - 2;0\}\)}{}{}}
\End

\Data\ID{9000123106}
\Updated{2022-04-23 05:04:20}
\Level{C}
\SubArea{Conic Sections}
\LANGen{
\Question{
Find the tangent line \(q\) to
the parabola \(4(y - 2) = (x + 1)^{2}\), so that the tangent \(q\)
is parallel to the line \(p\colon 4x - 5y + 17 = 0.\)}
{\(q\colon 20x - 25y + 54 = 0\)}{\(q\colon 20x - 25y - 27 = 0\)}{\(q\colon 4x - 5y + 27 = 0\)}{\(q\colon 4x -5y - 17 = 0\)}{}{}}
\LANGpl{
\Question{
Wskaż styczną \(q\) do
 paraboli \(4(y - 2) = (x + 1)^{2}\) tak, aby była
równoległa  do prostej  \(p\colon 4x - 5y + 17 = 0.\)}
{\(q\colon 20x - 25y + 54 = 0\)}{\(q\colon 20x - 25y - 27 = 0\)}{\(q\colon 4x - 5y + 27 = 0\)}{\(q\colon 4x -5y - 17 = 0\)}{}{}}
\LANGcs{
\Question{
Rovnice tečny paraboly \(4(y - 2) = (x + 1)^{2}\), která
je rovnoběžná s přímkou \(4x - 5y + 17 = 0\),
má tvar:}
{\(20x - 25y + 54 = 0\)}{\(20x - 25y - 27 = 0\)}{\(4x - 5y + 27 = 0\)}{\( 4x -5y - 17 = 0\)}{}{}}
\LANGsk{
\Question{
Rovnice dotyčnice paraboly \(4(y - 2) = (x + 1)^{2}\), ktorá
je rovnobežná s priamkou \(4x - 5y + 17 = 0\),
má tvar:}
{\(20x - 25y + 54 = 0\)}{\(20x - 25y - 27 = 0\)}{\(4x - 5y + 27 = 0\)}{\(4x -5y - 17 = 0\)}{}{}}
\LANGes{
\Question{
Encuentra la tangente \(q\) a la parábola \(4(y - 2) = (x + 1)^{2}\) que sea paralela a la recta \(p\colon 4x - 5y + 17 = 0.\)}
{\(q\colon 20x - 25y + 54 = 0\)}{\(q\colon 20x - 25y - 27 = 0\)}{\(q\colon 4x - 5y + 27 = 0\)}{\(q\colon 4x -5y - 17 = 0\)}{}{}}
\End

\Data\ID{9000123107}
\Updated{2021-10-28 09:10:18}
\Level{C}
\SubArea{Conic Sections}
\LANGen{
\Question{
In the following list identify a line such that the line has a unique intersection with
the hyperbola 
\[ 
 x^{2} - y^{2} = 5
\]
but the line is not the tangent to this hyperbola.}
{\(p\colon \frac{x}
{5} + \frac{y} 
{5} = 1\)}{\(p\colon y = 5x\)}{\(p\colon 2x + y = 5\)}{\(\begin{aligned}[t] p\colon x& = 1 &
\\y & = -1 + t\text{; }t\in \mathbb{R}
\\ \end{aligned}\)}{}{}}
\LANGpl{
\Question{
Wskaż prostą mającą dokładnie jeden punkt przecięcia z hiperbolą
\[ 
 x^{2} - y^{2} = 5
\]
tak, aby prosta nie była styczną do hiperboli.}
{\(p\colon \frac{x}
{5} + \frac{y} 
{5} = 1\)}{\(p\colon y = 5x\)}{\(p\colon 2x + y = 5\)}{\(\begin{aligned}[t] p\colon x& = 1 &
\\y & = -1 + t\text{; }t\in \mathbb{R}
\\ \end{aligned}\)}{}{}}
\LANGcs{
\Question{
Která z uvedených přímek má s hyperbolou
\(x^{2} - y^{2} = 5\) 
právě jeden společný bod a přitom není její tečna?}
{\(p\colon \frac{x}
{5} + \frac{y} 
{5} = 1\)}{\(p\colon y = 5x\)}{\(p\colon 2x + y = 5\)}{\(\begin{aligned}p\colon x& = 1; &
\\y & = -1 + t\text{, }t\in \mathbb{R}
\\ \end{aligned}\)}{}{}}
\LANGsk{
\Question{
Ktorá z uvedených priamok má s hyperbolou
\(x^{2} - y^{2} = 5\) 
práve jeden spoločný bod a pritom nie je jej dotyčnica?}
{\(p\colon \frac{x}
{5} + \frac{y} 
{5} = 1\)}{\(p\colon y = 5x\)}{\(p\colon 2x + y = 5\)}{\(\begin{aligned}p\colon x& = 1; &
\\y & = -1 + t\text{, }t\in \mathbb{R}
\\ \end{aligned}\)}{}{}}
\LANGes{
\Question{
En la siguiente lista identifica la recta que tiene una única intersección con la hipérbola \[ 
 x^{2} - y^{2} = 5
\]
pero al mismo tiempo la recta no es tangente a la hipérbola.}
{\(p\colon \frac{x}
{5} + \frac{y} 
{5} = 1\)}{\(p\colon y = 5x\)}{\(p\colon 2x + y = 5\)}{\(\begin{aligned}[t] p\colon x& = 1 &
\\y & = -1 + t\text{; }t\in \mathbb{R}
\\ \end{aligned}\)}{}{}}
\End

\Data\ID{9000123108}
\Updated{2022-04-23 05:04:20}
\Level{C}
\SubArea{Conic Sections}
\LANGen{
\Question{
Find all the tangents to the hyperbola
\(x^{2} - 2y^{2} = 8\) 
such that the angle between each tangent and the
\(x\)-axis
is \(45^{\circ }\).}
{\(y = x + 2\text{, }y = x - 2\text{, }y = -x + 2\text{, }y = -x - 2\)}{\(y = x + 2\text{, }y = x - 2\)}{\(y = x + 2\text{, }y = -x + 2\)}{\(y = x + 2\)}{}{}}
\LANGpl{
\Question{
Wskaż wszystkie styczne do hiperboli
\(x^{2} - 2y^{2} = 8\) 
tak, aby kąt pomiędzy nimi a osią 
\(x\)
wynosił  \(45^{\circ }\).}
{\(y = x + 2\text{, }y = x - 2\text{, }y = -x + 2\text{, }y = -x - 2\)}{\(y = x + 2\text{, }y = x - 2\)}{\(y = x + 2\text{, }y = -x + 2\)}{\(y = x + 2\)}{}{}}
\LANGcs{
\Question{
Všechny tečny hyperboly \(x^{2} - 2y^{2} = 8\),
jejichž odchylka s osou \(x\) 
je rovna \(45^{\circ}\),
mají rovnice:}
{\(y = x + 2\text{, }y = x - 2\text{, }y = -x + 2\text{, }y = -x - 2\)}{\(y = x + 2\text{, }y = x - 2\)}{\(y = x + 2\text{, }y = -x + 2\)}{\(y = x + 2\)}{}{}}
\LANGsk{
\Question{
Všetky dotyčnice hyperboly \(x^{2} - 2y^{2} = 8\),
ktorých odchýlka s osou \(x\) 
je rovná \(45^{\circ}\),
majú rovnice:}
{\(y = x + 2\text{, }y = x - 2\text{, }y = -x + 2\text{, }y = -x - 2\)}{\(y = x + 2\text{, }y = x - 2\)}{\(y = x + 2\text{, }y = -x + 2\)}{\(y = x + 2\)}{}{}}
\LANGes{
\Question{
Halla todas las tangentes a la hipérbola \(x^{2} - 2y^{2} = 8\) 
para las que el ángulo entre cada tangente y el eje \(x\)
es \(45^{\circ }\).}
{\(y = x + 2\text{, }y = x - 2\text{, }y = -x + 2\text{, }y = -x - 2\)}{\(y = x + 2\text{, }y = x - 2\)}{\(y = x + 2\text{, }y = -x + 2\)}{\(y = x + 2\)}{}{}}
\End
